

\documentclass[11pt]{amsart}
%\usepackage{amsfonts,amssymb,amscd,amsmath,enumerate,verbatim,calc}
\usepackage{amsmath}
\usepackage{amssymb}
\usepackage{amscd}
%\usepackage{showkeys}

%------    GENERAL MACROS    -----

\newcommand{\ext}{\operatorname{Ext}}
\newcommand{\reg}{\operatorname{reg}}
\def\height{\operatorname{ht}}
%\def\N{{\mathcal N}}
\newcommand{\sfrac}[2]{\frac{\displaystyle #1}{\displaystyle #2}}

\def\NZQ{\mathbb}               % the font for N,Z,Q,R,C
\def\NN{{\NZQ N}}
\def\QQ{{\NZQ Q}}
\def\ZZ{{\NZQ Z}}
\def\RR{{\NZQ R}}
\def\CC{{\NZQ C}}
\def\AA{{\NZQ A}}
\def\BB{{\NZQ B}}
\def\PP{{\NZQ P}}
\def\FF{{\NZQ F}}
\def\GG{{\NZQ G}}
\def\HH{{\NZQ H}}
\def\P{\mathcal P}
\def\mfp{\mathfrak p}
\def\mfa{\mathfrak a}
\def\Spec{\operatorname{Spec}}
\def \mm{{m_R}}
\def\Ass{\operatorname{Ass}}
\def\MinAss{\operatorname{MinAss}}
\def\height{\operatorname{height}}

\newcommand{\bfind}[1]{\index{#1}{\bf #1}}
\newcommand{\n}{\par\noindent}
\newcommand{\sn}{\par\smallskip\noindent}
\newcommand{\mn}{\par\medskip\noindent}
\newcommand{\bn}{\par\bigskip\noindent}
\newcommand{\pars}{\par\smallskip}
\newcommand{\parm}{\par\medskip}
\newcommand{\parb}{\par\bigskip}
\newcommand{\cal}{\mathcal}
\newcommand{\isom}{\simeq}
\newcommand{\ovl}[1]{\overline{#1}}
\newcommand{\dec}{^{\it dec}}
\newcommand{\ina}{^{\it in}}
\newcommand{\sep}{^{\it sep}}
\newcommand{\chara}{\mbox{\rm char}\,}
\newcommand{\trdeg}{\mbox{\rm trdeg}\,}
\newcommand{\Gal}{\mbox{\rm Gal}\,}
\newcommand{\Fix}{\mbox{\rm Fix}}
\newcommand{\rr}{\mbox{\rm rr}\,}
\newcommand{\crf}{\mbox{\rm crf}\,}
\newcommand{\tr}{\mbox{\rm Tr}\,}
\newcommand{\cO}{\mathcal{O}}
\newcommand{\cM}{\mathcal{M}}
\newcommand{\cE}{\mathcal{E}}
\newcommand{\cA}{\mathcal{A}}
\newcommand{\cC}{\mathcal{C}}
\newcommand{\rme}{\mbox{\rm e}\,}
\newcommand{\rmf}{\mbox{\rm f}\,}
\newcommand{\fs}{\!\!\uparrow\,}

%
\newcommand{\R}{\mathbb R}
\newcommand{\Q}{\mathbb Q}
\newcommand{\N}{\mathbb N}
\newcommand{\Z}{\mathbb Z}
\newcommand{\F}{\mathbb F}
\newcommand{\Qp}{{\mathbb Q}_p}



\newtheorem{theorem}{Theorem}[section]
\newtheorem{proposition}[theorem]{Proposition}
\newtheorem{lemma}[theorem]{Lemma}
\newtheorem{corollary}[theorem]{Corollary}
\theoremstyle{definition}
\newtheorem{definition}[theorem]{Definition}
\newtheorem{exercise}[theorem]{Exercise}
\newtheorem{example}[theorem]{Example}
%\theoremstyle{remark}
%\newtheorem{remark}[theorem]{Remark}
%\newtheorem{Remark}[theorem]{Remarks}
\renewcommand{\qedsymbol}{$\square$}


%------------------------------------------------
% Direct and inverse limits
%
%\opn\dirlim{\underrightarrow{\lim}}
%\opn\inivlim{\underleftarrow{\lim}}
%
%
% Names with a meaning
%
\let\union=\cup
\let\sect=\cap
\let\dirsum=\oplus
\let\tensor=\otimes
\let\iso=\cong
\let\Union=\bigcup
\let\Sect=\bigcap
\let\Dirsum=\bigoplus
\let\Tensor=\bigotimes

\newtheorem{Theorem}{Theorem}[section]
\newtheorem{Lemma}[Theorem]{Lemma}
\newtheorem{Corollary}[Theorem]{Corollary}
\newtheorem{Proposition}[Theorem]{Proposition}
\newtheorem{Remark}[Theorem]{Remark}
\newtheorem{Remarks}[Theorem]{Remarks}
\newtheorem{Example}[Theorem]{Example}
\newtheorem{Examples}[Theorem]{Examples}
\newtheorem{Definition}[Theorem]{Definition}
\newtheorem{Problem}[Theorem]{}
\newtheorem{Conjecture}[Theorem]{Conjecture}
\newtheorem{Question}[Theorem]{Question}
%
% We like the var forms of some greek letters (as taught in German schools)
%
\let\epsilon\varepsilon
\let\phi=\varphi
\let\kappa=\varkappa
%
%           We print on A4 paper
%
\textwidth=15cm \textheight=22cm \topmargin=-0.5cm
\oddsidemargin=0.5cm \evensidemargin=0.5cm \pagestyle{plain}
% ------    END OF GENERAL MACROS    -------
\begin{document}

\title{A general theorem on the analytic spread and associated primes of divisorial filtrations}
\author{Steven Dale Cutkosky}
\thanks{Partially supported by NSF grant DMS-2348849}

\address{Steven Dale Cutkosky, Department of Mathematics,
University of Missouri, Columbia, MO 65211, USA}
\email{cutkoskys@missouri.edu}



%\keywords{morphism, valuation, positive characteristic}
\subjclass[2000]{primary 13A15, 13A18;  secondary 13A02, 14C20} 


\begin{abstract} 
We prove that a classical theorem of McAdam about the  analytic spread of an ideal in a Noetherian local ring $R$ is true for rational divisorial filtrations on an
arbitrary Noetherian local ring. We deduce McAdam's original theorem from our result.
This generalizes earlier work by the author, showing that this theorem is true for rational divisorial filtrations, with the restriction that $R$
 is an equidimensional local ring which is essentially of finite type over an excellent local ring of dimension $\le 3$.
\end{abstract}

\maketitle 



\section{Introduction}

Notation used in this paper is explained in Section \ref{SecNot}.

McAdam proved the following important theorem in the theory of ideals. 
\begin{Theorem}\label{Theorem2}(\cite{McA}, Theorem 5.4.6 \cite{HS})  Let $R$ be a formally equidimensional local ring and $I$ be an ideal in $R$. Then $m_R\in \mbox{Ass}(R/\overline{I^n})$ if and only if the analytic spread $\ell(I)=\dim R$.
\end{Theorem}

In this paper, we generalize this theorem to divisorial filtrations in arbitrary Noetherian local rings. These filtrations are generally not Noetherian. Theorem \ref{Theorem1} is proven in Section \ref{SecThm1}.

 \begin{Theorem}\label{Theorem1} Let $R$ be a Noetherian local ring of dimension $d$, and $\mathcal I=\{I_n\}$ be a rational divisorial filtration of $R$. Then the following are equivalent. 
 \begin{enumerate}
 \item[1)] The analytic spread $\ell(\mathcal I)=d$.
 \item[2)] $m_R\in \mbox{Ass}_R(R/I_n)$ for some $n>0$.
 \item[3)] $m_R\in \mbox{Ass}_R(R/I_n)$ for all $n\gg 0$.
 \end{enumerate}
 \end{Theorem}
 
 Theorem \ref{Theorem2} is deduced from Theorem \ref{Theorem1} in Section \ref{SecThm2}.
 
 We reproduce a simple example from \cite{C1} showing that 
 the conclusion 2) implies 1) of Theorem \ref{Theorem1} do not hold for more general filtrations.

\begin{Example}(Example 1.5 \cite{C1}) There exist $\RR$-divisorial filtrations $\mathcal I$ such that 
the conclusions of Theorem \ref{Theorem1} are false. 
\end{Example}

\begin{proof} Let $R=k[[x]]$ be a power series ring in one variable over a field $k$. Let $I_n=(x^{\lceil n\pi\rceil})$ for $n\in \NN$, and $\mathcal I=\{I_n\}$. For fixed $n$, $r\lceil n\pi\rceil>\lceil rn\pi\rceil +1$ for some  $r\in \ZZ_{>0}$. Thus $f\in I_n$ implies $f^r\in m_RI_{nr}$. Thus $\ell(\mathcal I)=\dim R[\mathcal I]/m_RR[\mathcal I]=0<1=\dim R$.
\end{proof} 
 
 
 In \cite{C1}, Theorem \ref{Theorem1} is proven with the assumption that $R$ is an excellent local domain of equicharacteristic zero or of dimension $\le 3$ (or more generally when resolution of singularities holds). In \cite{C2}, the theorem is proven more generally with the assumption that $R$ is an equidimensional local ring which is essentially of finite type over an excellent local ring of dimension $\le 3$.

In the papers \cite{C1} and \cite{C2}, geometric methods are used to realize a quotient of 
\linebreak $R[\mathcal I]/m_RR[\mathcal I]$ as a subalgebra of a finitely generated $k$-algebra where  $k=R/m_R$ (a subfinite $k$-algebra), such that this quotient contains a finitely generated $k$-algebra of dimension $d$. By a theorem of Kemp and Trung in \cite{KT}, this implies that $\ell(\mathcal I)\ge d$. Since $\ell(\mathcal I)\le d$ for all graded families of ideals on $R$ by a lemma in \cite{CPS}, this proves 2) implies 1), with the assumptions of those papers. The proof of 2) implies 1) in \cite{C1} used a relative form of the theory of $\sigma$-decomposition of pseudo effective divisors on a projective nonsingular variety of Nakayama, developed in \cite{N}. This theory requires resolution of singularities, which is not known to be true in dimension larger than three. This problem is finessed in \cite{C2}, by using a theorem of Temkin on existence of regular alterations \cite{T}, which allowed the theorem to be extended to rings which are essentially of finite type over an excellent local ring of dimension three. 

The proof in this paper also proves 2) implies 1) by showing that a quotient of 
\linebreak $R[\mathcal I]/m_RR[\mathcal I]$ is 
a subfinite $k$-algebra  which contains a finitely generated $k$-algebra of dimension $d$. What is different in this paper is that we show that a divisorial valuation $v$ which is essential for the filtration $\mathcal I$ is centered at $m_R$ and is such that the associated graded ring ${\rm gr}_v(R)$ is subfinite.  We then show that a quotient $C$ of $R[\mathcal I]/m_RR[\mathcal I]$ is a subalgebra of ${\rm gr}_v(R)$, and then 
use the fact that $\mbox{trdeg}_{k}k(v)=d-1$ to prove that $C$ contains a polynomial ring over $k$ of dimension $d$. 

The proofs of 1) implies 3) in \cite{C1} and \cite{C2} follow from Theorem 4.6 and the following remark in \cite{CPS}, and the generalization in Theorem 1.4 \cite{C2}.

\begin{Theorem}(\cite[Theorem 4.6]{CPS} and Theorem 1.4 \cite{C2}))\label{TheoremN2} Suppose that $R$ is a Noetherian local ring and $\mathcal I=\{I_n\}$ is  an $\RR$-divisorial filtration on $R$ such that $\ell(\mathcal I)=\dim R$. Then there exists $n_0\in \NN$ such that $m_R\in \mbox{Ass}(R/\overline{I^n})$ for all $n\ge n_0$.
 \end{Theorem} 

1) implies 3) of Theorem \ref{Theorem1} also follows from Theorem \ref{TheoremN2}. We give a different proof of 1) implies 3) of Theorem \ref{Theorem1} in this paper, at the end of Section \ref{SecThm1}.

Theorem \ref{TheoremN2} above generalizes Burch's theorem for ideals (\cite{Bu} and Proposition 5.4.7 \cite{HS}) to real divisorial filtrations.
 

\section{Notation and preliminaries}\label{SecNot} Let $R$ be a commutative ring. If $P$ is a prime ideal of $R$ then we denote the residue field of $P$ by
$\kappa(P)=(R/P)_P= \mbox{QF}(R/P)$, where $\mbox{QF}(A)$ denotes the quotient field of a domain $A$. If $R$ is a local ring, we denote the maximal ideal of $R$ by $m_R$. If $v$ is a valuation of a field $K$ we denote the valuation ring of $v$ by $\mathcal O_v$, its maximal ideal by $m_v$ and its residue field by $\kappa(v)=\mathcal O_v/m_v$. We denote the value group of $v$ by $\Gamma_v$.

\begin{Definition}(Definition 9.3.1 \cite{HS}) Let $R$ be a Noetherian integral domain with field of fractions $K$ and let $v$ be a valuation of $K$ such that $v$ is nonnegative on $R$. Let $P=m_v\cap R$. If $\mbox{trdeg}_{\kappa(P)}\kappa(v)=\mbox{ht }P-1$, then $v$ is said to be a divisorial valuation with respect to $R$.
\end{Definition}
The valuation ring of a divisorial valuation is Noetherian and has value group $\Gamma_v\cong \ZZ$, as shown in \cite{HS}. 

\begin{Definition} Suppose that $R$ is a $d$-dimensional  local ring. A divisorial valuation of $R$ is a divisorial valuation $v$ of the quotient ring $R/P$ where $P$ is a minimal prime of $R$ such that $\dim R/P=d$. We extend $v$ to a pseudo valuation  $v:R\rightarrow \Gamma_v\cup\{\infty\}$ where $v(f)=\infty$ if $f\in P$. 
\end{Definition}

Suppose that  $R$ is a  local ring and $v$ is a divisorial valuation of $R$, which we may assume has value group $\Gamma_v=\ZZ$. We define valuation ideals in $R$ for $\lambda\in \RR_{\ge 0}$ by
$I(v)_\lambda=\{f\in R\mid v(f)\ge \lambda\}$. $I(v)_1$ is a prime ideal $P$, called the center of $v$ on $R$. It is the preimage of $m_v$ in $R$. We also will say that $v$ is centered at $P$ on $R$.
$I(v)_0=R$, $I(v)_{\lambda}=I(v)_1$ if $0<\lambda\le 1$ and $I(v)_{\lambda}$ is $I(v)_1$-primary if $\lambda>0$.
We define the associated graded ring of the valuation by 
$$
{\rm gr}_v(R)=\oplus_{n\ge 0}I(v)_n/I(v)_{n+1}.
$$
The ring ${\rm gr}_v(R)$ is an integral domain, since $v(ab)=v(a)+v(b)$ for all $a,b\in R$.
For $f\in R$, the initial form of $f$ with respect to $v$ is the class ${\rm in}_v(f)$ of $f$ in 
$$
I(v)_{v(f)}/I(v)_{v(f)+1}\subset {\rm gr}_v(R).
$$

A graded family of ideals $\mathcal I=\{I_n\}_{n\in \ZZ_{\ge 0}}$ in a ring $R$ is a family of ideals $I_n$ in $R$ such that $I_mI_n\subset I_{m+n}$ if $m,n\in \ZZ_{\ge 0}$. The Rees algebra of a graded family of ideals $\mathcal I$ is $R[\mathcal I]=\oplus_{n\ge 0}I_n$ and the analytic spread of $\mathcal I$ is 
$\ell(\mathcal I)=\dim R[\mathcal I]/m_RR[\mathcal I]$. In the case that $\mathcal I$ is an $I$-adic filtration $\mathcal I=\{I^n\}$ where $I$ is an ideal of $R$, these definitions give the Rees algebra $R[I]$ and the analytic spread $\ell(I)$ of the ideal $I$. In a $d$-dimensional Noetherian local ring $R$, and graded family of ideals $\mathcal I$, we have that 
$\ell(\mathcal I)\le d$ and all values from 0 to $d$ can be realized, as shown in \cite{CPS}.



 Suppose that $R$ is a Noetherian local ring. A graded family of ideals $\mathcal I$ in $R$ is called a divisorial filtration if there exist $r\in\ZZ_{>0}$, $a_1,\ldots, a_r$ in $\RR_{>0}$ and divisorial valuations $v_1\ldots,v_r$ of $R$ such that
\begin{equation}\label{eq1}
I_n=\{f\in R\mid v_i(f)\ge n a_i\mbox{ for }1\le i\le r\}.
\end{equation}
The expression (\ref{eq1}) is called a presentation of the divisorial filtration $\mathcal I$. In general, there are infinity many different representations of $\mathcal I$. If a representation has the property that for each $j$ with $1\le j\le r$ the ideal $\{f\in R\mid v_i(f)\ge n a_i\mbox{ for }i\ne j\}$ properly contains $I_n$, then the representation is said to be irredundant. 

$\mathcal I$ is said to be a rational divisorial filtration if there exists a   presentation (\ref{eq1})  such that $a_i\in \QQ_{>0}$ for all $i$. $\mathcal I$ is an integral divisorial filtration if all $a_i\in \ZZ$.


\section{Proof of Theorem \ref{Theorem1}}\label{SecThm1} 

\begin{Lemma}\label{Lemma1} Let $S$ be a Noetherian local domain of dimension $e$ and let $v$ be a divisorial valuation of $S$ whose center is $m_S$.
Identify the value group of $v$ with $\ZZ$. Then there exists a complete local domain $T$ which is a quotient of the $m_S$-adic completion $\hat S$ and a divisorial valuation $w$ of $T$ such that $w| S=v$ and ${\rm gr}_v(S)\cong {\rm gr}_w(T)$ as graded $k=\kappa(m_R)$-algebras.
\end{Lemma}

\begin{proof} Let $V=\mathcal O_v$, which is a rank 1 discrete valuation ring. Let $\hat V$ be the completion of $V$ at its maximal ideal, with natural valuation $\hat v$ extending $v$. Since the center of $v$ on $S$ is $m_S$, we have that $v(m_S)\ge 1$, so that the natural homomorphism $S\rightarrow V$ is continuous for the $m_S$-adic and $m_V$-adic topologies, and thus extends to a homomorphism $\hat\Phi:\hat S\rightarrow \hat V$. Let $H$ be the kernel of $\Phi$ and $T=\hat S/H$. Then $H$ is a prime ideal, $H\cap S=0$ and $T$ is a complete Noetherian local domain embedded in $\hat V$. Let $w$ be the restriction of $\hat v$ to $\mbox{QF}(T)$. $w$ is nonnegative on $T$ and is centered at the maximal ideal of $T$. We have that 
$\kappa(m_T)=\kappa(m_R)$, $\Gamma_w=\Gamma_v=\ZZ$ and $\kappa(w)=\kappa(v)$. By Abhyankar's inequality, \cite{Ab} or Theorem 6.6.7 \cite{HS},
$$
e=1+{\rm trdeg}_k\kappa(v)\le \dim T\le \dim \hat S=e.
$$
Thus $\dim T=e$ and so $w$ is divisorial on $T$. The inclusions $S\rightarrow T$ give injective maps
\begin{equation}\label{eq2}
I(v)_j/I(v)_{j+1}\rightarrow I(w)_j/I(w)_{j+1}
\end{equation}
For $y\in I(w)_j$, there exists $x\in S$ such that $x=y+h$ with $h\in m_T^n$, with $n$ sufficiently large that $nw(m_T)>j$. Then $v(x)=j$ and $y-x\in I(w)_{j+1}$, so that (\ref{eq2}) is an isomorphism, giving the isomorphism ${\rm gr}_v(S)\cong {\rm gr}_w(T)$.
\end{proof}

Let $k$ be a field. A $k$-algebra is subfinite if it is a subalgebra of a finitely generated $k$-algebra.

\begin{Lemma}\label{Lemma2}
Let $R$ be an excellent local domain of dimension $d$ and $v$ be a divisorial valuation of the quotient field $K={\rm QF}(R)$ of $R$ whose center on $R$ is $m_R$. Let $k=\kappa(m_R)$. Then ${\rm gr}_v(R)$ is a subfinite $k$-algebra.
\end{Lemma}

\begin{proof} Let $\pi:X\rightarrow \mbox{Spec}(R)$ be such that $X$ is normal and there is a prime divisor $E\subset X$ such that $\mathcal O_{X,E}=\mathcal O_v$. Such an $X$ exists by Lemma 3.3 \cite{C2}. Since the center of $v$ on $R$ is $m_R$, $\pi(E)=m_R$ so the closed integral subscheme 
$E$ of $X$ is projective over $k$. 
%Thus $\dim E=\mbox{trdeg}_k\kappa(v)=d-1$. 
Choose $g\in K^{\times}$ with $v(g)=1$ and set $B={\rm div}_X(g)-E$, where ${\rm div}_X(g)$ is the Weil divisor on $X$ associated to $g$, so the coefficient of $E$ in  $B$ is zero. Write $B=B_+-B_-$ where $B_+$ is the divisor of zeros of $B$ and $B_-$ is the divisor of poles. Since $X$ is projective over $\mbox{Spec}(R)$, there exists an effective ample Cartier divisor $H$ on $X$ such that $H\ge B_+$ and $E$ is not contained in the support of $H$. 

Suppose that  $f\in I(v)_j$. Then  $f$ is regular everywhere on $X$ and the order of vanishing of $f$ along $E$ is at least $j$. Thus ${\rm div}_X(f)\ge jE$, and so
$$
{\rm div}_X(\frac{f}{g^j})+jH={\rm div}_X(f)-j{\rm div}_X(g)+jH\ge j(H-B)\ge 0.
$$
Since $X$ is normal, this shows that
$$
\frac{f}{g^j}\in H^0(X,\mathcal O_X(jH)).
$$
Mapping $f$ to the image of $\frac{f}{g^j}$ after restriction to $E$, defines a map 
\begin{equation}\label{eq3}
I(v)_j\rightarrow H^0(E,\mathcal O_X(jH)\otimes_{\mathcal O_X}\mathcal O_E).
\end{equation}
$f\in I(v)_j$ is in the kernel of the map if and only if $v(\frac{f}{g^j})>0$ which holds if and only if $f\in I(v)_{j+1}$. Thus the kernel of (\ref{eq3}) is $I(v)_{j+1}$. These maps preserve multiplication. Thus they give an injective $k$-algebra homomorphism
$$
{\rm gr}_v(R)\rightarrow S=\oplus_{j\ge 0}H^0(E,\mathcal O_X(jH)\otimes_{\mathcal O_X}\mathcal O_E).
$$
The invertible sheaf $\mathcal O_X(jH)\otimes_{\mathcal O_X}\mathcal O_E$ is ample on the projective $k$-variety $E$. Hence $S$ is a finitely generated $k$-algebra proving Lemma \ref{Lemma2}.
\end{proof}

We now prove Theorem \ref{Theorem1}. We first prove that 2) implies 1). We can choose a presentation (\ref{eq1}) of $\mathcal I$ which is irredundant. Let $\mathfrak p_i$ be the center of $v_i$ on $R$, for $1\le i\le r$. Each ideal $I(v_i)_n$ is $\mathfrak p_i$-primary. The natural injection $R/I_n\rightarrow \oplus_iR/I(v_i)_n$ shows that $\mbox{Ass}_R(R/I_n)\subset \{\mathfrak p_1,\ldots,\mathfrak p_r\}$. By our assumption that 2) holds, we conclude that at least one valuation in the irredundant presentation is centered at $m_R$. 
Reindex the $v_i$ so that this valuation is $v_1$ and let $v=v_1$. Let 
$P=\{x\in R\mid v(x)=\infty\}$ and let $S=R/P$, so that $S$ is a $d$-dimensional local domain and $v$ is a divisorial valuation of $S$ centered at $m_S$. Let $a$ be the coefficient of $v=v_1$ in the defining inequalities of (\ref{eq1}). Let
$$
C=\sum_{n\ge 0}I_nt^n/(\sum_{n\ge 0}\{h\in I_n\mid v(h)>na\}t^n)=\sum_{n\ge 0}C_nt^n
$$
where $t$ is a variable of degree 1. We have a natural surjection of graded $k$-algebras
\begin{equation}\label{eq4}
R[\mathcal I]/m_RR[\mathcal I]\rightarrow C.
\end{equation}

If $na\in \ZZ$ then $I_n\cap I(v)_{na+1}=\{x\in I_n\mid v(x)>na\}$.
If $na\not\in \ZZ$ then every $h\in I_n$ has $v(h)\ge \lceil na \rceil>na$ so $C_n=0$. Thus we have a natural inclusion of $k$-algebras $C\rightarrow {\rm gr}_{v}S$. Thus $C$ is subfinite by Lemmas \ref{Lemma1} and \ref{Lemma2}. 

Since $v$ is divisorial and centered at $m_R$, we have that $\mbox{trdeg}_k\kappa(v)=d-1$. Choose algebraically independent elements 
$\xi_1,\ldots,\xi_{d-1}\in \kappa(v)$ over $k$. Let $\xi_0=1$. Lift $\xi_1,\ldots,\xi_d$ to elements of  $\mathcal O_v$, which necessarily have value zero, and lift $\xi_0$ to the 1 of $R$. After multiplying them by a suitable common element $h$ of $R$, we obtain $u_0,u_1,\ldots, u_{d-1}\in R$ which have the value
\begin{equation}\label{eq10}
v(u_0)=v(u_1)=\cdots=v(u_{d-1})=c>0,
\end{equation}
and the class of $\frac{u_i}{u_0}$ in $\kappa(v)$ is $\xi_i$ for $0\le i\le d-1$. Irredundancy of the $v_i$ implies that there exists an integer $N>0$ and $f\in R$ such that 
\begin{equation}\label{eq6}
v_i(f)\ge na_i\mbox{ for $i\ne 1$ and $v(f)<Na$}.
\end{equation} 
Set
\begin{equation}\label{eq7}
\delta=Na-v(f)\in \QQ_{>0}.
\end{equation}
Since $\delta$ is positive and rational, there exist positive integers $p$ and $q$ such that
\begin{equation}\label{eq12}
p\delta=qc.
\end{equation}
Let $h_i=f^pu_i^q$ for $0\le i\le d-1$. For every $v_i$ other than $v=v_1$, (\ref{eq6}) and nonnegativity on $R$ give
$$
v_j(h_i)=pv_j(f)+qv_j(u_i)\ge pNa_j.
$$
For $v$, equations  (\ref{eq10}), (\ref{eq7}) and (\ref{eq12}) give
$$
v(h_i)=pv(f)+qc=p(Na-\delta)+p\delta=pNa.
$$
Thus $h_i\in I_{pN}$ and each $h_it^{pN}$ has nonzero image $z_i$ in $C$. All the $z_i$ have the same positive degree $pN$.

We will now show that $z_0,\ldots,z_{d-1}$ are algebraically independent over $k$. Suppose otherwise, so that there exists a nonzero polynomial
$F(T_0,T_1,\ldots,T_{d-1})$ in the polynomial ring $k[T_0,\ldots,T_d]$ such that $F(z_0,\ldots,z_{d-1})=0$. Since $z_0,\ldots,z_{d-1}$ have the same degree, the homogeneous components of $F$ give homogeneous relations, so we may assume that $F$ is homogeneous, say of degree $b$. Now
$G={\rm gr}_v(S)$ is a domain. Since the natural map $C\rightarrow G$ is injective, $C$ is also a domain. Thus we have a natural inclusion
${\rm QF}(C)\rightarrow {\rm QF}(G)$ which takes $\frac{z_i}{z_0}$ to $\frac{{\rm in}_v(h_i)}{{\rm in}_v(h_0)}$.
Let 
$$
K_0(G)=\{\frac{a}{b}\in {\rm QF}(G)\mid a,b\mbox{ are nonzero homogeneous elements of $G$ of equal degree}\}\cup\{0\}.
$$
There is a canonical $k$-algebra isomorphism $K_0(G)\cong k(v)$ defined by 
$$
\frac{{\rm in}_v(f)}{{\rm in}_v(g)}\mapsto \overline{\left(\frac{f}{g}\right)},\mbox{ the class of 
$\frac{f}{g}$ in $\kappa(v)$},
$$
so we can identify $\frac{z_i}{z_0}$ with $\xi_i^q\in \kappa(v)$.
Since $F$ is homogeneous, $F(1,\frac{z_1}{z_0},\ldots,\frac{z_{d-1}}{z_0})=0$, so that $F(1,\xi_1^q,\ldots,\xi_{d-1}^q)=0$ in $\kappa(v)$. 
The algebraic independence of  $\xi_1^q,\ldots,\xi_{d-1}^q$ implies that $F(1,T_1,\ldots,T_{d-1})=0$. Since $F$ is homogeneous, we have that
$F(T_0,T_1,\ldots,T_{d-1})=0$ and so $F(T_0,T_1,\ldots,T_{d-1})=0$.
This contradiction shows that $z_0,z_1,\ldots,z_{d-1}$ are algebraically independent over $k$. Thus 
$$
k[T_0,\ldots,T_{d-1}]\cong k[z_0,\ldots,z_{d-1}]\subset C.
$$
Since $k[z_0,\ldots,z_{d-1}]\subset C$ are subfinite $k$-algebras, we have by Corollary 4.7 of \cite{KT} that 
$$
\dim C\ge \dim k[z_0,\ldots,z_{d-1}]=d.
$$ 
By (\ref{eq4}), we have that
$$
\ell_R(\mathcal I)=\dim R[\mathcal I]/m_RR[\mathcal I]\ge \dim C\ge d.
$$
Finally, $\ell_R(\mathcal I)\le \dim R=d$ by Lemma 3.6 \cite{CPS}, and so we have established that 2) implies 1).

We now prove that 1) implies 3). For $1\le i\le r$ and $n>0$, let
$$
\mathfrak p_i=\{h\in R\mid v_i(h)>0\}\mbox{ and }Q_{i,n}=\{h\in R\mid v_i(h)\ge na_i\}. 
$$
$\mathfrak p_i$ is the center of $v_i$ on $R$ and $Q_{i,n}$ is $\mathfrak p_i$-primary for $n>0$. 
Suppose that none of the $v_i$ have center $m_R$. By prime avoidance, there exists $x\in m_R\setminus\left((\cup_i\mathfrak p_i\right)\cup
\left(\cup_{P\in \mbox{Min}(R)}P\right)$, where the second union is over all minimal primes of $R$. Now $r\in (I_n:x)$ if and only if $v_i(xr)=v_i(r)\ge na_i$ for all $i$, so that 
$$(I_n:x)=I_n\mbox{ for all $n\ge 0$}.
$$
Thus 
\begin{equation}\label{eqL3}
I_n\cap xR=x(I_n:x)=xI_n.
\end{equation}
Let $R^*=R/xR$, $I_n^*=I_n(R/xR)$ and $\mathcal I^*=\{I_n^*\}$. By (\ref{eqL3}), there is a natural isomorphism
$$
R[\mathcal I]/xR[\mathcal I]\cong R^*[\mathcal I^*].
$$
Thus 
$$
R^*[\mathcal I^*]/m_RR[\mathcal I^*]\cong R[\mathcal I]/m_RR[\mathcal I]
$$
since $x\in m_R$. We have that $\dim R/xR\le d-1$ since $x$ avoids every minimal prime of $R$. Thus by Lemma 3.6 \cite{CPS}, we have that
$\ell(\mathcal I)=\ell(\mathcal I^*)\le d-1$, giving a contradiction. Thus some valuation, which we can assume is $v_1$, has center $\mathfrak p_1=m_R$.

Since our presentation of $\mathcal I$ is irredundant, there exists $N>0$ and $f\in R$ such that 
$$
v_i(f)\ge Na_i\mbox{ for $i>1$ and $c:=v_1(f)<Na_1$.}
$$
For each $n>0$, let $q_n=\lceil \frac{n}{N}\rceil$ and $h_n=f^{q_n}$. For every $i>1$,
$$
v_i(h_n)=q_nv_i(f)\ge q_nNa_i\ge na_i.
$$
Set $\epsilon=a_1-\frac{c}{N}>0$.  For $n\in\ZZ_{>0}$, 
$$
\left(\frac{n}{N}+1\right)c=n\frac{c}{N}+c=na_1-n\epsilon+c<na_1
$$
whenever $n>\frac{c}{\epsilon}$. Thus
$$
v_1(h_n)=q_nc< \left(\frac{n}{N}+1\right)c<na_1
$$
whenever $n>\frac{c}{\epsilon}$. Thus for all sufficiently large $n$,
$h_n\in\cap_{i\ne j}Q_{i,n}$, with $h_n\not\in Q_{1,n}$. Thus $h_n\not\in I_n$ and since $Q_{1,n}$ is $m_R$-primary, there exists $b_n>0$ such that $m_R^{b_n}\subset Q_{1,n}$ and so $m_R^{b_n}h_n\subset I_n$. Let $L=h_n(R/I_n)$. Let $s$ be the smallest power of $m_R$ such that $m_R^sL=0$. Let $0\ne z\in m_R^{s-1}L$. Then $m_Rz=0$ so $\mbox{Ann}(z)=m_R$. Thus $m_R\in \mbox{Ass}(R/I_n)$ for all $n$ sufficiently large, completing the proof of 1) implies 3).

The final implication 3) implies 2) is immediate, completing the proof of the theorem.

\section{Derivation of McAdam's Theorem \ref{Theorem2} from Theorem \ref{Theorem1}}\label{SecThm2}

Let $R$ be a formally equidimensional Noetherian local ring of dimension $d$ and $I$ be an ideal in $R$. Let $\mathcal I$ be the  filtration  $\mathcal I=\{\overline{I^n}\}$ of integral closures of powers of $I$.
Let 
$$
A=R[I]=\oplus_{n\ge 0}I^nt^n\mbox{ and }B=R[\mathcal I]=\oplus_{n\ge 0}\overline{I^n}t^n
$$
where $A$ and $B$ are graded by $\deg t=1$, be the Rees algebras of $I$ and $\mathcal I$ respectively. The natural inclusion $A\subset B$ is an integral extension. Thus the injection
$$
A/m_RB\cap A\rightarrow B/m_RB
$$
is also an integral extension. Let $\mathfrak p$ be a prime ideal of $A$ containing $m_RA$. Since $B$ is integral over $A$, there exists a prime ideal $\mathfrak q$ of $B$ such that $\mathfrak q\cap A=\mathfrak p$ by Theorem A.6 \cite{BH}. Thus $m_RB\cap A\subset \mathfrak q\cap A=\mathfrak p$, and so
$$
m_RB\cap A\subset \cap_{\mathfrak p \supset m_RA}\mathfrak p=\sqrt{m_RA}.
$$
Thus $\sqrt{m_RB\cap A}\subset \sqrt{m_RA}$. Certainly $\sqrt{m_RA}\subset \sqrt{m_RB\cap A}$, so that 
$$
\sqrt{m_RB\cap A}=\sqrt{m_RA}.
$$
Thus 
\begin{equation}\label{eqT1}
\ell(\mathcal I)=\dim B/m_RB=\dim A/m_RB\cap A=\dim A/m_RA=\ell(I),
\end{equation}
where the second equality is by Corollary A.8 \cite{BH}.

We first prove the theorem with the assumption that $I\not\subset P$ for every minimal prime ideal $P$ of $R$. 
The theory of Rees valuations of an ideal is developed in \cite{R} and Section 10 of \cite{HS}.
We have that $R/P$ is formally equidimensional of dimension $d$ (by Corollary B.4.3 \cite{HS}) and every localization of a formally equidimensional local ring ring is formally equidimensional (by Theorem B.5.2 \cite{HS}).  By Proposition 10.4.3 \cite{HS}, there exist $r>0$ and divisorial valuations $v_i$ on $R/P_i$ where $P_i$ is a minimal prime of $R$ and $a_i\in \ZZ_{>0}$ such that 
$$
\overline{I^n}=\{f\in R\mid v_i(f)\ge na_i\mbox{ for }1\le i\le r\}.
$$
It is implicitely assumed that $IV\ne 0$ for  $V=\mathcal O_v$ in this Proposition. 
This condition is satisfied since $I$ is not contained in a minimal prime ideal of $R$. The valuations $v_1,\ldots,v_r$ are called the Rees valuations of $I$.
The filtration $\mathcal I=\{\overline{I^n}\}$ is thus an integral divisorial filtration on $R$, so by Theorem \ref{Theorem1}, $m_R\in \mbox{Ass}(R/\overline{I^n})$ for some $n$ if and only if $\ell(\mathcal I)=d$. Since $\ell(\mathcal I)=\ell(I)$ by (\ref{eqT1}), Theorem \ref{Theorem2} is proven when $I$ is not contained in a minimal prime ideal of $R$.

Now assume that $I$ is any proper ideal of $R$ (possibly contained in a minimal prime of $R$). Let $R[X]$ be a polynomial ring over $R$ and let $S=R[X]_{\mathfrak n}$ where $\mathfrak n=(m_R,X)R[X]$. Then $S$ is a local ring with maximal ideal $m_S=\mathfrak n S$. Let $J=IS+XS$. The $m_S$-adic completion of $S$ is $\hat S\cong \hat R[[X]]$. Since $\hat R$ is equidimensional of dimension $d$, $\hat S$ is equidimensional of dimension $d+1$. $J$ contains the nonzero divisor $X$, so $J$ is not contained in any minimal prime of $S$. Thus 
\begin{equation}\label{eqT5}
m_S\in \mbox{Ass}_S(S/\overline{J^n})\mbox{ for some $n$ if and only if }\ell(J)=d+1
\end{equation}
by the first part of this proof. 

Let $\overline X$ be the class of $X$ in $F_J(S)=S[J]/m_SS[J]$.
$F_J(S)=\oplus_{n\ge 0}J^n/m_SJ^n$ where
$$
J^n/m_SJ^n\cong \oplus_{j=0}^n\left(I^{n-j}/m_RI^{n-j}\right)\overline X^j.
$$
Thus there is a natural ring isomorphism of the polynomial ring $F_I(R)[Z]$ over $F_I(R)=R[I]/m_RR[I]$ with $F_J(S)$, obtained by mapping $Z$ to $\overline X$. Hence
\begin{equation}\label{eqT6}
\ell(J)=\ell(I)+1.
\end{equation}
$$
\overline{J^n}=\overline{I^n}S+\overline{I^{n-1}}XS+\cdots+\overline IX^{n-1}S+X^nS.
$$
Thus as an $R$-module,
\begin{equation}\label{eqT8}
S/\overline{J^n}\cong \oplus_{j=0}^{n-1}(R/\overline{I^{n-j}})X^j.
\end{equation}
Let $M=S/\overline{J^n}$. Equation (\ref{eqT8}) implies that
\begin{equation}\label{eqT10}
L:=0:_M m_R\ne 0 \mbox{ if and only if }m_R\in \mbox{Ass}_R(R/\overline{I^s})\mbox{ for some }1\le s\le n.
\end{equation}
$X^jL\subset L$ for all $j\ge 0$ and $X^n\in J^n\subset \overline{J^n}$ since $X\in J$. Thus $X^nM=0$, and so $X^nL=0$.

Suppose that $L=0:_Mm_R\ne 0$ and let $z$ be a nonzero element of $L$.  Let $s\ge 1$ be the smallest positive integer such that $X^sz=0$. Then $w=X^{s-1}z\ne 0$. Thus $w\in 0:_Mm_R$ and $Xw=0$ so that $0\ne w\in 0:_Mm_S$ and thus $m_S\in\mbox{Ass}_S(M)=\mbox{Ass}_S(S/\overline{J^n})$.

Certainly $m_S\in \mbox{Ass}_S(M)$ implies $0:_Mm_R\ne 0$, so that 
$$
m_S\in \mbox{Ass}_S(M)\mbox{ if and only if }0:_Mm_R\ne 0.
$$
The conclusions of Theorem \ref{Theorem2} for arbitrary $I$ now follow from (\ref{eqT10}), (\ref{eqT5}) and (\ref{eqT6}).

\begin{thebibliography}{1000000000}


\bibitem{Ab} S. Abhyankar, On the valuations centered in a local domain. Amer.
J. Math. 78 (1956), 321–348.

\bibitem{Bu} L. Burch, Codimension and analytic spread, Proc. Cambridge Phil. Soc. 72 (1972), 369-373.

\bibitem{BH} W. Bruns and J. Herzog, Cohen Macaulay Rings, revised edition, Cambridge University Press 1993.


\bibitem{C1} S. D. Cutkosky, Analytic spread and associated primes of divisorial filtrations,
Proceedings of the AMS 154 (2026), 2721 - 2732.

\bibitem{C2} S.D. Cutkosky, The Rees algebra and analytic spread of a divisorial filtration,
Advances Math. 479 (2025).

\bibitem{CPS} S.D. Cutkosky and P. Sarkar, Analytic spread of filtrations and symbolic algebras, Journal of the London Math. Soc. 106 (2022), 2635 - 2662.

\bibitem{KT} G. Kemper and N.V. Trung, Krull dimension and monomial orders, J. Algebra 399 (2014), 782 - 800.

\bibitem{McA}  S. McAdam, Asymptotic prime divisors and analytic spread, Proc. AMS 90 (1980), 555-559.

\bibitem{N} N. Nakayama, Zariski-decomposition and abundance, MSJ Memoirs 14, Math. Soc. Japan,
Tokyo, 2004.

\bibitem{R} D. Rees, Valuations associated to ideals II, J. London Math. Soc. 31 (1956), 221 - 228.

\bibitem{HS} I. Swanson and C. Huneke, Integral Closure of Ideals, Rings and Modules, Cambridge University Press, 2006.

\bibitem{T} M. Temkin, Tame distillation and desingularization by p-alterations, Ann. of Math.
186 (2017), 97 - 126.

\end{thebibliography}

\end{document}

%\[
%\boxed{
%\overline{I^n}
%=
%K\cap
%\bigcap_{P\in\mathcal A}\;
%\bigcap_{j=1}^{r_P}
%\left\{
%f\in R:
%\nu_{P,j}(f\bmod P)\ge n a_{P,j}
%\right\}.
%}
%\tag{1}
%\] 
%The factor \(K\) records exactly the components on which \(I\) is zero. Indeed, if \(I\subseteq P\) and \(v=\nu\circ(R\to R/P)\), then
%\[
%v(I)=\infty,\]


Let $R$ be a Noetherian local ring and $\mathcal I$ be a graded $m_R$-family. We define the multiplicity $e(\mathcal I)$ of 
$\mathcal I$ as
$$
e(\mathcal I)=\limsup_{p\rightarrow\infty}\frac{e(I_p)}{p^d}.
$$

Although the volume $\mbox{vol}(\mathcal I)$ does not  exist in general as a limit on an arbitrary Noetherian local ring $R$,  the multiplicity $e(\mathcal I)$ always does exist. 

\begin{Theorem}(Theorem 1.3 [C1])\label{Theorem1}
  Let $R$ be a $d$-dimensional Noetherian local ring and $I = \{I_n\}$ be a graded family of $m_R$-primary ideals. Then the multiplicity
$$
e(\mathcal I)=\lim_{n\rightarrow\infty}\frac{e(I_n)}{n^d}
$$
exists as a limit.
\end{Theorem}




\section{Multiplicities of  graded families of ideals on families of schemes}\label{QSec1}
\thispagestyle{empty}

The questions of this section are an application of the general multiplicity theory defined and developed in  [C1], which is valid for 
graded families of $m_R$-primary ideals in a Noetherian local ring. This multiplicity was defined in Theorem \ref{Theorem1}.

We propose to extend some classical theorems of Teissier on an ideal on  a family of analytic spaces to graded families of ideals on families of schemes. It is expected that these questions will have  positive answers, by combining techniques of [C1] and [Li1]. 



In this section, we  have the following assumptions, generalizing  those of Section 3 of [Li1] from a single ideal $I$ to a graded family $\mathcal I$. 

Let $\phi:S\rightarrow R$ be a homomorphism of commutative Noetherian local rings.   For every prime ideal $P$ in $S$, let $k(P)$ be the fraction field
 of $S/P$ and  set $R(P)=R\otimes_Sk(P)$. If $\mathfrak m$ is a prime ideal in $R$, let $k(\mathfrak m)$ be the fraction field of $R/\mathfrak m$.
 
Let $J$ be an ideal of $R$ such that $R/J$ is finite over $S$ and $\phi^{-1}(J)=0$. Define 
$$
d_P(J)=\max_{\mathfrak m\in V(J(P))}\dim R(P)_{\mathfrak m}
$$
and
$$
e_P(J)=\sum_{\mathfrak m\in V(J(P)),\,\, \dim R(P)_{\mathfrak m}=d_P(J)}[k(\mathfrak m):k(P)]e(J(P)_{\mathfrak m})
$$
where $e(J(P)_{\mathfrak m})$ is the multiplicity of the ideal $J(P)_{\mathfrak m}$ in $R(P)_{\mathfrak m}$. 
 \begin{Proposition}(Proposition 3.1 [Li])\label{Prop1}
 Suppose that $P\subset Q$ are prime ideals in $S$. Then $d_P(J)\le d_Q(J)$ and $d_P(J)=d_Q(J)$ implies $e_P(J)\le e_{Q}(J)$.
 \end{Proposition}
 
 The following corollary  follows from Proposition 3.1 [Li1], using generic freeness of $\mbox{gr}_I(R)$, a property which is discussed in the  sketch of the proof of Proposition 3.1 in [Li1]. In Theorem 3.4 [Sm], a result is proven from which this result for $I$-adic filtrations follows more directly. Upper semicontinuity of multiplicity (with the assumption that $d_P$ is constant) was earlier proven by Teissier for a coherent ideal sheaf  $I$ with finite cosupport over $Y$ ($\mathcal I$ is the $I$-adic filtration) in a reduced equidimensional analytic family $X\rightarrow Y$ in Proposition 3.1, Section 3, Chapter II [T1]. 
  We give a self contained  proof here. 
 
 \begin{Corollary}\label{Cor1} The map $P\mapsto (d_P(J),e_P(J))$ for $P\in \mbox{Spec}(S)$ is upper semicontinuous in the lex order.
 \end{Corollary}
 
 \begin{proof} To show that the function is upper semicontinuous, we must show that it is nondecreasing under specialization and is generically constant. It is nondecreasing under specialization by Proposition \ref{PropA}.
 
 For $P\in \mbox{Spec}(S)$, let 
 $$
 \lambda_{J,P}(n)=\dim_{k(P)}R(P)/J(P)^n.
 $$
 We have an expansion
 $$
 \lambda_{J,P}(n)=\sum_{\mathfrak m\in V(J(P))}[k(\mathfrak m):k(P)]\ell_{R(P)_{\mathfrak m}}(R(P)_{\mathfrak m}/(J(P)_{\mathfrak m})^n).
 $$
 For $n\gg 0$, $\lambda_{J,P}(n)$ is a polynomial of degree $d_P(J)$ and the coefficient of $n^{d_P(J)}$ is $\frac{e_P(J)}{d_P(J)!}$.
 
 We will prove that $P\mapsto (d_P(J),e_P(J))$ is generically constant. Let $\mathfrak p\in \mbox{Spec}(S)$ and let $Y=V(\mathfrak p)\subset \mbox{Spec}(S)$. We will show that the function is constant on a dense open subset of $Y$. 
 
 Let $\overline S=S/\mathfrak p$, $\overline R=R/\mathfrak p$ and $\overline J=J\overline R$. Let $Q\in Y$ and $\overline Q=Q/\mathfrak p\in \mbox{Spec}(\overline S)$. For all $n\ge 0$, we have
 $$
 R(Q)/J(Q)^n\cong (R/J^n)\otimes_Sk(Q)\cong (\overline R/\overline J^n)\otimes_{\overline S}k(\overline Q)\cong \overline R(\overline Q)/J(\overline Q)^n.
 $$
 Thus 
 \begin{equation}\label{eqC1}
 \dim_{k(Q)}R(Q)/J(Q)^n =\dim_{k(\overline Q)}\overline R(\overline Q)/J(\overline Q)^n.
 \end{equation}
 Let 
 $$
 G=\oplus_{n\ge 0}G_n={\rm gr}_{\overline J}(\overline R)=\oplus_{n\ge 0}\overline J^n/\overline J^{n+1}.
 $$
 $G_0=\overline R/\overline J$ is finite over $\overline S$ and $G$ is generated over $G_0$ be the finite $G_0$-module $\overline J/\overline J^2$.
 By generic freeness, after inverting some nonzero $s\in \overline S$,  $G_s$ is a free over $\overline S_s$. Each homogeneous component $(R_n)_s$ is thus flat over $\overline S_s$. From the exact sequences
 $$
 0\rightarrow \overline J^{n-1}/\overline J^n\rightarrow \overline R/\overline J^n\rightarrow \overline R/\overline J^{n-1}\rightarrow 0,
 $$
  we obtain inductively that $(\overline R/\overline J^n)_s$ is flat over $\overline S_s$. Since it is finite, it is locally free. Thus, using (\ref{eqC1}),
  $$
  \lambda_{J,Q}(n)=\dim_{k(Q)}R(Q)/J(Q)^n=\dim_{k(\overline Q)}\overline R(\overline Q)/\overline J(\overline Q)^n=\lambda_{\overline J,\overline Q}(n)
  $$ 
 is constant for $Q$ in a dense open subset of $Y$, simultaneously for every $n$. So the entire Hilbert function, and hence the pair $(d_Q(J),e_Q(J))$ is constant on a dense open subset of $Y$, establishing generic constancy, and completing the proof of the corollary.
 
  \end{proof}
 
 
 
 
 \subsection{Graded families}
 %Suppose that $I$ is an ideal in $S$ such that $R/I$ is a finite $S$-module and $I\cap S=0$. For a prime $P$ in $S$, let $I(P)=IR(P)$. 
%For large $n$, $\lambda_P(n):= \dim_{k(P)}(R(P)/I(P)^n)$ is a polynomial
%$$
%\lambda_P(n)=\frac{e_P(I)}{d_P(I)!}n^{d_P(I)}+\mbox{lower order terms in $n$}
%$$
%where 
%$$
%d_P(I)=\dim R(P)_{1+I(p)}=\max\dim R(P)_m
%$$
%with the max over the finite set of prime ideals $m$ in $V(I(P))\subset \mbox{MaxSpec}(R(P))$ and
%$$
%e_P(I)=\sum_{m\in V(I(P))}[k(m):k(P)]e(I(P)_m).
%$$
%Lipman proves the following in Proposition 3.1 [Li1].
 %\begin{Proposition}(Lipman) If $P\subset Q$ are two distinct prime ideals in $S$, then  either 
%$$
%\dim R(P)=d_P( I)<d_Q( I)
%$$
 %or 
 %$$
 %d_P( I)=d_Q( I)\mbox{ and }e_P(I)\le e_Q( I).
 %$$
%\end{Proposition}



Now let $\mathcal I=\{I_n\}$ be a graded family of ideals in $R$ such that $V(I_n)=V(I_1)$ for all $n$ (which always holds if $\{I_n\}$ is a filtration), $R/I_n$ are finite $S$-modules, $\phi^{-1}(I_n)=0$ for all $n$ and
 for $P$ a prime ideal in $S$, let $I_n(P)=I_nR(P)$ and
 $\mathcal I(P)=\{I_n(P)\}$, a graded family of ideals on $R(P)$. 
 We will use the following extension of definitions in [Li1] from an ideal to 
  a graded family of $m_R$-primary ideals.
 For $P\in\mbox{Spec}(S)$, let  
 $$
 \Sigma_P=\Sigma_P(\mathcal I)=V(I_1(P))=V(I_n(P))=\subset \mbox{MaxSpec}(R(P))
 $$ 
 for all $n>0$ and define 
 $$
 d_P(\mathcal I)=\max_{m\in \Sigma_P}\dim R(P)_m.
 $$
 Let
 $$
 \Sigma_P^{\rm top}=\Sigma_P(\mathcal I)=\{\mathfrak m\in \Sigma_P\mid \dim R(P)_{\mathfrak m}=d_P(\mathcal I)\}.
 $$ 
 Define 
 $$
 e_P(\mathcal I)=\sum_{m\in \Sigma_P^{\rm top}}[k(m):k(P)]e(\mathcal I(P)_m)
 $$
 where the multiplicity $e(\mathcal I(P)_m)$ is defined by Theorem \ref{Theorem1}.
 
 
 
 
 
 
 
 
 %where $e((I(P)_n)_m)$ is understood to be zero if $I(P)_n=R(P)_m$. 
 
 
\begin{Proposition} (Monotonicity)\label{PropA} If $P\subset Q$ are two distinct prime ideals in $S$, then  either 
$$
d_P(\mathcal I)<d_Q(\mathcal I)
$$
 or 
 $$
 d_P(\mathcal I)=d_Q(\mathcal I)\mbox{ and }e_P(\mathcal I)\le  e_Q(\mathcal I).
 $$
\end{Proposition}

In the case that $\mathcal I$  is an $I$-adic filtration $\mathcal I=\{I^n\}$ this is the statement of Proposition Proposition 3.1 of [Li1] (restated as Proposition \ref{Prop1} of this paper).
\begin{proof} Proposition \ref{Prop1} and the assumption that $\Sigma_P=(V(I_n(P))$ for all $n$ implies that 
$$
d_P(\mathcal I)=d_P(I_n)\le d_Q(I_n)=d_Q(\mathcal I)
$$
for all $n$. Suppose that $d_P(\mathcal I)=d_Q(\mathcal I)$. Then $d_P(I_n)=d_Q(I_n)$ for all $n$ and thus   
$$
\begin{array}{l}
\sum_{\mathfrak m\in \Sigma_P^{\rm top}}[k(\mathfrak m):k(P)]e(J(P)_{\mathfrak m})=e_P(I_n)\\
\le e_Q(I_n)=\sum_{\mathfrak m\in \Sigma_Q^{\rm top}}[k(\mathfrak m):k(Q)]e(J(Q)_{\mathfrak m})
\end{array}
$$
for all $n$ by Proposition \ref{Prop1}. Divide both sides of this equation by $n^{d_P(J)}$ and take the limit as $n$ goes to infinity. By Theorem \ref{Theorem1}, the limit exists, giving the conclusions of Proposition \ref{PropA}.

\end{proof}

\begin{Proposition}\label{PropB}(Upper semicontinuity of multiplicity) The map $P\mapsto (d_P(\mathcal I),e_P(\mathcal I))$ for $P\in \mbox{Spec}(S)$ is upper semicontinuous in the lex order.
\end{Proposition}

\begin{proof}
By assumption, $V(I_n(P))=\Sigma_P$ and so
\begin{equation}\label{eqB1}
d_P(I_n)=d_P(\mathcal I)
\end{equation}
for all $n$ and $P\in \mbox{Spec}(S)$. Now by [C2],
If $\mathcal K=\{K_n\}$ is a graded family of $m_A$-primary ideals in a Noetherian local ring $A$ of dimension $d$, then 
\begin{equation}\label{eqB3}
e(\mathcal K)=\lim_{n\rightarrow\infty}\frac{e(K_n)}{n^d}=\inf_{n\ge 1}\frac{e(K_n)}{n^d}.
\end{equation}
Since $[k(\mathfrak m):k(P)]\ge 0$ for all $\mathfrak m\in V(J(P))$ and $e(\mathcal I(P)_{\mathfrak m})\le \frac{e(I_n(P)_{\mathfrak m})}{n^{d_P(\mathcal I)}}$ 
for all $n$ and $\mathfrak m\in V(J(P))$ by equation (\ref{eqB3}) and (\ref{eqB1}),
\begin{equation}\label{eqB2}
e_P(\mathcal I)= \lim_{n\rightarrow \infty}\frac{e_P(I_n)}{n^{d_P(\mathcal I)}}=\liminf_{n\rightarrow \infty}\frac{e_P(I_n)}{n^{d_P(\mathcal I)}}. 
\end{equation}

By (\ref{eqB1}) and (\ref{eqB2}), given $(b,c)\in \Z_{\ge 0}\times \RR_{\ge 0}$,
$$
(d_P(\mathcal I),e_P(\mathcal I))<_{\rm  lex}(b,c)\mbox{ if and only if }(d_P(I_n),\frac{e_P(I_n)}{n^{d_P(I_n)}})<{\rm lex}(b,c)\mbox{ for some $n$}.
$$
Thus 
$$
\{P\mid (d_P(\mathcal I),e_P(\mathcal I))<_{\rm  lex}(b,c)\}=\cup_{n\ge 1}\{P\mid \frac{e_P(I_n)}{n^{d_P(I_n)}})<{\rm lex}(b,c)\}.
$$
Each set $\{P\mid \frac{e_P(I_n)}{n^{d_P(I_n)}})<{\rm lex}(b,c)\}$ is open by Corollary \ref{Cor1}. Thus the set 
$$
\{P\mid (d_P(\mathcal I),e_P(\mathcal I))<_{\rm  lex}(b,c)\}
$$
 is open, implying semicontinuity.
\end{proof}

With the above notation, for $P\in \mbox{Spec}(R)$ and $\mathfrak m\in \Sigma_P^{\rm top}$, let $\widetilde{\mathcal I(P)_{\mathfrak m}}$ be the saturation of $\mathcal I(P)_{\mathfrak m}$ defined in [C2]. 

\begin{Theorem}\label{TheoremC} Let assumptions be as above. Suppose that there exist constants $a\ge 0$ and $c\ge 0$ such that $d_P(\mathcal I)=a$ and $e_P(\mathcal I)=c$ for every $P\in \mbox{Spec}(S)$. 

Suppose that $h\in \sqrt{I_1}$ and there exists a Zariski open subset $U$ of $\mbox{Spec}(S)$ such that the natural image $h(P)$ of $h$ in $R(P)$ satisfies
$h(P)\in \widetilde{\mathcal I(P)}_1$ for all $P\in U$. Then 
\begin{equation}\label{eqT}
h(P)\in \widetilde{\mathcal I(P)}_1\mbox{ for all }P\in \mbox{spec}(S).
\end{equation}
\end{Theorem}

\begin{proof} Let $\mathcal J=\{J_n\}$ be the graded family of ideals on $R$ where $J_n=\sum_{i=0}^nh^iI_{n-i}$. $\mathcal I\subset \mathcal J$. Since $h\in\sqrt{I_1}$, $\sqrt{I_n}=\sqrt{I_1}$ and $I_n\subset J_n\subset \sqrt{I_1}$ for all $n$, we have that 
\begin{equation}\label{eqT1}
\sqrt{J_n}=\sqrt{I_n}=\sqrt{I_1}
\end{equation}
for all $n$.  For $P\in \mbox{Spec}(S)$, $\Sigma_P(\mathcal J)=\Sigma_P(\mathcal I)$, $\Sigma_P^{\rm top}(\mathcal J)=\Sigma_P^{\rm top}(\mathcal I)$ and $d_P(\mathcal J)=d_P(\mathcal I)$. Also $R/J_n$ is finite over $S$ for all $n$ since it is a quotient of $R/I_n$. Further, we have that  $\phi^{-1}(J_n)\subset \sqrt{0}$ by (\ref{eqT1}) and sincwe $\phi^{-1}(I_n)=0$.

For $P\in U$ and $\mathfrak m\in \Sigma_P^{\rm top}$, $\mathcal I(P)_{\mathfrak m}\subset \widetilde{\mathcal I(P)_{\mathfrak m}}$ by our assumption on $h$. Thus $e(\mathcal I(P)_{\mathfrak m})=e(\mathcal I(P)_{\mathfrak m})$ by Theorem 1.12 [C2]. Thus $e_P(\mathcal J)=c$ for all $P\in U$. 


Now fix $P\in \mbox{Spec}(S)$ and choose a minimal prime $\eta\subset P$. Since $U$ is dense and open in $\mbox{Spec}(S)$, we have that $\eta\in U$. 


We have that 
$\mbox{Spec}(S/\sqrt{0_S})=\mbox{Spec}(S)$, so Lipman's specialization inequality Proposition \ref{Prop1} applied to $J_n(R/\sqrt{0_S}R)$ for each $n$ gives the inequalities $d_{\eta}(J_n)\le d_P(J_n)$ since $d_{\eta}(J_n)=d_P(J_n)=a$. After dividing by $n^a$ and taking limits, we obtain that 
\begin{equation}\label{eqT2}
c=e_{\eta}(\mathcal J)\le e_P(\mathcal J).
\end{equation}
Since $I_n\subset J_n$ for all $n$,
\begin{equation}\label{eqT3}
e_P(\mathcal J)\le e_P(\mathcal I)=c.
\end{equation}
Thus by (\ref{eqT2}), equality holds in (\ref{eqT3}), and so $e(\mathcal J(P)_{\mathfrak m})=e(\mathcal I(P)_{\mathfrak m})$ for all $\mathfrak m\in \Sigma_P^{\rm top}$ since $e(\mathcal J(P)_{\mathfrak m})\le e(\mathcal I(P)_{\mathfrak m})$. 
By Theorem 1.12 [C2], the saturations of $\mathcal I(P)_{\mathfrak m}$ and $\mathcal J(P)_{\mathfrak m}$ are equal. Since $h\in J_1$, this implies (\ref{eqT}).

\end{proof}




\end{document}

\thispagestyle{empty}
%\centerline{Project Summary}
%\vskip .2truein
%\noindent{Overview}
%\vskip .2truein
%Problems on multiplicities of families of ideals and on  singularities will be investigated. Multiplicity is  the main invariant distinguishing different singularities and ideals. 

%\vskip.2truein
%\noindent{Intellectual Merit}
%\vskip .2truein
 %It is proposed to extend classical theorems on upper semi continuity of multiplicity  and specialization of integral dependence  to graded families of ideals. Extension of  theories of multiplicity  for non m-primary ideals to graded families of ideals we also be investigated. 

%The project seeks to construct a global resolution of singularities from the knowledge of good local uniformizations along all monomial valuations. Local uniformizations along such valuations always exist, even in positive characteristic, as they are Abhyankar valuations.  We also propose to give a short and accessible proof of embedded resolution of reduced excellent two dimensional schemes. We propose to give a counterexample to local monomialization in positive characteristic p in a degree $p^2$ extension. The role of valuation theory in these problems will be emphasized.

%The project will involve training and mentoring of graduate students. 
%\vskip .2truein
%\noindent{Broader Impacts}
%\vskip .2truein
%The proposer and Ana Bravo are organizing an SLMath Summer Graduate School on Resolution of Singularities, which will be held the two weeks of  August 2 - 13, 2027 at the Kavli Institute of the University of Tokyo in Kavli, Japan. SLMath will support students from US institutions to attend the summer school.

%The proposer will train and mentor graduates working towards a Ph.D.  Organization of conferences and writing expository articles is part of the project. 




%\vfill\eject


\thispagestyle{empty}
\centerline{\bf Results From Prior NSF Support}
\vskip .2truein
NSF award number: DMS-2348849, Amount: \$209,771

 Period of support: 08/15/2024 - 07/31/2027.

Title of project: Multiplicities and Valuations in Commutative Algebra.

\vskip .1truein


\noindent{\bf Intellectual Merit} 

We will discuss the results of two papers written during the period of the grant, [C1] and [C7].
Let $R$ be a $d$-dimensional Noetherian local ring with maximal ideal $m_R$. Definitions used here are presented in Section \ref{SecGrad} of the Project Description.

\subsection{Mutiplicities of graded families of $m_R$-primary ideals on an arbitrary Noetherian local ring.}


In [C1], we give a generalization of the multiplicity $e(I)$ of an $m_R$-primary ideal $I$ of $R$ to a multiplicity $e(\mathcal I)$ of a graded family of $m_R$-primary ideals $\mathcal I$ in $R$. 

This multiplicity gives the classical multiplicity $e(I)$ if $\mathcal I=\{I^n\}$ is the $I$-adic filtration, and agrees with the volume, $\displaystyle \lim_{n\rightarrow \infty}d!\frac{\lambda(R/I_n) }{n^d}$ for $R$ such  that the volume always exists as a limit. 

We  show in this paper that many of the classical theorems for the multiplicity of an ideal generalize to this multiplicity, including mixed multiplicities, the Rees theorem and the Minkowski inequality and equality. We give simple proofs which are independent of the theory of volumes and Okounkov bodies for all of our results, with the one exception being  the proof of the Minkowski equality. We do this by interpreting the multiplicity of graded families of $m_R$-primary ideals as a limit of  intersection products on the family of $R$-schemes which are obtained by blowing up $m_R$-primary ideals in $R$.

The volume 
$\mbox{vol}(\mathcal I)$ of a graded $m_R$-family $\mathcal I$ is defined as
\begin{equation}\label{eq34}
\mbox{vol}(\mathcal I)=\limsup_{n\rightarrow\infty} d!\frac{\lambda(R/I_n)}{n^d}.
\end{equation}







This limsup is always a nonnegative real number (since there exists a positive integer $s$ such that $m_R^{sn}\subset I_n$ for all $n$).

The problem of existence of such limsups (\ref{eq34}) as limits  has been considered by Ein, Lazarsfeld and Smith [ELS] and Musta\c{t}\u{a} [Mus].
When the ring $R$ is a domain and is essentially of finite type over an algebraically closed field $k$ with $R/m_R=k$, Lazarsfeld and Musta\c{t}\u{a} [LM] showed that
the limit exists for all graded $m_R$-families.  We have the following general theorem.
\begin{Theorem}(Theorem 1.1 [C16])\label{TheoremVolEx}
Suppose that $N(\hat R)$ is the nilradical of the $m_R$-adic completion $\hat R$ of $R$. Then the limit
\begin{equation}\label{eq10}
\lim_{n\rightarrow\infty} d!\frac{\lambda_R(R/I_n)}{n^d}
\end{equation}
exists for any graded $m_R$-family $\mathcal I=\{I_n\}$ on $R$ if and only if $\dim N(\hat R) < d$.
\end{Theorem}
 The condition $\dim N(\hat R)<d$ holds for instance if $R$ is analytically unramified ($\hat R$ is reduced). It was shown by Hailong Dao and Ilya Smirnov that if $\dim N(\hat R)=d$ then there exists a graded $m_R$-family $\mathcal I$ on $R$ such that $\mbox{vol}(\mathcal I)$ is not a limit.
 %The proof of Theorem \ref{TheoremVolEx} uses the theory of Okounkov bodies, which was introduced in the papers \cite{Ok}, \cite{KK} and \cite{LM}.
 
 The following ``Volume = Multiplicity'' formula is proven in the generality stated below in [C14]. It was earlier proven  in some cases in [ELS] and  when the ring $R$ is a domain that is essentially of finite type over an algebraically closed field $k$ with $R/m_R=k$ by   Lazarsfeld and Musta\c{t}\u{a} in [LM]. 
 %The proof of Theorem \ref{volmult} uses the theory Okounkov bodies.
  
  \thispagestyle{empty}  
 \begin{Theorem}\label{volmult}(Theorem 1.1 [C14])  Suppose that $\dim N(\hat R)<d$ and $\mathcal I=\{I_n\}$ is a graded $m_R$-family on $R$. Then 
\begin{equation}\label{eq33}
\mbox{\rm vol}(\mathcal I)=\lim_{p\rightarrow\infty}\frac{e(I_p)}{p^d}.
\end{equation}
\end{Theorem}
Let $R$ be a Noetherian local ring and $\mathcal I$ be a graded $m_R$-family. We define the multiplicity $e(\mathcal I)$ of 
$\mathcal I$ as
$$
e(\mathcal I)=\limsup_{p\rightarrow\infty}\frac{e(I_p)}{p^d}.
$$

Although the volume $\mbox{vol}(\mathcal I)$ does not  exist in general as a limit on an arbitrary Noetherian local ring $R$,  the multiplicity $e(\mathcal I)$ always does exist. 

\begin{Theorem}(Theorem 1.3 [C1])\label{Theorem1}
  Let $R$ be a $d$-dimensional Noetherian local ring and $I = \{I_n\}$ be a graded family of $m_R$-primary ideals. Then the multiplicity
$$
e(\mathcal I)=\lim_{n\rightarrow\infty}\frac{e(I_n)}{n^d}
$$
exists as a limit.
\end{Theorem}

\thispagestyle{empty}


We give  a simple direct proof of Theorem \ref{Theorem1} in [C1], using limits of intersection products. This proof  does not use the theories of Volume or of Okounkov bodies.
The beginning point of this proof is the following interpretation of the multiplicity of an $m_R$-primary ideal as an intersection product.  


\begin{Theorem}\label{Theorem20}(Theorem A, [CM1])
Let $R$ be a $d$-dimensional Noetherian local ring  and $I\subset R$ be an $m_R$-primary ideal. 
 Let $\pi:Y\rightarrow \mbox{Spec}(R)$ be a birational proper morphism such that $I\mathcal O_Y$ is invertible. Then 
\begin{equation}
e(I)=-((I\mathcal O_Y)^d).
\end{equation}
\end{Theorem}

 In [Rm] Ramanujam obtained the above formula when there is a good intersection product defined on birational proper morphisms to $\mbox{Spec}(R)$. In [Laz1], the formula is stated for the
case that R is the local ring of a closed point on an algebraic scheme over an algebraically closed
field and in [BLQ] the formula is proven when R is an analytically irreducible domain and Y is normal. The fact that the good intersection product exists above arbitrary Noetherian local rings $R$ follows from 
the theory of rational equivalence developed in [Th] as applied to the intersection theory developed in [Fl] and [Kl].  This general intersection theory is developed in [CM1], and it is used to give a geometric interpretation of Rees' degree functions, originally developed by 
Rees in  [Re2] and [Re3].
%A survey of intersection theory above a noetherian local ring is given in section \ref{SecInt}.



Here is a short outline of our proof of Theorem \ref{Theorem1}.
Let $Y_n$ be the blowup of the ideal $I_n$ above $\mbox{Spec}(R)$. Then $I_n\mathcal O_{Y_n}=\mathcal O_{Y_n}(-
F_n)$ for some effective Cartier divisor $F_n$ on $Y_n$ with exceptional support. We compute $e(I_n)=-((-F_n)^d)$, the self intersection of $-F_n$ on $Y_n$.
Using only the fact that  $\mathcal I$ is a graded $m_R$-family, we show as a consequence of  the more general Theorem 2.3 [C1], that 
\begin{equation}\label{eqIn11}
\lim_{n\rightarrow\infty}\frac{e(I_n)}{n^d}=\lim_{n\rightarrow \infty}-\left(\left(\frac{-F_n}{n}\right)^d\right)
\end{equation}
exists, establishing Theorem \ref{Theorem1}.

 We always have that
 $$
 {\rm vol}(\mathcal I)\le e(\mathcal I)
 $$
 as is shown in Proposition 4.1 [C1]. However, there exist examples where the inequality is strict, as is shown in Example 4.2 [C1].
 In particular, the  volume = multiplicity formula of Theorem \ref{volmult} for rings $R$ such that $\dim N(\hat R)<d$ does not extend to more general rings.
  
  
  We establish the following generalization of Rees' theorem (Rees
[Re], Theorem 11.3.1[SH) for $m_R$-primary ideals to arbitrary graded $m_R$-families.  The saturation $\tilde{\mathcal I}$ of a graded family of $m_R$-primary ideals $\mathcal I$ is defined in Definition \ref{SatDef}.
The essential case of analytically irreducible local rings in [BLQ].  
  \begin{Theorem}\label{Theorem14*}(Theorem 1.12 [C1])  Suppose that $R$ is a Noetherian local ring and $\mathcal I\subset \mathcal J$ are graded $m_R$-families. Then $e(\mathcal I)=e(\mathcal J)$ if and only if $\tilde{\mathcal I}=\tilde{\mathcal J}$. In particular, $\tilde{\mathcal I}$ is the largest $m_R$-filtration $\mathcal J$ such that $\mathcal I\subset \mathcal J$ and $e(\mathcal J)=e(\mathcal I)$.
\end{Theorem}

\thispagestyle{empty}
We also establish the Minkowski inequalities and characterization of equality in the Minkowski equality, generalizing classical results for $m_R$-primary ideals
(Teissier [T1], [T3], [T4], Rees and Sharp [RS], D. Katz [Ka]).

Let $\mathcal I=\{I_n\}$ be a graded $m_R$-family. 

\begin{equation}\label{eqIn3}
\mbox{Given $c\in \RR_{>0}$, let $\mathcal I^{(c)}$ be the family $\{I_{\lceil cn\rceil}\}$,}
\end{equation}
where $\lceil cn\rceil$ is the round up of the real number $cn$. If $c=0$, then $\mathcal I^{(0)}$ is the filtration such that $(I^{(0)})_n=m_R$ for $n>0$.

\begin{Theorem}\label{Minkgen*} Suppose that $R$ is a $d$-dimensional Noetherian local ring with $d\ge 2$ and $\mathcal I$, $\mathcal J$ are graded $m_R$-families such that 
$e(\mathcal J)>0$. 
Then there is equality in the Minkowski inequality
$$
e(\mathcal I\mathcal J)^{\frac{1}{d}}= e(\mathcal I)^{\frac{1}{d}}+e(\mathcal J)^{\frac{1}{d}}
$$
if and only if $\tilde{\mathcal I}=\widetilde{(\mathcal J^{(c)})}$ for some $c\in \RR_{\ge 0}$.
\end{Theorem}


The theorem relies essentially on  Theorem 10.3 [C12].
  
  
  %A graded $m_R$-filtration $\mathcal I$ is bounded if its integral closure $\overline{\mathcal I}$ is a divisorial filtration. The $I$-adic filtration of an $m_R$-primary ideal $I$ on a formally equidimensional local ring is an example of a  bounded filtration. 


  
  %\begin{Theorem}\label{BoundMink}(Proven in Theorem 1.4 \cite{C8} when $R$ is analytically irreducible  or an excellent local domain)
  %Suppose that $R$ is a $d$-dimensional Noetherian  local ring of dimension $d\ge 2$ and $\mathcal I$ and $\mathcal J$ are $\QQ$-bounded $m_R$-filtrations.
%Then the  Minkowski equality
%$$
%e(\mathcal I\mathcal J)^{\frac{1}{d}}=e(\mathcal I)^{\frac{1}{d}}+e(\mathcal J)^{\frac{1}{d}}
%$$
%holds if and only if there exist positive integers a,b such that there is equality of integral closures
%$\overline{\mathcal I^{(a)}}=\overline{\mathcal J^{(b)}}$.
%\end{Theorem}


%%%%%%%%%%%%%

\subsection{ McAdam's theorem for divisorial filtrations}

In this subsection, we summarize a theorem from [C7].

 If $I$ is an ideal in $R$, then the Rees algebra of $I$ is the graded $R$-algebra
$R[It]=\sum_{n\ge 0}I^nt^n$ and the analytic spread of $I$ is 
$$
\ell(I)=\dim R[It]/m_RR[It].
$$ 
%The classical theory of the analytic spread of an ideal is expounded in [SH]. 

%It is natural to extend the definition of analytic spread of an ideal to arbitrary graded filtrations of $R$. A graded filtration $\mathcal I=\{I_n\}$ of $R$ is a filtration of $R$ such that   $I_mI_n\subset I_{m+n}$ for all $m,n$.
Let $\mathcal I=\{I_n\}$ be a graded filtration on  $R$. The Rees algebra of  $\mathcal I$ is $R[\mathcal I]=\oplus_{n\ge 0}I_n$ and the analytic spread of the filtration $\mathcal  I$ is defined to be
\begin{equation}\label{eqN10}
\ell(\mathcal I)=\dim R[\mathcal I]/m_RR[\mathcal I].
\end{equation}
Thus if $I$ is an ideal in $R$ and $\mathcal I=\{I^n\}$ is the $I$-adic filtration, then $\ell(I)=\ell(\mathcal I)$.

We always have (Lemma 3.6 [CSa2]) that 
\begin{equation}\label{eqZ1}
\ell(\mathcal I)\le \dim R,
\end{equation}
so the classical inequality $\ell(I)\le \dim R$ of ideals (Proposition 5.1.6 [SH]) extends to arbitrary graded filtrations.


We have the following remarkable theorem by McAdam.

\begin{Theorem}\label{TheoremC4}([McA], Theorem 5.4.6 [HS]) Let $R$ be a formally equidimensional local ring and $I$ be an ideal in $R$. Then $m_R\in \mbox{Ass}(R/\overline{I^n})$ for some $n$ if and only if $\ell(I)=\dim(R)$.
\end{Theorem}




To see how Theorem \ref{TheoremC4} can be extended, we  consider divisorial filtrations. Divisorial filtrations are defined in Subsection \ref{SubSecVal} of the Project Description. %The following theorem is a generalization of Theorem \ref{TheoremC4} to divisorial filtrations. 
	If $\mathcal I=\{I_n\}$ is a divisorial filtration on an excellent local ring $R$, then $I_n=\overline{I_n}$ for all $n$ and $R$ is formally equidimensional if $R$ is equidimensional (Scholie IV.7.8.3 (x) [EGAIV]).
	
\begin{Theorem}\label{Theorem3}(Theorem 1.5 [C7]) Let $R$ be an equidimensional local ring which is essentially of finite type over an excellent local ring of dimension less than or equal to 3. Let $\mathcal I=\{I_n\}$ be a $\QQ$-divisorial filtration on $R$. Then the following are equivalent.
\begin{enumerate}
\item[1)] The analytic spread of $\mathcal I$ is $\ell(\mathcal I)=\dim R$.
\item[2)] There exists $n_0\in \ZZ_{>0}$ such that $m_R\in \mbox{Ass}(R/I_n)$ if $n\ge n_0$.
\item[3)] $m_R\in \mbox{Ass}(R/I_{m_0})$ for some $m_0\in \ZZ_{>0}$.
\end{enumerate}
\end{Theorem}

 Theorem \ref{Theorem3} is proven for excellent equicharacteristic zero local domains in Theorem 1.4 [C5], using the existence of resolutions of singularities of excellent equicharacteristic zero domains ([H1] or [Te]). 
 %The proof in [C5] shows that Theorem \ref{Theorem3} is true  for excellent local domains of dimension $\le 3$, using the existence of resolutions of singularities of  excellent domains of dimension $\le 3$ proven in [CoP]. 
 The function (\ref{ex2*}) and a generalization of the theory of $\sigma$-decomposition of pseudo-effective divisors of [Na] to nonsingular projective schemes which birationally dominate $\mbox{Spec}(R)$ play an essential role in the proof in [C5]. 

The generalization of Theorem 1.4 [C5] to Theorem \ref{Theorem3} in [C7] uses
Temkin's theorem on existence of regular alterations, Theorem 1.2.5 [Te1]. Temkin's theorem shows that if $X$ is an integral scheme which is of finite type over an excellent domain of dimension less than or equal to three, then there exists a regular alteration $X'\rightarrow X$. %Alterations and regular alterations are defined in the following definition.
 
 %\begin{Definition}\label{AltDef}(Definition 2.20 [dJ])
%Let $X$ be a Noetherian integral scheme. An Alteration $X'$ of $X$ is an integral scheme $X'$ with a morphism $\phi:X'\rightarrow X$ which is dominant, proper and such that for some open $U\subset X$, $\phi^{-1}(U)\rightarrow U$ is finite.
%\end{Definition}

%A scheme $X$ is nonsingular if $\mathcal O_{X,q}$ is a regular local ring for all $q\in X$. An alteration $\phi:X'\rightarrow X$ is regular if $X'$ is nonsingular. 





The conclusions of 3) implies 1) of Theorem \ref{Theorem3} do not hold for more general filtrations. In fact, there exist $\RR$-divisorial filtrations $\mathcal I$ on an excellent equicharacteristic zero local domain such that 
the conclusions of Theorem \ref{Theorem3} are false, as shown in Example 1.5 [C5].

%The existence of regular alterations over a local domain which is essentially of finite type over a field was first proven by de Jong
%in Theorem 4.1[dJ].
\vskip .2 truein
\noindent{\bf Broader impacts}

\thispagestyle{empty}
The KUMUNU Conference on Commutative Algebra was coorganized by the proposer with  Ian Aberbach, Rankeya Datta and Hema Srinivasan In Fall of 2023, 2024 and 2025 (the 2024 and 2025 meetings were during the time of this grant).   Each year, there were about 100 participants, mostly graduate students, postdocs and faculty from colleges and universities in the mid midwest. These  conferences were  partially supported by a conference grant from NSF.
The KUMUNU conference is a regional conference which alternates every three years between the Universities of Missouri, Kansas and Nebraska. It brings in some national and international speakers in addition to speakers from the three hosting universities. 

Stephen Landsittel earned his Ph.D. at the University of Missouri, under the direction of the proposer, in April 2025. He is currently a PostDoc at  SEAS (School of Engineering and Applied Sciences) at Harvard University and Hebrew University of Jerusalem, working with Eran Nevo. During this time period, Mifron Fernandez began working with the proposer towards a Ph.D. The PI meets every week individually with graduate students working with him. 


The proposer wrote an expository article on deeply ramified fields,  ``Almost mathematics, K\"ahler differentials and deeply ramified fields'', which is  to appear in ``The landscape of commutative algebra'' of the Lecture Note Series of the London Math. Soc.,  listed as [C3] in publications acknowledging this award.  The article  includes an exposition of  material on almost mathematics  from the book ``Almost Ring Theory'' by Gabber and Romero, using only elementary methods of commutative algebra.

Coorganizer, since October 2022,  of The Fellowship of the Ring, a virtual seminar in commutative algebra, with
Michael  Brown, David Eisenbud, Irena Peeva, Claudia Polini, Karl Schwede, Gregory Smith and Bernd Ulrich.

The PI has been an editor of the Journal of Algebra since 2005.


\vskip .1truein


\noindent{\bf Publications resulting from the NSF award}\thispagestyle{empty}









\begin{enumerate}

\item[{[C0]}] S.D. Cutkosky, Limits of Amao multiplicities and epsilon multiplicity,\linebreak arXiv:2608.12238v2.

\item[{[C1]}] S.D. Cutkosky,  Multiplicities of graded families of ideals on Noetherian local rings, arXiv:2603.06844.

\item[{[CK]}] S. D. Cutkosky and F.-V. Kuhlmann,, K\"ahler differentials of extensions of valuation rings and deeply ramified fields,  arXiv:2306.04967.

\item[{[CM1]}] S.D. Cutkosky and  J. Monta\~no, Degree functions of graded families of ideals,  arXiv:2503.14429, to appear in Collectanea Mathematica.

\item[{[C2]}] S.D. Cutkosky, Asymptotic behavior of Hilbert functions of divisors above a singularity'' to appear in a volume in honor of Bernard Teissier.

\item[{[C3]}] S.D. Cutkosky, Almost mathematics, K\"ahler differentials and deeply ramified fields, arXiv:2310.09581, to appear in ``The landscape of commutative algebra'', edited by M. Rossi, N.V. Trung, B. Urich and J. Verma, Lecture Note Series of the London Math. Soc., Cambridge Univ. Press.

\item[{[C4]}] S.D. Cutkosky,  The Proj of the Rees Algebra of a graded family of ideals, J. Pure Appl. Algebra 230 (2026).

\item[{[C5]}]  S.D. Cutkosky, Analytic spread and associated primes of divisorial filtrations,  Proceedings of the AMS 154 (2026), 2721 - 2732.

\item[{[AC]}] R. Ahmadian and S.D. Cutkosky, Generating sequences of valuations on simple extensions of domains,   Michigan Math. Journal, 76 (2026), 689 - 705.

\item[{[CM2]}] S.D. Cutkosky and  J. Monta\~no, Multiplicities and degree functions in local rings via intersection products,  
J. London Math. Soc 112 (2025).


\item[{[C6]}] S.D. Cutkosky,  Rees algebras and the reduced fiber cone of divisorial filtrations on two dimensional normal local rings, 
 Comm. in Algebra 53 (2025), 3364 - 3387.

\item[{[C7]}] S.D. Cutkosky, The Rees algebra and analytic spread of a divisorial filtration,  Adv. Math. 479 (2025).

\item[{[CL]}] S.D. Cutkosky and S. Landsittel,  Epsilon multiplicity is a limit of Amao multiplicities,  Journal of Algebra and its applications 24 (2025), 13 -14.


\item[{[CKR]}] S.D. Cutkosky, F.-V. Kuhlmann and A. Rzepka, On the Computation of K\"ahler Differentials and Characterization of Galois extensions with independent defect,   Mathematische Nachrichten 298 (2025), 1549 - 1577.

\item[{[CSa1]}] S.D. Cutkosky and P. Sarkar, Epsilon multiplicity and analytic spread of filtrations,    Illinois Journal of Math. 68 (2024), 189 - 210.



\item[{[C8]}] S.D. Cutkosky,  A counterexample to strong local monomialization in a tower of two independent defect Artin-Schreier extensions, Annales de la Facult\'e des Sciences Toulouse 33 (2024) 915 - 935.

\item[{[CPr]}] S.D. Cutkosky and S. Praharaj,  The Asymptotic Samuel Function of a Filtration, Acta Mathematica Vietnamica 49 (2024),  61 - 81.

\end{enumerate}

\vfill\eject

\noindent{\bf Project Description}
\vskip .1truein
\thispagestyle{empty}
\section{Graded families of ideals}\label{SecGrad}
Sections \ref{QSec1} and \ref{QSec2}  of this proposal discuss questions about graded families of ideals. We review some definitions and basic properties here.   $\lambda(N)$ will denote the length of a module $N$ over a ring $R$.

\begin{Definition} Suppose that $R$ is a ring. A graded family of ideals of $R$ is a family $\mathcal I=\{I_n\}_{n\in \NN}$ of ideals $I_n$ in $R$ such that $I_0=R$ and $I_mI_n\subset I_{m+n}$ for all $m,n\in \NN$. 
\end{Definition}

A graded family of ideals $\mathcal I$ is a filtration if $I_m\subset I_{m+1}$ for all $m$. The case of a single ideal $I$ is realized as the $I$-adic filtration $\mathcal I=\{I^n\}$.


\subsection{Valuations, integral closure and saturation}\label{SubSecVal}

Let $R$ be a $d$-dimensional Noetherian local ring. A valuation $v$ of $R$ is defined  in Definition VI.3.1, page 386 [Bou]
 as follows. There is a prime ideal $P$ of $R$ and a valuation $v$ of the quotient field of $R/P$ such that $v$ is nonnegative on $R/P$. We will restrict to minimal primes $P$ of $R$. Elements of $P$ have value $\infty$, which is larger than any element of the value group of $v$.
A valuation $v$ of $R$ is 
 a divisorial valuation of $R$ if letting $P$ be the preimage of the maximal ideal of the valuation ring of $v$ in $R$, $\mbox{trdeg}_{k(P)}k(v)=\mbox{ht} P-1$ (Definition 9.3.1 [SH]). The value group of a divisorial valuation is $\ZZ$. A valuation $v$ is called an $m_R$-valuation (as defined by Rees) if $v(m_R)>0$ and $v$ is a divisorial valuation of $R$.
 
If $R$ is an analytically unramified local domain, such as a complete local domain, then the divisorial valuations of $R$ are the valuations of $R$
 such that the valuation ring $\mathcal O_v$ is essentially of finite type over $R$ (this follows from  Theorem 9.3.2 [SH]). 

 If $R$ is a Noetherian local ring and $v$ is a divisorial valuation of $R$,  we define a graded filtration by $\mathcal I(v)=\{I(v)_n\}$ where 
 \begin{equation}\label{eqvalfin}
 I(v)_n=\{ f\in R\mid v(f)\ge n\}.
 \end{equation}
  A (real) divisorial filtration is a graded filtration $\mathcal I=\{I_n\}$ where there exist divisorial valuations $v_1,\ldots,v_s$ of $R$ and $a_1,\ldots,a_s\in \RR_{\ge 0}$ such that $I_n=\{f\in R\mid v_i(f)\ge na_i\mbox{ for }1\le i\le s\}$. 



Define $v(\mathcal I)$ for $\mathcal I$ a graded filtration of $m_R$-primary ideals and a divisorial valuation  $v$ of $R$ by
\begin{equation}\label{ex2*}
v(\mathcal I)=\inf_{m\in \ZZ_{>0}}\frac{v(I_m)}{m}\in \RR_{\ge 0}.
\end{equation}
 %It is   shown in [JM] that 
%$v(\mathcal I)=\displaystyle\lim_{m\rightarrow\infty}\frac{v(I_m)}{m}$. 



If $I$ is an ideal in a local ring $R$, then $\overline I$ denotes its integral closure. 

Suppose that $R$ is a local ring and $\mathcal I=\{I_n\}$ is a graded family of ideals on $R$. The proof in Lemma 5.6 [C12] extends immediately 
to show that the integral closure $\overline{\sum_{n\ge 0}I_nt^n}$ of $\sum_{n\ge 0}I_nt^n$ in $R[t]$ is the graded $R$-algebra 
$\overline{\sum_{n\ge 0}I_nt^n}=\sum_{n\ge 0}I^*_nt^n$,  where $\overline{\mathcal I}=\{I^*_n\}$ is the filtration
$$
I^*_n=\{f\in R\mid f^r\in \overline{I_{rn}}\mbox{ for some }r>0\}.
$$
If $\mathcal I=\{I^n\}$ is the adic-filtration of the powers of a fixed ideal $I$, then $I^*_n=\overline{I^n}$ for all $n$. However, it is possible for the integral 
closure $\overline{\mathcal I}$ of a graded family $\mathcal I=\{I_n\}$ to strictly contain the graded family $\{\overline{I_n}\}$ of integral closures of the ideals $I_n$.
 

%\begin{Lemma}\label{LemmaIC2} Suppose that $R$ is a local ring, $\mathcal I=\{I_n\}$ is a filtration of $R$ and $p\in \mbox{Spec}(R)$. Let $\overline{R[\mathcal I]}=\oplus_{n\ge 0}K_n$. Then the integral closure $\overline{\sum_{n\ge 0}I_nR_pt^n}$ of $\sum_{n\ge 0}I_nR_pt^n$ in $R_p[t]$ is $\sum_{n\ge 0}K_nR_pt^n$.
%\end{Lemma}


The following definition is made in Definition 3.1 [BLQ] on  analytically irreducible local rings $R$ for graded families of $m_R$-primary ideals. 
 We extend the definition to graded families of $m_R$-primary ideals on Noetherian  local rings.
\begin{Definition}\label{SatDef}(Definition 1.7 [C1]) Let $R$ be a Noetherian local ring  and $\mathcal I$ be a graded family of $m_R$-primary ideals on $R$. The saturation $\tilde{\mathcal I}$ of $\mathcal I$ is the filtration
$\tilde{\mathcal I}=\{\tilde I_n\}$ defined by 
$$
\tilde I_n=\{f\in m_R\mid v(f)\ge n v(\mathcal I)\mbox{ for all divisorial valuations $v$ of $R$}\}.
$$
\end{Definition}

It is shown in Lemma 5.2 [C1] that
the graded $R$-algebra $R[\tilde{\mathcal I}]=\sum_{n\ge 0}\tilde I_nt^n$ is integrally closed in $R[t]$.
  We have inclusions 
$$
\mathcal I\subset \{\overline{I_n}\}\subset\overline{\mathcal I}\subset \tilde{\mathcal I}
$$
where it may be that all of the inclusions are proper. 

                    
Let $I$ be an $m_R$-primary  ideal of $R$ and $\mathcal I=\{I^n\}$ be the $I$-adic filtration. Rees showed that 
there are  discrete valuations $v_1,\ldots,v_r$ dominating $R$  and $a_1,\ldots,a_r\in\ZZ_{>0}$ such that 
$$
I^*_n=\overline{I_n}=\{f\in R\mid v_i(f)\ge na_i\mbox{ for }1\le i\le r\}
$$
for all $n\ge 0$. The valuations $v_i$ are  called the Rees valuations of $I$. If $R$ is formally equidimensional, then the Rees valuations are all divisorial  (as follows from Proposition 10.2.6 [SH])
and $\tilde{\mathcal I}=\overline{\mathcal I}$ is thus a divisorial filtration.

If $R$ is not formally equidimensional, then there exist $m_R$-primary ideals $I$ of $R$ such that the filtration
 $\{\overline{I^n}\}$ is not real divisorial, which will occur if the Rees valuations of $I$ are not all divisorial valuations. 



It is more difficult to find examples where $\tilde{\mathcal I}$ is not real divisorial on a formally equidimensional local ring. 
We give such an example in Section 6 [C1],  showing that there exists a saturated filtration on a two dimensional regular local ring which is not real divisorial. 






\section{Multiplicities of  graded families of ideals on families of schemes}\label{QSec1}
\thispagestyle{empty}

The questions of this section are an application of the general multiplicity theory defined and developed in  [C1], which is valid for 
graded families of $m_R$-primary ideals in a Noetherian local ring. This multiplicity was defined in Theorem \ref{Theorem1}.

We propose to extend some classical theorems of Teissier on an ideal on  a family of analytic spaces to graded families of ideals on families of schemes. It is expected that these questions will have  positive answers, by combining techniques of [C1] and [Li1]. 



In this section, we  have the following assumptions, generalizing  those of Section 3 of [Li1] from a single ideal $I$ to a graded family $\mathcal I$. 

Let $\phi:S\rightarrow R$ be a homomorphism of commutative Noetherian local rings.   For every prime ideal $P$ in $S$, let $k(P)$ be the fraction field
 of $S/P$ and  set $R(P)=R\otimes_Sk(P)$. 
 %Suppose that $I$ is an ideal in $S$ such that $R/I$ is a finite $S$-module and $I\cap S=0$. For a prime $P$ in $S$, let $I(P)=IR(P)$. 
%For large $n$, $\lambda_P(n):= \dim_{k(P)}(R(P)/I(P)^n)$ is a polynomial
%$$
%\lambda_P(n)=\frac{e_P(I)}{d_P(I)!}n^{d_P(I)}+\mbox{lower order terms in $n$}
%$$
%where 
%$$
%d_P(I)=\dim R(P)_{1+I(p)}=\max\dim R(P)_m
%$$
%with the max over the finite set of prime ideals $m$ in $V(I(P))\subset \mbox{MaxSpec}(R(P))$ and
%$$
%e_P(I)=\sum_{m\in V(I(P))}[k(m):k(P)]e(I(P)_m).
%$$
%Lipman proves the following in Proposition 3.1 [Li1].
 %\begin{Proposition}(Lipman) If $P\subset Q$ are two distinct prime ideals in $S$, then  either 
%$$
%\dim R(P)=d_P( I)<d_Q( I)
%$$
 %or 
 %$$
 %d_P( I)=d_Q( I)\mbox{ and }e_P(I)\le e_Q( I).
 %$$
%\end{Proposition}



Now let $\mathcal I=\{I_n\}$ be a graded family of ideals in $R$ such that $V(I_n)=V(I_1)$ for all $n$ (which always holds if $\{I_n\}$ is a filtration), $R/I_n$ are finite $S$-modules, $I_n\cap S=0$ for all $n$ and
 for $P$ a prime ideal in $S$, let $I_n(P)=I_nR(P)$ and
 $\mathcal I(P)=\{I_n(P)\}$, a graded family of ideals on $R(P)$. We will use the following extension of definitions in [Li1] from an ideal to 
  a graded family of $m_R$-primary ideals.
 
 

 For $P\in\mbox{Spec}(S)$, let  $\Sigma_P=V(I_1(P))\subset \mbox{MaxSpec}(R(P))$ and define 
 $$
 d_P(\mathcal I)=\max_{m\in \Sigma_P}\dim R(P)_m,
 $$
 $$
 e_P(\mathcal I)=\sum_{m\in \Sigma_P, \dim R(P)_m=d_P(\mathcal I)}[k(m):k(P)]e(\mathcal I(P)_m)
 $$
 where the multiplicity $e(\mathcal I(P)_m)$ is defined by Theorem \ref{Theorem1}.
 %where $e((I(P)_n)_m)$ is understood to be zero if $I(P)_n=R(P)_m$. 
 
 The philosophy of the following questions is that asymptotically, Rees valuations of the blowups of the $I_n$ whose center on $S$ is not an irreducible component of $S$, should determine a jump in $e_P(\mathcal I)$.

\begin{Question} (Monotonicity)\label{QuestionA} If $P\subset Q$ are two distinct prime ideals in $S$, is it true that  either 
$$
d_P(\mathcal I)<d_Q(\mathcal I)
$$
 or 
 $$
 d_P(\mathcal I)=d_Q(\mathcal I)\mbox{ and }e_P(\mathcal I)\le  e_Q(\mathcal I)?
 $$
\end{Question}

In the case that $\mathcal I$  is an $I$-adic filtration $\mathcal I=\{I^n\}$ this is the statement of Proposition 3.1 of [Li1] and is proven there.
\begin{Question}\label{QuestionB}(Upper semicontinuity of multiplicity) Is the map $P\rightarrow (d_P(\mathcal I),e_P(\mathcal I))$ for $P\in \mbox{Spec}(S)$ upper semicontinuous in the lex order?
\end{Question}

In the case that $\mathcal I$ is an $I$-adic filtration $\mathcal I=\{I^n\}$ the answer is yes. In fact this follows from Proposition 3.1 [Li1], using generic freeness of $\mbox{gr}_I(R)$, a property which is discussed in the  sketch of the proof of Proposition 3.1 in [Li1]. In Theorem 3.4 [Sm], a result is proven from which this result for $I$-adic filtrations follows more directly. 

Upper semicontinuity of multiplicity (with the assumption that $d_P$ is constant) was earlier proven by Teissier for a coherent ideal sheaf  $I$ with finite cosupport over $Y$ ($\mathcal I$ is the $I$-adic filtration) in a reduced equidimensional analytic family $X\rightarrow Y$ in Proposition 3.1, Section 3, Chapter II [T1].

To formulate the next question, we must generalize the notion of saturation of a graded family from the case of the saturation of a graded family of $m_R$-primary ideals, defined in Definition \ref{SatDef}. Let $R$ be a Noetherian ring and let $\mathcal I=\{I_n\}$ be a graded family of ideals on $R$ such that $V(I_n)=V(I_1)$ for all $n$. Define the saturation of $\mathcal I$ to be $\tilde{\mathcal I}=\{\tilde I_n\}$ where
$$
\tilde I_n=\{f\in \sqrt{I_1}\mid v(f)\ge  n v(\mathcal I)\mbox{ for all divisorial valuations v whose center lies in $V(I_1)$}\}.
$$
If $\mathcal I=\{I^n\}$ is the $I$-adic filtration then $\tilde I_n=\overline{I^n}$ for all $n$ and if $\mathcal I$ is a divisorial filtration then $\tilde{\mathcal I}=\mathcal I$.



We need to add an equidimensionality condition on the fibers. Namely, we must assume that $\dim R(P)_m=d_P(\mathcal I)$ for all $P\in\mbox{Spec}(R)$ and $m\in \Sigma_P$. This is implied by the condition that $R(P)$ is of finite type over $k(P)$ and equidimensional of the same dimension for all $P$.

\begin{Question}\label{QuestionC} Suppose that $d_P(\mathcal I)$ and $e_P(\mathcal I)$ are constant on $\mbox{Spec}(S)$. Assume that $\dim R(P)_m=d_P(\mathcal I)$ for all $P\in\mbox{Spec}(R)$ and $m\in \Sigma_P$. 
Let $h\in R$ be such that for all $P$ in a Zariski open dense subset of $Y$, the restriction of $h$ to $R(P)$ is in the saturation $\widetilde{\mathcal I(P)}$ of $\mathcal I(P)$. Then is $h$ in the saturation $\tilde{\mathcal I}$ of $\mathcal I$?
\end{Question}

The Rees type  theorem for saturation of graded families of $m_R$-primary ideals, Theorem \ref{Theorem14*},  should play a major role in a proof affirming Question \ref{QuestionC}.

\thispagestyle{empty}
Question \ref{QuestionC} has a positive answer when $\mathcal I=\{I^n\}$ is an $I$-adic filtration on an equidimensional family of reduced equidimensional analytic spaces $Y$, as is shown by Teissier [T1], Appendix 1 of [T2] and in Section (1.3) of Gaffney and Gassler [GG]. In the case of an $I$-adic filtration on $Y$, $\mathcal I=\{I^n\}$, the saturation $\mathcal I=\{\overline{I^n}\}$ is the filtration of integral closures of powers of $I$. These papers in fact show that if $h$ is in the integral closure $\overline{I(P)}$ for all $P$ in a Zariski open dense subset of $Y$ then $h\in \overline I$. Thus showing ``Specialization of integral dependence''.  The proof uses valuation theoretic methods, which generalize integral closure to saturation in the case of general graded families of ideals. An algebraic proof of this statement (for the case of ideals) is obtained from Theorem 4b of Section 3 and the Scholium on page 135 of [Li1]. 

It is expected that it will not be difficult to find an example of a graded family of ideals for which ``specialization of integral dependence'' fails. For instance there is an example of this type for Rees' theorem for multiplicities of graded families after Theorem 6.9  [CSS].

%\begin{Question} \label{QuestionD} Assume that Question \ref{QuestionC} has a positive answer. Can Rees' theorem for graded families of $m_R$-primary ideals, Theorem \ref{Theorem14*}, be deduced as a corollary of the conclusions of Question \ref{QuestionC}?
%\end{Question}

%The proof would be by taking the graded family $\{(I_n+tJ_n)T$ on $T=R[t]$, regarded as an extension of $R$. This is the deformation used (for ideals $I$ and $J$) in the proof of Corollary 4.9 [GG], [PTUV] and [CPU]. 




 \section{Epsilon multiplicities and Amao multiplicities}
 \thispagestyle{empty} 
 Epsilon multiplicity is defined for ideals, modules and graded algebras by Bernd Ulrich and Javid Validashti in their paper [UV]. A motivation is to find criteria for 
 equisingularity in families of arbitrary singularities, extending work of Teissier, Gaffney, Kleiman and others. 
 
 
 In their paper, they define the epsilon multiplicity as a limsup, which is always finite since it is bounded from above by the $j$-multiplicity.
 They show that if $R$ is a Noetherian local ring of dimension $d$ which is equidimensional and universally catenary,  and $U\subset E$ are rank $e$ submodules of a free module of finite rank $e$, then  $E$ is integral over $U$ if and only if there is equality of $\epsilon$ multiplicities 
 $\epsilon(U_q)=\epsilon(E_q)$ for all prime ideals $q$ of $R$. 
 
 The epsilon multiplicity can be an irrational number, even for an ideal in a regular local ring, as shown in [CHST].
 
 An important question in this theory is if the epsilon multiplicity is actually a limit, and not just a limsup. We first consider the case of ideals (rank $e=1$).
 
  
 Let $R$ be a $d$-dimensional Noetherian local ring with maximal ideal $m_R$ and $I\subset R$ be an ideal. The epsilon multiplicity of $I$ (the case $e=1$) is defined in [UV] as 
$$
\epsilon(I)=
\limsup_{n\rightarrow\infty} \frac{\ell_R((I^n)^{\rm sat}/I^n)}{n^d/d!},
$$
where $\ell_R(M)$ is the length on an $R$-module $M$.
The saturation $J^{\rm sat}$ of an ideal $J$ of $R$ is 
$$
J^{\rm sat}=J:m_R^{\infty}=\cup_{i=0}^{\infty}J:m_R^i
$$
where $m_R$ is the maximal ideal of $R$.

We have recently shown that the $\epsilon$-multiplicity  always exists as a limit for all ideals in an arbitrary Noetherian local ring $R$ in [C0].

\begin{Theorem}\label{limthm}([C0])
Suppose that $R$ is a $d$-dimensional Noetherian local ring and $I$ is an ideal of $R$. Then
the epsilon multiplicity of an ideal in $R$ exists as a limit; that is, 
$$
\epsilon(I)=
\lim_{n\rightarrow\infty} \frac{\ell_R((I^n)^{\rm sat}/I^n)}{n^d/d!}.
$$
\end{Theorem}


Theorem \ref{limthm} was earlier proven in Corollary 6.3  [C16] for ideals in a $d$-dimensional unramified local ring $R$.  This proof uses Okounkov body methods, which are not applicable for rings whose nilradical has dimension $d$.

We propose the following question.

\begin{Question} \label{QE} Does the epsilon multiplicity exist as a limit  for modules and graded algebras  over an arbitrary Noetherian local ring?
\end{Question} 

Theorem \ref{limthm} shows that this question has a positive answer for ideals. In [sD], Suprajo Das shows that the question 
has a positive answer for reduced standard graded algebras over an excellent local ring. This last paper uses Okounkov body methods, which are not applicable for algebras which are not close to being reduced. 

Our method of investigation of Question \ref{QE} is to combine the methods of [C0] and [sD]. 

The proof in [C0] is accomplished by defining another multiplicity, $a\epsilon(I)$, as a limsup of normalized Amao multiplicities. 
The  Amao multiplicity is defined and developed in [Am], [R5], [KR]  and on page 332 of [SH].
We show in [C0] that $a\epsilon(I)$ always exists as a limit. We use the fact  shown in [CL] with Stephen Landsittel, that if $I$ is a ideal in an analytically unramified local ring $R$,
  then $\epsilon(I) =a\epsilon(I)$. 
We then show that $a\epsilon(I)$ satisfies an associativity formula and show that  the corresponding formula is an inequality for $\epsilon(I)$. Then using the fact that $a\epsilon(I)\le\epsilon(I)$ we deduce that $\epsilon(I)=a\epsilon(I)$. We then further deduce that $\epsilon(I)$ is itself a limit, proving Theorem \ref{limthm}.

The Amao multiplicity is defined for graded algebras and modules. We propose to develop the theory of $a\epsilon$ for modules and graded algebras, and compare with $\epsilon$-multiplicity in our investigation of Question \ref{QE}.





%It is shown in Corollary 6.3 \cite{C2}, using Okounkov body techniques,  that if $R$ is analytically unramified, then the epsilon multiplicity of an ideal in $R$ exists as a limit; that is, 
%$$
%\epsilon(I)=
%\lim_{n\rightarrow\infty} \frac{\ell_R((I^n)^{\rm sat}/I^n)}{n^d/d!}.
%$$
%However, the epsilon multiplicity can be an irrational number, as shown in \cite{CHST}.


%Some recent papers on epsilon multiplicity are Suprajo Das \cite{sD}, Cutkosky and Sarkar \cite{CS}, Suprajo Das, S. Dubey, S. Roy and J. Verma \cite{DDRV} and Suprajo Das, S. Roy and V. Trivedi \cite{DRT}, Cutkosky and Landsittel \cite{CL}, Landsittel \cite{L}.

%The  Amao multiplicity is defined and developed in \cite{Am}, \cite{R} and on page 332 of \cite{HS}. Suppose that $R$ is a $d$-dimensional Noetherian local ring and $K\subset J$ are ideals in $R$ such that $J/K$ has finite length. Then $J^n/K^n$ has finite length for all $n\ge 0$ and there exists a numerical polynomial $p(n)$ such that for $n$ sufficiently large, $p(n)=\ell_R(J^n/K^n)$. The Amao multiplicity $a(K,J)$ is defined to be $d!$ times the coefficient of the degree $d$ term of $p(n)$. That is,
%$$
%a(K,J)=\lim_{n\rightarrow\infty}\frac{\ell_R(J^n/K^n)}{n^d/d!}.
%$$
%The Amao multiplicity is always a nonnegative integer.

%In  Theorem 1.5 \cite{L},  it is  proven that the epsilon multiplicity of an ideal in a Noetherian $d$-dimensional local ring is a limit of Amao multiplicities, if the dimension of the nilradical of $\hat R$ is less than the dimension of $R$.

%\begin{Theorem}\label{TheoremA}(Theorem 1.5 \cite{L})  Let $R$ be a  $d$-dimensional Noetherian local ring and let $I\subset R$ be an ideal. 
%Suppose that the dimension of the nilradical of $\hat R$ is less than $d$.
%Then
%$$
%\epsilon(I)=
%\lim_{m\rightarrow\infty}\frac{a(I^m,(I^m)^{\rm sat})}{m^d}.
%$$
%\end{Theorem}

%Theorem \ref{TheoremA} 
 %is  a generalization of  the volume equals multiplicity formula of ordinary multiplicity to epsilon multiplicity. The volume = multiplicity formula of ordinary multiplicity is  proven in increasing generality in Ein, Lazarsfeld and Smith \cite{ELS}, Musta\c{t}\u{a} \cite{M1}, Lazarsfeld and Musta\c{t}\u{a} \cite{LM} and  Cutkosky, Theorem 6.5 \cite{C2} and Theorem 1.4 \cite{C4}. Recently, volume = multiplicity formulas have been proven for $p$-families and weakly graded families, by Sudipta Das \cite{SD}, Sudipta Das and Cheng Meng \cite{DM} and Sarkar \cite{PS}.
 
 %The proof of Theorem \ref{TheoremA} in \cite{L} is by reduction to the case that $R$ is analytically unramified. The case when $R$ is analytically unramified is proven in Theorem 1.1 \cite{CL} using Okounkov body methods.

%In our recent paper ****, we show that the limit 
%$$
%\lim_{m\rightarrow\infty}\frac{a(I^m,(I^m)^{\rm sat})}{m^d}
%$$
% exists generally in rings which are essentially of finite type over a field, even when they are not generically reduced, as is required in the equality of Theorem \ref{TheoremA}. 

%\begin{Theorem}\label{TheoremAmlim}(Theorem \ref{Theorem3} of Section \ref{SecAmlim})
%Let $R$ be a $d$-dimensional local ring which is essentially of finite type over a field $k$ and $I\subset R$ be an ideal. Then the limit
%$$
%\lim_{n\rightarrow\infty} \frac{a(I^n,(I^n)^{\rm sat})}{n^d}\in \RR
%$$
%exists. 
%\end{Theorem}



%It follows from Theorem 5.7 \cite{KT} that Amao multiplicities $a(I,J)$ can be computed using intersection theory, specifically as a sum of suitable Segre products obtained by blowing up $I$ and $J$. This ex  
 
 
 
 
 
%\begin{Question}\label{Question2} Does epsilon multiplicity always exist as a limit for ideals in a Noetherian local ring? If so, is it a limit of Amao multiplicities?
%\end{Question}

%\begin{Question}\label{Question3} Does the limit of Amao multiplicities of $(I^m)^{\rm sat}/I^m$ exist for an ideal $I$ in an arbitrary Noetherian local ring?
%\end{Question}

%\section{Some questions on Segre numbers of graded families of ideals}\label{QSec2}

%The generalization of the multiplicity of an $m_R$-primary ideal to an arbitrary ideal is the set of Segre numbers.

%In Section 2 of the paper [GG] of Gaffney and Gassler, Segre cycles are defined for ideal sheaf germs $I$ on a germ of a reduced equidimensional analytic space $X$ of dimension $d$, as are their Segre numbers $e_k(I,X)$. They show that these numbers can be computed from intersection products on the blowup of $mI$ where $m$ is the ideal of the point of the germ. Their definitions and the proofs of these formulas extend with a little modification to the case of a Noetherian local ring. In Corollary 4.9 to their Theorem 4.7, ``The principal of specialization'', they show that if $I\subset J$ then $e_k(I,X)=e_k(J,X)$ for $1\le k\le d$ if and only if $\overline I=\overline J$, giving a generalization of Rees' theorem to ideals which are not $m_R$-primary, in a reduced equidimensional analytic local ring.

%The multiplicity sequence of an ideal $I$ in a $d$-dimensional Noetherian local ring $R$
%$$
%(c_0(I),c_1(I),\ldots,c_d(I))
%$$
% is defined by Achilles and Manaresi [AM]. Although given a different algebraic formulation,  they show in [AM2] that these are the same as the Segre numbers of [GG]. We will write $c_i$ for the Segre numbers. This theory is developed in [PTUV] and [CPU]. The  Rees theorem in  [GG] is extended  to ideals in Noetherian local rings in Theorem 4.4 [PTUV].
 
%\begin{Question}\label{Question4} Does there exist a good theory of multiplicities for graded families of  ideals (which are not maximal) in an analytically unramified  Noetherian local ring, or even in a reduced local ring which is essentially of finite type over a field?
%\end{Question}  
%\thispagestyle{empty}
% The questions which we ask are only interesting in the case that the Rees algebra $\oplus_{n\ge 0}I_n$ is not a finitely generated $R$-algebra. If it is finitely generated, everything can be reduced to the (known) case of an $I$-adic filtration. 


%We present  our recent progress on this problem and ask some questions. 

%If $I$ is a fixed ideal and $\mathcal I=\{I^n\}$ is the $I$-adic filtration of $R$ then $c_i(I^n)=n^ic_i(I)$, so it is natural to define the Segre numbers of a graded family of ideals on $R$ as 
%$$
%c_i(\mathcal I)=\limsup_{n\rightarrow \infty}\frac{c_i(I_n)}{n^i}.
%$$

%Now $c_1(\mathcal I)$ is generally very well behaved. If we assume that  all irreducible components of $X=\mbox{Spec}(R)$ are not contained in $V(I_n)$ for $n\gg 0$, then
%$c_1(\mathcal I)$ exists as a  finite limit
%\begin{equation}\label{S2}
%c_1(\mathcal I)=\lim_{n\rightarrow \infty}\frac{c_1(I_n)}{n}.
%\end{equation}
%This follows from Fekete's lemma. There exist examples where the limit does not exist, if an irreducible component of $\mbox{Spec}(R)$ is contained in $V(I_n)$ for infinitely many $n$ and is also not contained in $V(I_n)$ for infinitely many $n$. This however is a situation which we can avoid just by making the reasonable assumption on the graded family that $V(I_n)$ never contains a minimal prime of $R$, which we will do for the rest of this discussion.

%However, things get much more wild for $i>1$. In fact, given an increasing function $\sigma(n)$, there exists a filtration $\{I_n\}$ in a two-dimensional regular local ring such that 
%\begin{equation}\label{S1}
%\frac{c_2(I_n)}{\sigma(n)}\mapsto \infty
%\end{equation}
%as $n\rightarrow \infty$. Using formulas (3.4) of [GG] (which extend to Noetherian local rings), we get examples of regular local rings of dimension $d$ for any $d\ge 2$ such that the conclusions of (\ref{S1}) hold. 

%Further, there are examples of filtrations on regular local rings $R$ of dimension 2 such that  $\limsup_{n\rightarrow \infty}\frac{c_2(I_n)}{n^2}$ is finite but the limit $\lim_{n\rightarrow \infty}\frac{c_2(I_n)}{n^2}$ does not exist. Again, this yields examples of this type on regular local rings of arbitrary dimension $\ge 2$. 

%It is tempting to try to use techniques from [C1], involving intersection theory, to try to prove existence of limits.  The exceptional divisor of the blow up of 
%$I_n$ is a Cartier divisor $D_n$ which has a decomposition $D_n=D_n^{X\setminus 0}+D_n^0$, where $D_n^{X\setminus 0}$ is the cycle of components of $D_n$ which do not contract to $m_R$ and $D_n^0$ is the cycle of components of $D_n$ which contract to $m_R$.


%We know from [C1] that there is a finite limit if we are taking intersections of families of divisors $\{G_n\}$ all of which are nef (or antinef) and satisfy inequalities 
%\begin{equation}\label{S3}
%G_n+G_m\le G_{n+m}
%\end{equation}
% (on a suitable birational model). However, each of these inequalities can fail for $G_n=D_n^{0}$. In dimension 2, $D_n^{0}$ is always antinef, but (\ref{S3}) can fail, which is the reason for the existence of the example giving (\ref{S1}). In dimension 3 and higher,  $D_n^{0}$ and $D_n^{X\setminus 0}$ can be neither nef nor antinef, so there is a further obstruction to the existence of a limit. Of course the limit always exists for $c_1$, even if there is failure of nefness. The proof of this equation uses algebraic methods and not geometric methods.

%While the proposer has not worked out examples, Segre numbers $c_i$ get more complicated as $i$ increases, so it is expected that this bad behavior continues to hold for $i> 2$ and arbitrary graded families, and it should not be very difficult to find examples of this type. 

%\thispagestyle{empty}
%We  ask if the limit exists for divisorial filtrations, an important ``geometric'' case. 
%\begin{Question}\label{QuestionS1} Suppose that $R$ is an  equidimensional excellent local ring and $\mathcal I=\{I_n\}$ is a divisorial filtration of ideals on $R$. Then are the Segre numbers $c_i(\mathcal I)$ always less than infinity? Is
%$$
%c_i\mathcal (I)=\lim_{n\rightarrow\infty}\frac{c_i(I_n)}{n^i}
%$$
%for $1\le i\le d$?
%\end{Question}
%Of course we already know that the limit exists and is finite for $i=1$, so the question is really for $i\ge 2$.



%Some positive evidence exists, as we have shown that Question \ref{QuestionS1} has a positive answer if $R$ has dimension two and $R$ is normal. 

%A positive answer to Question \ref{QuestionS1} should follow from a positive answer for the case when $R$ is normal.
%Divisorial filtrations really live on $R$ mod its nilradical, so we may assume that $R$ is reduced. Techniques used in Section 5.3 [C12] to reduce related results on the normalization of a ring to the ring itself could be used here.

 %The proposer will use Okounkov body methods and other methods, such as from intersection theory, to address this question. 

%Regardless of the answer to Question \ref{QuestionS1}, 
%we can ask if the Rees type theorem of Theorem 4.4 [PTUV] extends to divisorial filtrations.

%\begin{Question} Suppose that $R$ is an excellent equidimensional local ring and $\mathcal I=\{I_n\}\subset \mathcal J=\{J_n\}$ are divisorial filtrations. Are the conditions 
%$$
%\overline{ I^n}=\overline{J^n}\mbox{ for all $n\ge 0$}
%$$
 %and 
 %$$
 %c_i(\mathcal I)=c_i(\mathcal J)\mbox{ for $0\le i\le d$}
 %$$
 % equivalent?
%\end{Question}

%At least in dimension 2, and for $R$ normal, where we know the Segre numbers exist as finite limits, it is expected that the question has a positive answer. Methods of intersection theory should be applicable to address this. 







\section{Local uniformization of Abhyankar valuations and Global resolution of singularities after a reembedding and a toric morphism}

\begin{Question} \label{Question5} Suppose that $R$ is a reduced excellent local ring and $\mbox{Spec}(R)$ has a local uniformization for all Abhyankar valuations of $R$. Does there exist a proper birational morphism $X\rightarrow\mbox{Spec}(R)$ such that $X$ is nonsingular?
\end{Question}
\thispagestyle{empty}

Question \ref{Question5} is expected to have a positive answer for quasi-ordinary singularities, as we now outline. 

A characteristic zero hypersurface germ  with equation f = 0 in $k^n$ is quasi-ordinary if there exists a finite projection to $k^{n-1}$
 such that the discriminant locus is a germ of a normal crossing divisor ([Li2]).


Let $R$ be a Noetherian local ring and $v$ be a valuation of $R$. In general, the valuation group $\Gamma_v$ of $v$ can be very big. For instance, even on an algebraic function field it can be $\QQ^n$ with the lex order for any $n$. The valuation ideals of $v$ in $R$ are defined for $\gamma\in \Gamma_v$ by
$$
I(v)_{\gamma}=\{f\in R\mid v(f)\ge \gamma\}
\mbox{ and }
I(v)^+_{\gamma}=\{f\in R\mid v(f)>\gamma\}.
$$
The associated graded ring $\mbox{gr}_v(R)$ is the graded ring
$$
 \mbox{gr}_v(R)=\oplus_{\gamma\in \Gamma_v}I(v)_{\gamma}/I(v)^+_{\gamma}.
 $$
 In general $\mbox{gr}_v(R)$ is not Noetherian, even when $\Gamma_v=\ZZ$.

Let $R=k[x_1,\ldots,x_n]/(f)$ where $k$ is an algebraically closed field and 
$$
f=x_n^d+a_1x_n^{d-1}+\cdots +a_d
$$
 with all $a_i\in k[x_1,\ldots,x_{n-1}]$ is irreducible. The algorithm of [CMT] constructs the associated graded ring 
 $\mbox{gr}_v(R)$ of a valuation dominating $R$, as long as the characteristic of $k$ does not divide $d$. $\mbox{gr}_v(R)$  is always a finitely generated $k$-algebra with this restriction. This algorithm has been extended to arbitrary ground fields by Razieh Ahmedian and the proposer in [AC].

The region $\RR^n$ parametrizes monomial valuations. It is divided into different chambers where the algorithm produces the same ring $\mbox{gr}_v(R)$, but with different gradings.  



With Hussein Mourtada and Bernard Teissier, we have recently shown that with the additional assumption that $R$ is a quasi-ordinary singularity, there is only one chamber, and we can use the common ring $\mbox{gr}_v(R)$ to construct an embedding of $\mbox{Spec}(R)$ into $\AA^m_k$ for some $m$ so that the embedded singularity is Newton nondegenerate; that is, the intersection of the singular locus $\mbox{Sing}(V(\mbox{In}_v(I))\cap (K^*)^m$ is empty for all monomial valuations $v$, so the singular locus is disjoint from the torus of $\AA^m$. Here $\mbox{In}_v(I)$ is the ideal of initial forms of $I$ with respect to the valuation $v$. We can then apply the theorem of Aroca-G\'omez-Morales and Mourtada [AGM] (generalizing results of Khovanskii, Kouchnirenko [Kou] and Varchenko [V]) to resolve $\mbox{Spec}(R)$ after one toric blowup of $\AA_k^m$.


\begin{Question}\label{Question6} Can the above proposed proof of resolution of quasi-ordinary singularities be extended to arbitrary singularities $R=k[x_1,\ldots,x_n]/(f)$?
\end{Question}

The significance of this is that local uniformization (resolution of singularities along a valuation) is true for Abhyankar valuations as shown in [KK] and [C11] and  monomial valuations are Abhyankar valuations. However, local uniformization is not known for arbitrary valuations. The essential obstruction to local uniformization is defect and Abhyankar valuations do not have defect [Ku]. A positive answer to Question \ref{Question6} would show that it suffices to only consider the better behaved Abhyankar valuations 
for resolution of singularities. 

\section{An easily  accessible proof of embedded resolution of surface singularities}\label{SecRes}

In this section, we address the following question.

\begin{Question}\label{Question7} Can a relatively simple, generally accessible and short proof of resolution and embedded resolution of singularities of noetherian reduced excellent 2 dimensional schemes be given?
\end{Question}

Such a proof would be of general interest since there are many applications of this result in different fields of mathematics. Earlier the proposer gave in [C18] an accessible and short proof of resolution of singularities of projective threefolds over an algebraically closed field of characteristic $p>5$. This gives a much simpler proof of Abhyankar's resolution theorem in [Ab3].


A scheme $X$ is regular if all local rings $\mathcal O_{X,Q}$ of $X$ are regular local rings. 

\thispagestyle{empty}
\begin{Definition} Let $X$ be a reduced noetherian and excellent scheme. A resolution of singularities of $X$ is an everywhere regular scheme $X'$ and a surjective proper birational morphism $\pi:X'\rightarrow X$.
\end{Definition}

\begin{Definition} Let $X$ be a reduced closed subscheme of a regular noetherian excellent scheme $Y$. An embedded resolution of singularities of $X$ is a proper birational morphism $\pi:Y'\rightarrow Y$ such that $Y$ is regular and $\pi^{-1}(X)$ is a simple normal crossings divisor. 
\end{Definition}






The first proof of resolution of singularities of a noetherian reduced excellent finite type scheme over a field of characteristic zero was by Hironaka [H1]. His proof is in fact  for schemes of finite type over an excellent local ring of equicharacteristic zero. The most general possible statement in equicharacteristic zero, resolution of singularities of noetherian reduced quasi-excellent schemes of equicharacteristic zero, was proven by Temkin [Te]. These papers also prove embedded resolution.

The first proof of nonembedded resolution of  singularities for a projective surface over an algebraically closed  field of characteristic zero is by Abhyankar [Ab2]. Abhyankar gave the first proof of embedded resolution of singularities for a projective surface over an algebraically closed field of positive characteristic in [Ab3]. The first edition of this book appeared in 1966. Most of the local calculations for the proof are given in a series of earlier papers. A survey of his proof is given in [C17].

A proof of nonembedded resolution of noetherian excellent surface singularities is given by Lipman in [Li3]. 

Hironaka gave a very terse outline of a proof of embedded resolution of an irreducible surface embedded in a nonsingular three dimensional variety over an algebraically closed field in [H2]. At the end of the proof, he makes some comments explaining that the proof extends to prove embedded resolution of an excellent surface. In Chapter 7 of [C21], the proposer gives a proof where the details of Hironaka's proof are worked out in the case that Hironaka addresses in his notes: embedded resolution of an irreducible surface embedded in a nonsingular three dimensional variety over an algebraically closed field. In [CJS] Cossart, Jannsen and Saito give a proof of embedded resolution of a noetherian excellent two dimensional scheme, using techniques of Hironaka. However, the proof is extremely long (247 pages), partly because it  proves  many results in the full generality of arbitrary dimension, in the hope that these methods will be useful in proving resolution in all dimensions. 

The strongest established result on resolution of singularities in positive characteristic is the theorem of Cossart and Piltant [CoP]. They proved that resolution of singularities is true for reduced excellent schemes of dimension $3$.

The goal of this project is to simplify the proof in [CJS], limiting the development to dimension two,  and perhaps use other methods in the goal of giving a much shorter and more accessible proof giving a positive answer to Question \ref{Question7}.

In  equicharacteristic zero  there is the miracle of the existence locally of hypersurfaces of maximal contact (first noticed by Abhyankar), which allows an induction argument in resolution on dimension. This fails completely in residue characteristic $p>0$, and is the reason that resolution in positive characteristic is so difficult.  

The starting point in this project is the proof in Chapter 7 of [C21], which gives a positive answer to the question for surfaces locally embedded in a nonsingular threefold over an algebraically closed field of characteristic zero.   Some of the machinery developed by Hironaka in [H1] Is applicable here, namely the use of the Hilbert Samuel function to give a suitable stratification of the singular locus of the surface. We also need to use Hironaka's theory of the characteristic polyhedra [H3], and directrix, and their behavior under blowing up.  
\section{Failure of local monomialization in positive characteristic}
\thispagestyle{empty}

A fundamental problem in algebraic geometry is to find a general notion of simple morphism, and to then establish that it is possible to simplify a dominant morphism $\phi:X\rightarrow Y$ of varieties by performing  suitable monoidal transforms (blow ups of nonsingular subvarieties) above $X$ and $Y$ to obtain such a simple morphism.  

A local formulation of this problem can be given by asking for an appropriate simple form of the extension of local rings $\mathcal O_{Y,q}\rightarrow \mathcal O_{X,p}$ if $p\in X$ and $\phi(p)=q$. The notion of simplification after performing sequences of monoidal transforms above $\mathcal O_{Y,q}$ and $\mathcal O_{X,p}$ can then be expressed by fixing a valuation dominating $\mathcal O_{X,p}$. Then the valuation will dominate unique local rings on any sequences of monoidal transforms of $X$ and $Y$. 

A possible notion of a local simple morphism is that of monomial inclusion of regular local rings. We call such a map a monomial extension. Simplification by performing monoidal transforms is  called local monomialization. In characteristic zero, local monomialization is a good definition of simplification 
as a local monomialization can always be achieved along any valuation ([C22] and [C20]). In [C19] a global form of local monomialization is established for morphisms of projective 3-folds in characteristic zero.

In characteristic $p>0$, it is shown in [CP] that local monomialization is true for generically finite maps of surfaces over an algebraically closed field, if the extension of valuation rings is defectless (a condition which is always true in characteristic zero). It is further shown in [C13] that this result holds for arbitrary defectless extensions of two dimensional regular local rings. 

It is natural to ask if local monomialization can always be achieved in positive characteristic.

\begin{Question}\label{Question8} Can an example of a degree $p^2$ extension of two dimensional regular local rings of characteristic $p>0$ be given above which local monomialization fails?
\end{Question}

\thispagestyle{empty}
This is the smallest degree in which it is possible for local monomialization to fail. In [C16], an example is given of a  degree $p^2+1$ extension $R\rightarrow S$ of two dimensional regular local rings where local monomialization fails. However, it is desirable to have an example of degree $p^2$ in which it is clear how defect determines the failure of local monomialization. 

Suppose that $R\rightarrow S$ is an extension of two dimensional regular local rings and $v$ is a zero dimensional valuation which dominates $S$.
By Abhyankar's theorem [Ab1], given a valuation $v$ dominating $S$, the set of all regular local rings which birationally dominate $S$ and are dominated by $v$ is exactly the chain 
\begin{equation}\label{eqM1}
S=S_0\rightarrow S_1\rightarrow  \cdots S_n\rightarrow S_{n+1}\rightarrow \cdots
\end{equation}
 of the blowups at the center of the valuation $v$ on each $S_n$. The zero dimensional condition just means that the center of $v$ on $S_n$ is always the maximal ideal. Similarly, the set of 
all regular local rings which birationally dominate $S$ and are dominated by $v$ is exactly the chain 
\begin{equation}\label{eqM2}
R=R_0\rightarrow R_1\rightarrow  \cdots R_n\rightarrow R_{n+1}\rightarrow \cdots.
\end{equation}
Thus $R\rightarrow S$ does not have a local monomialization along $v$ means exactly that there is never a factorization $R_i\rightarrow S_j$ which has a monomial form.

 The natural way to construct such an example is within a chain of two defect Artin Schreier extensions. It is known that the sequences (\ref{eqM1}) and (\ref{eqM2}) are very wild within defect Artin Schreier extensions ([CP],   [C9] and [C10]), and this wildness is understood as a consequence of the defect. 

In the example  in [C15], it is shown that there is an $n_0$ such that there are nonmonomial factorizations $R_{mn_0}\rightarrow S_{mn_0}$ for all $m$ 
but there is no factorization $R_{mn_0+r}\rightarrow S_{(m+1)n_0}$ for all $m$ and $1<r<n_0$.
This is accomplished by showing that there is periodicity in the sequences $R_{mn_0}\rightarrow R_{(m+1)n_0}$ and $S_{mn_0}\rightarrow S_{(m+1)n_0}$
which have different but the same ``type'' of equations. This is done by ad hoc calculations which do not exhibit the defect. Only after the fact can it be determined using ramification theory that the defect of the extension is $p^2$. Looking more closely at the extension, one sees that  in $R\rightarrow \overline R$, the integral closure of $R$ in the quotient field of $S$, the maximal ideal of $R$ splits into two points above $m_R$. The one whose local ring is $S$, and another one which is actually an unramified extension. Thus the map on completions $\hat R\rightarrow \hat S$ is a degree $p^2$ extension, which has $p^2$ as its defect. However, it would be nice to make this extension explicit. Probably it is an extension of two defect Artin-Schreier extensions, and it would be interesting to interpret in terms of the theory of defect Artin-Schreier extensions. Further, it would be nice to give an ``algebraic'' example. 



\section{Use of AI}  
The proposer proposes to use  AI for preliminary research. 
%Some of the examples of bad behavor of limits of Segre numbers given in Section \ref{QSec2} were found as a result of conversations with ChatGPT. 
The proposer is using  version 6, which is a fairly high level version. The proposer asked ChatGPT  to explain parts of the paper [GG] by Gaffney and Gassler to him (specifically 
Section 4 on ``Specialization of Integral Closure'').  After a little going back and forth, it provided  very clear proofs and explanation of the paper, which the proposer finds absolutely amazing. The proposer is now preparing with AI assistance a survey article ``Rees like theorems in analytic geometry and commutative algebra'' comparing techniques of analytic geometry and commutative algebra in Rees like theorems characterizing equality of integral closure of ideals. 
\vskip .1truein
\noindent{\bf Broader Impacts}
\thispagestyle{empty}

The proposer and Ana Bravo are organizing an SLMath Summer Graduate School on Resolution of Singularities, which will be held the two weeks of  August 2 - 13, 2027 at the Kavli Institute of the University of Tokyo in Kavli, Japan. SLMath will support students from US institutions to attend the summer school.
Daily lectures will be held by the Proposer and Ana Bravo. 

The PI intends to be active in organization of workshops and conferences during the period of this proposal. 

Currently, Mifron Fernandez and Austin Spratt are working with the proposer towards a Ph.D at the University of Missouri. The PI meets every week individually with graduate students working with him. 

The proposed simplified proof of embedded resolution of excellent surface singularities in Section \ref{SecRes} would be useful to many mathematicians. 

The proposer is preparing  a survey article ``Rees like theorems in analytic geometry and commutative algebra'' comparing techniques of analytic geometry and commutative algebra in Rees like theorems characterizing equality of integral closure of ideals. 


Coorganizer, since October 2022,  of The Fellowship of the Ring, a virtual seminar in commutative algebra, with
Michael  Brown, David Eisenbud, Irena Peeva, Claudia Polini, Karl Schwede, Gregory Smith and Bernd Ulrich.


The PI has been an editor of the Journal of Algebra since 2005.




 
\vfill\eject
\end{document}

\thispagestyle{empty}
\begin{enumerate}\thispagestyle{empty}


\item[{[Ab1]}] S. Abhyankar, On the valuations centered in a local domain, Amer. J. Math. 78 (1956), 321 - 348.

\item[{[Ab2]}] S. Abhyankar, Local uniformization of algebraic surfaces over ground fields of characteristic $p\ne 0$, Annals of Math. 63 (1956), 491 - 526.

\item[{[Ab3]}] S. Abhyankar, Resolution of singularities of embedded algebraic surfaces, second edition, Springer Verlag, 1998.

\item[{[AC]}] R. Ahmadian and S.D. Cutkosky, Generating sequences of valuations on simple extensions of domains, to appear in Michigan Math. J., arXiv:1904.10702.

\item[{[AM]}] R. Achilles and M. Manaresi, Multiplicities of a bigraded ring and intersection theory, Math. Ann. 309 (1997),
573–591.

\item[{[AM2]}] R. Achilles and M. Manaresi, Mixed multiplicities, Segre numbers and Segre classes, J. Algebra 525 (2019), 390 - 415.

\item[{[AGM]}] F. Aroca, M. G\'omez-Morales and H. Mourtada, Groebner fan and embedded resolutions of ideals on toric varieties, Beitr. Algebra Geom. 65 (2024).

\item[{[BLQ]}] H. Blum, Y. Liu, and L. Qi, Convexity of multiplicities of filtrations on local rings, Compos. Math. 160 (2023),
878 - 914.

\item[{[Bou]}] N. Bourbaki, Commutative Algebra, Chapters 1-7, Springer Verlag, 1989.

\item[{[CPU]}] Y. Cid-Ruiz, C. Polini and  B. Ulrich, Multidegrees, families and integral dependence, arXiv:2405.07000.


\item[{[CJS]}] V. Cossart, U. Jannsen and S. Saito, Desingularization: Invariants and strategy,  Lecture Notes in Mathematics 2270, Springer 2020.

\item[{[CoP]}] V. Cossart and O. Piltant, Resolution of singularities of arithmetical threefolds, J. Algebra  529   (2019),    268 -535.


\item[{[C1]}] S.D. Cutkosky, Multiplicities of graded families of ideals on Noetherian local rings, arXiv:2603.06844.

\item[{[C9]}] S.D. Cutkosky, A counterexample to strong local monomialization in a tower of
two independent defect Artin-Schreier extensions, Annales de la
Facult\'e des Sciences Toulouse 33 (2024) 915 - 935.

\item[{[C10]}] S.D. Cutkosky, Erratic birational behavior of mappings in positive characteristic,
Mathematische Nachrichten 296 (2023), 5123 - 5156.

\item[{[C11]}] S.D. Cutkosky, Local Uniformization of Abhyankar Valuations, Michigan Math-
ematical Journal 71 (2022), 859 - 891.

\item[{[C12]}] S.D. Cutkosky, The Minkowski equality of filtrations, Advances in Math 338 (2021).

\item[{[C13]}] S.D. Cutkosky, Ramification of valuations and local rings in positive characteris-
tic, Communications in Algebra 44 (2016), 2828-2866.

\item[{[C14]}] S.D. Cutkosky, A general volume = multiplicity formula, Acta Math. Vietnamica 40 (2015), 139 - 147.



\item[{[C15]}] S.D. Cutkosky, Counterexamples to local monomialization in positive character-
istic, Math. Annalen. 362 (2015), 321 - 334.


\item[{[C16]}] S.D. Cutkosky, Asymptotic multiplicities of graded families of ideals and linear series,  Advances in Mathematics 264 (2014), 55 - 113.

\item[{[C17]}] S.D. Cutkosky, A Skeleton Key to Abhyankar’s Proof of Embedded Resolution
of Characteristic p Surfaces, Asian Journal of Math. 15 (2011),
369 - 416, Hironaka Volume.

 \item[{[C18]}] S.D. Cutkosky, Resolution of Singularities for 3-folds in positive characteristic,
American J. of Math. 131 (2009), 59-127.

\item[{[C19]}] S.D. Cutkosky, Toroidalization of dominant morphisms of 3-folds, Memoirs of the
AMS, Volume 190, American Mathematical Society, Providence
(2007), 222 pages.

\item[{[C20]}] S.D. Cutkosky, Local monomialization of transcendental extensions, Ann. Inst. Fourier 55 (2005), 1517 - 1586.


\item[{[C21]}] S.D. Cutkosky, Resolution of Singularities, Graduate Studies in Mathematics 63, American Mathematical Society, 2004.

\item[{[C22]}] S.D. Cutkosky, Local factorization and monomialization of morphisms, Ast\'erisque
260, 1999.






\item[{[CM1]}] S.D. Cutkosky and J. Monta\~no, Multiplicities and degree functions in local rings via intersection
products, J. London Math. Soc 112 (2025).



\item[{[CMT]}]  S.D. Cutkosky,  H. Mourtada and B. Teissier,    On the construction of valuations and generating sequences, 
 Algebraic Geometry, 8 (2021), 705 - 748.
 
 \item[{[CP]}] S.D. Cutkosky and O. Piltant, Ramification of valuations, Advances in Math. 183 (2004), 1-79.



\item[{[CSa2]}] S.D. Cutkosky and P. Sarkar, Analytic Spread of filtrations and symbolic algebras,
Journal London Math. Soc. 106 (2022), 2635 - 2662.

\item[{[CSS]}] S.D. Cutkosky, P. Sarkar and H. Srinivasan, Mixed multiplicities of filtrations, Transactions Amer. Math. Soc. 372 (2019),
6183 - 6211.

\thispagestyle{empty}

\item[{[sD]} S. Das, Epsilon multiplicity for graded algebras, J. Pure Appl. Algebra 225, 2021

\item[{[ELS]}] L. Ein, R. Lazarsfeld and K. Smith, Uniform Approximation of Abhyankar valuation ideals
in smooth function fields, Amer. J. Math. 125 (2003), 409 - 440.

%\item[{[JM]}] M. Jonsson and M. Musta\c{t}\u{a}, Valuations and asymptotic invariants for sequences of ideals, Ann. Inst. Fourier 62 (2012), 2145 - 2209.

\item[{[Fl]}] W. Fulton, Intersection Theory, second edition, Springer, 1998.

\item[{[GG]}] T. Gaffney and R. Gassler, Segre numbers and hypersurface singularities, J. Algebraic Geom. 8 (1999), 695–736.

\item[{[EGAIV]}] A. Grothendieck and J. Dieudonn\'e, El\'ements de G\'eom\'etrie Alg\'ebrique IV, Etude 
locale des sch\'emes et des morphismes de sch\'emes (seconde partie) Publ. Math.
IHES 24 (1965).

\item[{[H1]}] H. Hironaka, Resolution of singularities of an algebraic variety over a field of characteristic zero: I-II Ann. Math. 79, 109 - 326 (1964).

\item[{[H2]}] H. Hironaka, Desingularization of excellent surfaces, in Appendix to Resolution of Surface Singularities, V. Cossart, J. Giraud and U. Orbanz, Lecture Notes in Math. 1101, Springer, 1984.

\item[{[H3]}] H. Hironaka, Characteristic polyhedra of singularities, J. Math. Kyoto Univ. 7, 251 - 293 (1967).

%\item[{[dJ]}] A. J. de Jong, Smoothness, semi-stability and alterations, Publ. Math. I.H.E.S. 83 (1996), 51 - 93.



\item[{[Ka]}] D. Katz, Note on multiplicity, Proc. Amer. Math. Soc. 104 (1988), 1021 - 1026. 

\item[{[Kl]}] S. Kleiman, The Picard scheme, Fundamental Algebraic Geometry, 235 - 321, Math. Surveys Monogr. 123,
Amer. Math. Soc. 2005.

\item[{[KK]}] H. Knaf and F.-V. Kuhlmann, Abhyankar places admit local uniformization in any characteristic, Ann. Sci. Ec. Norm. Sup\'er. (4) 38 no. 6 (2005), 833–846. 

\item[{[Kou]}] A.G. Kouchirenko, Poly\`edres de Newton et nombres de Milnor, Invent. Math. 32 (1976), 1-31.


\item[{[Ku]}] F.-V. Kuhlmann, Elimination of ramification I: The generalized stability theorem, Trans.
AMS 362 (2010), 5697-5727.

\item[{[Laz1]}] R. Lazarsfeld, Positivity in Algebraic Geometry I, Springer 2000.

\item[{[LM]}] R. Lazarsfeld and M. Musta\c{t}\u{a}, Convex bodies associated to linear series, Ann. Sci. Ec. Norm.
Super 42 (2009) 783 - 835.



\item[{[Li3]}] J. Lipman,  Desingularization of 2-dimensional schemes, Annals of Math. 107 (1978), 115 -- 207.

\item[{[Li1]}] J. Lipman, Equimultiplicity, Reduction and Blowing Up.

\item[{[Li2]}] J. Lipman, Quasi-ordinary singularities of surfaces in $\CC^3$,  Proc. Sympos. Pure Math., 40, Amer.
Math. Soc., Providence, RI, 1983.

\item[{[McA]}] S. McAdam, Asymptotic prime divisors and analytic spread, Proc. AMS 90 (1980), 555-559.

\item[{[Mus]}] M. Musta\c{t}\u{a}, On multiplicities of graded sequences of ideals, J. Algebra 256 (2002), 229-249.

\item[{[Na]}] N. Nakayama, Zariski-decomposition and abundance, MSJ Memoirs 14, Math. Soc. Japan,
Tokyo, 2004.

\item[{[PTUV]}] C. Polini, N. V. Trung, B. Ulrich, and J. Validashti, Multiplicity sequence and integral dependence, Math. Ann. 378 (2020), 951–969.

\item[{[Rm]}] C.P. Ramanujam, A geometric interpretation of multiplicity, Inventiones Math. 22 (1973), 85 - 147.


\item[{[Re]}] D. Rees, A-transforms of local rings and a theorem on multiplicities of ideals, Proc. Cam-
bridge Philos. Soc. 57 (1961), 8 - 17.

\item[{[Re2]}] D. Rees, Degree functions in local rings, Proceedings of the Cambridge Phil. Soc. 57 (1961), 1-7.

\item[{[Re3]}] D. Rees, Lectures on the asymptotic theory of ideals, London Math Soc. Lecture Note Series 113, Cambridge
Univ. Press, 1988.

\item[{[RS]}] D. Rees and R. Sharp, On a Theorem of B. Teissier on Multiplicities of Ideals in Local Rings,
J. London Math. Soc. 18 (1978), 449-463.

\item[{[Sm]} ]I. Smirnov, Semicontinuity of mutliplicities in families, Documenta Mathematica 25 (2020), 381 - 399.
  
  \item[{[SH]}] I. Swanson and C. Huneke, Integral Closure of Ideals, Rings and Modules, Cambridge University Press, 2006.
  
 \item[{[T1]}] B. Teissier,   Cycles \'evanescents, sections planes et conditions de Whitney, Ast\'erisque 7-8 (1973) 285 - 362.                                                                    

\item[{[T2]}] B. Teissier, Cycles \'evanescents et resolution simultan\'ee, I et II, S\'eminaire sur les singularities 1976 - 1977, Springer Lecture Notes in Mathematics 777, 1980, 71 - 146.

\thispagestyle{empty}


\item[{[T3]}] B. Teissier, Sur une in\'egalit\'e \`a la Minkowski pour les multiplicit\'es (Appendix to a paper by
D. Eisenbud and H. Levine), Ann. Math. 106 (1977), 38 - 44.

\item[{[T4]}] B. Teissier, On a Minkowski type inequality for multiplicities II, In C.P. Ramanujam - a
tribute, Tata Inst. Fund. Res. Studies in Math. 8, Berlin - New York, Springer, 1978.

\item[{[Te]}] M. Temkin, Desingularization of quasi-excellent schemes schemes of characteristic zero, Adv. Math. 219 (2008), 488 - 522.

\item[{[Te1]}] M. Temkin, Tame distillation and desingularization by p-alterations, Ann. of Math.
186 (2017), 97 - 126.

\item[{[Th]}] A. Thorup, Rational equivalence theory on arbitrary Noetherian schemes, In Enumerative Geometry, Sebastian
Xambó-Descamps Ed, LNM 1436, Springer, 1990.


\item[{[UV]}] B. Ulrich and J. Validashti, Numerical criteria for integral dependence, Math. Proc. Cambridge Philos. Soc. 151
(2011), 95–102.

\item[{[V]}] A.N. Varchenko, Zeta-function of monodromy and Newton's diagram, Invent, Math. 37 (1976), 253 - 262.
\end{enumerate}

\end{document}








\centerline{\bf Results From Prior NSF Support}
\vskip .2truein
NSF award number: DMS-2054394, Amount: \$316,364

 Period of support: 07/01/2021 - 06/30/2024.

Title of project: Interaction of  Commutative Algebra, Valuations and Geometry.

\vskip .2truein
%Multiplicities and Valuations in Commutative Algebra

%\noindent{\bf Intellectual Merit.}

\section{Filtrations, divisors, analytic spread, epsilon multiplicity, mixed multiplicity and ramification}
\thispagestyle{empty}

\subsection{Analytic spread} \thispagestyle{empty}
Basics on Rees algebras $R[\mathcal I]$ and the analytic spread $\ell(\mathcal I)$ of a filtration $\mathcal I$ are given in Section \ref{SecAnS} of the Project Description.
We have the following remarkable theorem.




\begin{Theorem}\label{TheoremC4}([McA], Theorem 5.4.6 [SH]) Let $R$ be a formally equidimensional local ring and $I$ be an ideal in $R$. Then $m_R\in \mbox{Ass}(R/\overline{I^n})$ for some $n$ if and only if $\ell(I)=\dim(R)$.
\end{Theorem}

The assumption of being formally equidimensional is not required for the if direction of Theorem \ref{TheoremC4} (this is Burch's theorem, [Bu], Proposition 5.4.7 [SH]).

 
To see how Theorem \ref{TheoremC4} can be extended, we  consider divisorial filtrations, which are defined in Subsection \ref{SecAnS} of the Project Description. 

We prove that Theorem \ref{TheoremC4} is true for $\QQ$-divisorial filtrations in an excellent local domain of equicharacteristic zero or of dimension $\le 3$. Resolution of singularities is used in the proof which is the reason for the restriction of equicharacteristic zero or of dimension $\le 3$.

\begin{Theorem}(Theorem 1.4 [C10])\label{Theorem2} Let $R$ be an excellent local domain of equicharacteristic 0, or  of dimension $\le 3$. Let $\mathcal I=\{I_n\}$ be a $\QQ$-divisorial filtration on $R$. Then the following are equivalent.
\begin{enumerate}
\item[1)] The analytic spread of $\mathcal I$ is $\ell(\mathcal I)=\dim R$.
\item[2)] There exists $n_0\in \ZZ_{>0}$ such that $m_R\in \mbox{Ass}(R/I_n)$ if $n\ge n_0$.
\item[3)] $m_R\in \mbox{Ass}(R/I_{m_0})$ for some $m_0\in \ZZ_{>0}$.
\end{enumerate}
\end{Theorem}

1) implies 2) is true for Noetherian local domains (Theorem 4.6 [CSa2]). 3) implies 1) is Theorem 1.3 of [C10]. Example 1.5 of [C10] shows that 3) implies 1) is false for $\RR$-divisorial valuations. 

 
\subsection{Epsilon multiplicity and analytic spread}\label{SubsecEp}\thispagestyle{empty} The epsilon multiplicity of an ideal $I$ in a $d$-dimensional local ring $R$ with maximal ideal $m_R$ is defined by Ulrich and Valadashti in [UV] to be 
\begin{equation}\label{eq12}
\epsilon(I)=d!\limsup_n\frac{\lambda_R(H^0_{m_R}(R/I^n))}{n^d}.
\end{equation}
It is shown there, by comparison with the $j$-multiplicity,  that $\epsilon(I)$ is always a real number. It is shown in [CHST] that $\epsilon(I)$ can be an irrational number. 
In a local ring $R$, we have that 
$$
H^0_{m_R}(R/I)=I:m_R^{\infty}/I.
$$



For ideals in analytically unramified local rings, the epsilon multiplicity is a limit (Corollary 6.3 [C13]). 

%\begin{Theorem}( \cite[Corollary 6.3]{C1})\label{TheoremI1} Suppose that $I$ is an ideal in an analytically unramified local ring $R$. Then $\epsilon(I)$ is actually a limit, 
%$$
%\epsilon(I)
%=d!\lim_{n\rightarrow \infty}\frac{\lambda_R(I^n:m_R^{\infty}/I^n)}{n^d}.
%$$
%\end{Theorem}



%In [CSa3], we extend epsilon multiplicity to filtrations and obtain some results generalizing theorems about epsilon multiplicities for ideals. 

In [CSa3] we extend the definition of epsilon multiplicity to  filtrations by defining the epsilon multiplicity of a filtration $\mathcal I=\{I_n\}$ of ideals in a $d$ dimensional local ring $R$ to be
\begin{equation}\label{epdef}
\epsilon(\mathcal I)=d!\limsup_n\frac{\lambda_R(H^0_{m_R}(R/I_n))}{n^d}.
\end{equation}

We prove in [CSa3] that the epsilon multiplicity is a limit for many naturally occuring filtrations, including $\RR$-divisorial filtrations in  analytically unramified local rings. 


The following theorem about epsilon multiplicity of ideals follows from results of Ulrich and Validashti.

\begin{Theorem}\label{TheoremI2} Suppose that $R$ is a universally catenary Noetherian local ring of dimension $d$  and $I$ is an ideal of $R$. Then the analytic spread
$\ell(I)=d$ if and only if $\epsilon(I)>0$.
\end{Theorem}

%The following corollary is immediate  from Theorem \ref{Theorem2} and \cite{MA} or Theorem 5.4.6 \cite{HS}.


%\begin{Corollary}\label{CorI3} Suppose that $R$ is a universally catenary Noetherian local ring of dimension $d$ and $I$ is an ideal of $R$. Let $\nu_1,\ldots,,\nu_r$ be the Rees valuations of $R$. Then $\epsilon(I)>0$ if and only if the center of at least one of the $\nu_i$ on $R$ is the maximal ideal.
%\end{Corollary}




  
  

 We obtain the following generalization of Theorem \ref{TheoremI2} to filtrations.  Divisorial filtrations and the analytic spread $\ell(\mathcal I)$ of a filtration are defined in Subsection \ref{SecAnS}. Theorem \ref{TheoremI4} gives a further equivalent condition to Theorem \ref{Theorem2}.
 

\begin{Theorem}\label{TheoremI4}(Theorem 1.5 [CSa3]) Let $R$ be a $d$-dimensional excellent normal local domain of equicharacteristic zero, or an arbitrary excellent local domain of dimension $\le 3$.  Let $\mathcal I$ be a 
$\QQ$-divisorial filtration of $R$. Then 
 $\epsilon(\mathcal I)>0$ if and only if 
the analytic spread $\ell(\mathcal I)=d$.
\end{Theorem}

 Theorem \ref{TheoremI4} is not true for general filtrations. Example 4.2 of [CSa3]
   gives an example of an $\RR$-divisorial filtration $\mathcal I$ in a regular local ring such that $\epsilon(\mathcal I)>0$ but $\ell(\mathcal I)<\dim R$.
   
   

\subsection{The Minkowski equality} \thispagestyle{empty}
Given an  $m_R$-filtration $\mathcal I=\{I_n\}$ in a $d$-dimensional analytically unramified local ring $R$, the multiplicity
\begin{equation}\label{I5}
e_R(\mathcal I)=\lim_{n\rightarrow\infty}d!\frac{\lambda_R(R/I_n)}{n^d};
\end{equation}
this limit is proven to exist in Theorem 1.1 [C13] and Theorem 4.2 [C14]. The problem of existence of such limits (\ref{I5}) has been considered by Ein, Lazarsfeld and Smith [ELS] and Musta\c{t}\u{a} [Mus].
When the ring $R$ is a domain and is essentially of finite type over an algebraically closed field $k$ with $R/m_R=k$, Lazarsfeld and Musta\c{t}\u{a} [LM] showed that
the limit exists for all $m_R$-filtrations.




Mixed multiplicities of filtrations are defined and proven to exist as limits  in analytically unramified local rings in [CSS]. 




%%%%%%%%%%%%%%%%%%%%%%%%%%%%%
%The Minkowski inequalities hold for $m_R$-filtrations. 


The Minkowski inequalities were formulated and proven for   $m_R$-primary ideals in reduced equicharacteristic zero local rings by Teissier [T1], [T2] and proven for  $m_R$-primary ideals in full generality, for  local rings,  by Rees and Sharp [RS]. The same inequalities hold for filtrations. This is proven in Theorem 6.3 [CSS]. A consequence of this result is that for $m_R$-filtrations $\mathcal I(1)$ and  $\mathcal I(2)$ in an analytically unramified $d$-dimensional local ring $R$, we have that 
\begin{equation}\label{eqMinkIn}
e_R(\mathcal I(1)\mathcal I(2))^{\frac{1}{d}}\le e_R(\mathcal I(1))^{\frac{1}{d}}+e_R(\mathcal I(2))^{\frac{1}{d}}
\end{equation}
where $\mathcal I(1)\mathcal I(2)=\{I(1)_mI(2)_m\}$.
This equation is called The Minkowski Inequality.


There is a beautiful characterization of when equality holds in the Minkowski inequality (\ref{eqMinkIn}) for $m_R$-primary ideals $I$ (the $I$-adic filtrations) by 
Teissier [T3], Rees and Sharp [RS] and Katz [Ka].


  
 
 In Theorems 13.2 and 14.5 of [C1],     
 we give a complete generalization of the Teissier, Rees and Sharp, Katz Theorem for bounded $m_R$-filtrations. Bounded  $m_R$-filtrations  are a class of filtrations which include the $I$-adic filtrations and the divisorial filtrations.  
 
 
\begin{Theorem}\label{NewTRSKT}(Theorem 13.2 and Theorem 14.5 [C1]) Suppose that $R$ is a $d$-dimensional analytically irreducible or excellent local domain and $\mathcal I(1)$ and $\mathcal I(2)$ are bounded $m_R$-filtrations. Then the following are equivalent
\begin{enumerate}
\item[1)]  The Minkowski equality
$$
e_R(\mathcal I(1)\mathcal I(2))^{\frac{1}{d}}=e_R(\mathcal I(1))^{\frac{1}{d}}+e_R(\mathcal I(2))^{\frac{1}{d}}
$$
holds.
\item[2)] There exist positive integers $a,b$ such that there is equality of  integral closures
$$
\overline{\sum_{n \ge 0}I(1)_{an}t^n}=\overline{\sum_{n\ge 0}I(2)_{bn}t^n}
$$
 in $R[t]$.
\end{enumerate}
 \end{Theorem}
 
It is proven in  Theorem 8.4 [C1] that 2) implies 1) is true for $m_R$-filtrations $\mathcal I(1)$, $\mathcal I(2)$ in a $d$-dimensional analytically unramified local ring. An example is given in Remark 6.4 [CSS], showing that the  converse statement, that  the Minkowski equality 1) implies condition 2)   is not true for filtrations, even in a regular local ring.
   


\subsection{Ramification, K\"ahler differentials and deeply ramified fields}\thispagestyle{empty}
Background information on valuations and extensions of valuations is given in Subsection \ref{SecRam} of the Project Description.
In our papers [CKR] and [CK], we make  explicit constructions of  extensions 
of valuation rings $\cO_L |\cO_K$ where $(L|K,v)$ is a Kummer
extension of valued fields of prime degree  or an Artin-Schreier extension, and give a characterization 
in terms of invariants of the valuation $v$ of when we have the vanishing 
$\Omega_{\cO_L|\cO_K}=0$. 

Defect extensions are discussed in Subsection \ref{SecRam}. In the case of defect extensions, the characterization is as follows.


\begin{theorem} (Theorem 1.2 [CKR])                          \label{Ocharind}
Take a valued field $(K,v)$ with $\chara Kv>0$; if $\chara K=0$, then assume that $K$
contains all $p$-th roots of unity. Further, take a Galois extension $(L|K,v)$ of 
prime degree with nontrivial defect. Then the
extension has independent defect if and only if 
\begin{equation}                         %\label{}
\Omega_{\cO_L|\cO_K} \>=\> 0\>.
\end{equation}
\end{theorem}

\pars


Our classification allows us  to give an
explicit construction of the extension of valuation rings $\cO_L|\cO_K$ in a finite 
Galois extension $L|K$ of valued fields, and to give an explicit computation of  
ramification and the K\"ahler differential $\Omega_{\cO_L|\cO_K}$ of this extension.
\pars

\parm
We define a nontrivially valued field $(K,v)$ to be a \bfind{deeply ramified field} 
if the following conditions hold:
\sn
{\bf (DRvg)} whenever $\Gamma_1\subsetneq\Gamma_2$ are convex subgroups of the value
group $vK$, then $\Gamma_2/\Gamma_1$ is not isomorphic to $\Z$ (that is, no
archimedean component of $vK$ is discrete);
\sn
{\bf (DRvr)} if $\chara Kv=p>0$, then the homomorphism
\begin{equation}                          \label{homOpO}
\cO_{\hat K}/p\cO_{\hat K} \ni x\mapsto x^p\in \cO_{\hat K}/p\cO_{\hat K}
\end{equation}
is surjective, where $\cO_{\hat K}$ denotes the valuation ring of the completion 
$\hat K$ of $(K,v)$. 
\pars
Take a valued field $(K,v)$ with valuation ring $\cO_K\,$. Choose any extension of $v$ 
to $K\sep$ and denote the valuation ring of $K\sep$ with respect to this extension by 
$\cO_{K\sep}\,$. In Theorem 6.6.12 of [GR], Gabber and Ramero prove the following result 
%
\begin{theorem}              \label{GRthm}
A valued field $(K,v)$ satisfies conditions (DRvg) and (DRvr) if and only if
\begin{equation}                         \label{GRdefdr}
\Omega_{\cO_{K\sep}|\cO_K} \>=\> 0\>.
\end{equation}
This does not depend on the chosen extension of the 
valuation from $K$ to $K\sep$.
\end{theorem}
\pars
In [CK] we give an alternative proof of Theorem~\ref{GRthm}. Our proof relies on our construction of relative K\"ahler differentials of
extensions of valuation rings in finite separable extensions. Our methods are valid for valuations of arbitrary rank. The proof of Gabber and Romero uses an induction argument to reduce to the case of rank 1 valuations. 

\pars
%Our proof of Theorem~\ref{GRthm} yields the following new result:
%
%\begin{theorem}(Theorem 1.4 [CK])                              \label{GRthm+}
%Let $(K,v)$ be a valued field of residue characteristic $p>0$. If $K$ has characteristic
%$0$, then assume in addition that it contains all $p$-th roots of unity. Then the 
%following assertions are equivalent:
%\sn
%a) $\,\Omega_{\cO_L|\cO_K}=0$ for all Galois extensions $(L|K,v)$ of degree $p$,  
%\sn
%b) $\,\Omega_{\cO_{K\sep}|\cO_K}=0$,
%\sn
%c) $(K,v)$ is a deeply ramified field.
%\end{theorem}

\noindent{\bf Broader Impacts.}\thispagestyle{empty}

During the time of this award, three students at the University of Missouri earned their PhDs under my direction. 
Roberto N\'u\~nez graduated in 2021. His thesis is ``Volumes of Line Bundles on Schemes''. He is currently a PostDoc in the Department of Mechanical and Aerospace Engineering at the  University of Missouri. Navaneeth Chenicheri Chathoth graduated
in 2022. His thesis is ``Simultaneous Local Resolution in Positive Characteristic''. Smita Praharaj graduated in 2023. Her thesis is 
``The Asymptotic Samuel Function of a Filtration''. Currently, Stephen Landsittel and Marc Abi Aad are working with me towards a PhD. 

I mentored young mathematicians during this award. 
Arpan Dutta (PhD under my direction in 2018) is now an Assistant Professor at
IIT Bhubaneswar. The Indian Institutes of Technology (IIT) are the premier engineering and science universities in India. 
Suprajo Das (PhD under my direction in 2020) is currently a PostDoc at IIT Bombay. Parangama Sarkar (PostDoc at Univ. Missouri in 2019) is now an Assistant Professor at IIT Palakkad. During the period of this NSF award, we wrote two papers on the analytic spread of a filtration, [CSa2] and [CSa3]. Former PhD student Razieh Ahmadian was promoted to Senior PostDoctoral Researcher at IPM (Institute for Research in Fundamental Sciences) in Tehran, Iran. She earned her PhD from IPM in 2015. I was her joint advisor with Rahim Zaare Nahandi. In this last year, I introduced her to a new area in mathematics from her previous research. We wrote the paper [AC]. This paper gives a major extension of  a theorem of myself with Mourtada and Teissier from our paper [CMT]. In [AC] we prove in complete generality the existence of a generating sequence for a valuation in an arbitrary Noetherian domain $A$ dominated by
a valuation with some mild conditions on the valuation. Most of the above mentoring (including the theses of Chathoth and Praharaj) was done on Zoom, which was  extremely difficult psychologically for young people, and introduced new challenges for all of us. 


\vskip .1truein
\noindent{\bf Conferences held during the time of the award that I coorganized:} \thispagestyle{empty}
\begin{enumerate}
\item MSRI Graduate Summer School on Commutative Algebra and its interaction with Algebraic Geometry,  University of Notre Dame, May - June 2023. Coorganized with Claudia Poloni, Claudiu Raicu, Steven Sam and Kevin Tucker. 
\item ``Resolutions in local algebra and singularity theory'', Oberwolfach Workshop,  Feb. 5 - Feb. 11, 2023, Oberwolfach Germany (Coorganized with Aldo Conca, Eleonore Faber and Srikanth Iyengar).
\item Conference on ``Singularities in Positive Characteristic'', CIRM, Luminy Marseille,  July 12 - 16, 2021 (Scientific Committee).
\end{enumerate}

\noindent Coorganizer, from October 2022,  of The Fellowship of the Ring, a virtual seminar in commutative algebra, with Christine Berkesch, David Eisenbud, Mel Hochster, Srikanth Iyengar, Claudia Miller, Irena Peeva, Karl Schwede and Bernd Ulrich.




%\noindent{\bf Current Ph.D. Students}
%\begin{enumerate}





%\item Stephen Landsittel

%\item Marc Abi Aad

%\end{enumerate}
%\vskip .2truein
%\noindent{\bf  Ph.D. Students}
%\vskip .2truein
%\begin{enumerate}
%\setcounter{enumi}{2}
%\item Smita Praharaj. Thesis: ``The Asymptotic Samuel Function of a Filtration'', May 2023.
%\item Navaneeth Chenicheri Chathoth. Thesis: ``Simultaneous Local Resolution in Positive Characteristic'', July, 2022.
%\item Roberto N\'u\~nez. Thesis: ``Volumes of Line Bundles on Schemes'', July, 2021.
%\end{enumerate}
%\begin{enumerate}
%\item Smita Praharaj. Thesis: ``The Asymptotic Samuel Function of a Filtration'', May 2023.
%\item Navaneeth Chenicheri Chathoth. Thesis: ``Simultaneous Local Resolution in Positive Characteristic'', July, 2022.
%\item Roberto N\'u\~nez. Thesis: ``Volumes of Line Bundles on Schemes'', July, 2021.
%\item Suprajo Das. Thesis: ``Epsilon Multiplicity for Graded Algebras'', May, 2020.
%\item Arpan Dutta. Thesis: ``Generating Sequences and Semigroups of Valuations'', May 2018.
%\item Razieh Ahmadian. Thesis: ``Toroidal Morphisms of 3-Folds'', December, 2015.
%\end{enumerate}


%%%%%%%%%%%%%%%%
%\vskip 2truein


%Arpan Dutta earned his Ph.D. under my direction in May 2018, with his thesis ``Generating Sequences and Semigroups of Valuations''. He is currently a PostDoc at IISER, Mohali (India).     Suprajo Das earned his Ph.D. under my direction in May 2019, with his thesis ``Epsilon Multiplicity for Graded Algebras ''. He is currently a PostDoc at the Chennai Mathematical Institute in India.  Roberto N\'u\~nez is expected to graduate in May, 2021. 
% Navaneeth Chenicheri Chathoth and Smita Prahraj
%are currently working with me towards a Ph.D.   

%I mentored young mathematicians during this grant.  When Parangama Sarkar was a PostDoc at the University of Missouri on a Government of India INSPIRE Fellowship we wrote our joint paper [CSS] with Hema Srinivasan. Josnei Novacoski visited the University of Missouri where we wrote our joint paper [CNo]. When Hussein Mourtada visited the University of Missouri we wrote our paper [CM]. Later, we wrote the paper [CMT] with Bernard Teissier. 
%Quang Hoa Tran also visited the University of Missouri. We wrote with Marc Chardin the joint paper [ChCH].


%\vskip .1truein
%\noindent{\bf Conferences organized:} 
%\begin{enumerate}

%\item``Singularities and Homological Aspects of Commutative Algebra'', Oberwolfach workshop,  Oberwolfach, Germany, Feb. 10 - Feb. 16, 2019,  Organized with 
%Aldo Conca and Srikanth Iyengar.
 

%\item``Holomorphic Foliations and Singularities of Mappings and Spaces'',   Special Session of the Mathematical Congress of the Americas, Montreal, July 23 - 28, 2017.
%Organized with Marcio Soares, Fernando Cukierman, Jawad Snoussi, Lev Birbrair and Marcelo Jose Saia.

%\item``The Prospects for Commutative Algebra'', Osaka, Japan,  July 10 - 14, 2017,   Organized with Takayuki Hibi, David Eisenbud and Juergen Herzog. 

%\item``Second International Meeting in Commutative Algebra and its Related Areas'' (in honor of the 75th birthday of Aron Simis),  University of Sao Paulo, Brazil, July 7 - July 21, 2017, Member of Scientific Committee. 
%\end{enumerate}



%I have been an editor of the Journal of Algebra since 2005. 
%I was an editor, with David Eisenbud, Shiro Goto, J\"urgen Herzog, Takayuki Hibi and  Ng\^o Vi\^et Trung, of a special volume of  Acta Mathematica Vietnamica, Volume 44, 2019, dedicated to recent developments in commutative algebra and related fields.
%My book ``Introduction to Algebraic Geometry'' [C9]  was published by the  American Mathematical Society in 2018.
%The book is  based on my lectures at the University of Missouri.



\vskip .2truein


\centerline{\bf Publications resulting from the NSF award}\thispagestyle{empty}
\vskip .1truein

\begin{enumerate}\thispagestyle{empty}
\item[{[AC]}] Generating sequences of valuations on simple extensions of domains, with Razieh Ahmadian, arXiv:1904.10702.

\item[{[CK]}] K\"ahler differentials of extensions of valuation rings and deeply ramified
fields, with Franz-Viktor Kuhlmann, arXiv:2306.04967.

\item[{[CSa3]}] Epsilon multiplicity and analytic spread of filtrations, with Parangama Sarkar, arXiv2305.17532. 

\item[{[CKR]}] Characterizations of Galois extensions with independent defect, with Franz-Viktor Kuhlmann and Anna Rzepka, arXiv:2305.10022.


\item[{[CPr]}] The Asymptotic Samuel Function of a Filtration, with Smita Praharaj,

\noindent arXiv:2211.13205.

\item[{[C11]}] A counterexample to strong local monomialization in a tower of two independent defect Artin-Schreier extensions, to appear in Annales de la Facult\'e des Sciences Toulouse, arXiv:2302.07710.


\item[{[C10]}] Analytic spread and associated primes of divisorial filtrations, arXiv:2204.03047, to appear in Proceedings of the AMS.
\item[{[C9]}] Erratic birational behavior of mappings in positive characteristic, to appear in Mathematische Nachrichten, arXiv:2111.12818

\item[{[C8]}] The Minkowski equality of big divisors, to appear in Michigan Math. Journal, arXiv:2017.05091.


\item[{[C7]}] Errata to ``A simpler proof of toroidalization from 3-folds to surfaces'', to appear in Annales l'Institut Fourier.


\item[{[C6]}] Analytic spread of filtrations on two dimensional normal local rings,  Nagoya Math. Journal 249 (2023), 239 - 268.

\item[{[CSa2]}] Analytic Spread of Filtrations and Symbolic Algebras, with Parangama Sarkar, Journal of the London Math. Soc. (2) 106 (2022), 2635 - 2662.


\item[{[CSa1]}] Multiplicities and mixed multiplicities of arbitrary filtrations, with Parangama Sarker, 
 Research in the Mathematical Sciences 9, Volume in honor of J\"urgen Herzog, (2022).

\item[{[C5]}] Abhyankar's fundamental work on resolution of singularities, to appear in The Mathematics Student.

\item[{[C4]}] Essential Finite Generation of Valuation Rings in Characteristic Zero Algebraic Function Fields,  Algebra and Number Theory 16 (2022), 290 - 310.

\item[{[C3]}] Local Uniformization of Abhyankar Valuations,  Michigan Mathematical Journal 71 (2022), 859 - 891.


\item[{[CSr]}] Multiplicites and Mixed Multiplicites of Filtrations, with Hema Srinivasan, Commutative Algebra, dedicated to David Eisenbud on the occasion of his 75th birthday, Irena Peeva editor, Springer, Cham, 2021, 299 - 316.

\item[{[CMT]}] On the construction of valuations and generating sequences, with Hussein Mourtada and Bernard Teissier,  Algebraic Geometry, 8 (2021), 705 - 748.

\item[{[CN]}] Volumes of line bundles on schemes, with Roberto N\'u\~nez,  Proc. AMS. 149 (2021), 4099 - 4108.

\item[{[C2]}] Examples of multiplicities  and mixed multiplicities of filtrations, in Commutative Algebra - 150 years with Roger and Sylvia Wiegand, Contemporary Math. 773 (2021), 19 - 34.



\item[{[C1]}] The Minkowski Equality of Filtrations, Advances in Mathematics 388 (2021).

\item[{[CN]}] Essentially finite generation of valuation rings in terms of classical invariants, with Josnei Novacoski,   Mathematische Nachrichten 294 (2021), 15 - 37.


\item[{[CCT]}] Fibers of rational maps and Jacobian matrices, with Marc Chardin and Quang Hoa Tran, Journal of Algebra 571 (2021), 40 - 54.

\end{enumerate}
\thispagestyle{empty}
\vfill\eject


\vskip .2truein
\centerline{\bf Project Description}
%\noindent{\bf Intellectual Merit.}



\section{Analytic spread of filtrations, Hilbert functions, generalized multiplicities, inseparable local uniformization and ramification}
\thispagestyle{empty}
\subsection{Analytic spread of filtrations}\label{SecAnS}\thispagestyle{empty}
 Let $R$ be a Noetherian local ring with maximal ideal $m_R$. If $M$ is an $R$-module, we denote the length of $M$ by $\lambda_R(M)$ or $\lambda(M)$.
 Let $\mathcal I=\{I_n\}$ be a filtration on  $R$. A filtration $\mathcal I$ is called an $m_R$-filtration if $I_n$ is $m_R$-primary for $n\ge 1$.
 
 The Rees algebra of the filtration $\mathcal I$ is $R[\mathcal I]=\oplus_{n\ge 0}I_n$,
 the fiber cone of $\mathcal I$ is $F[\mathcal I]=R[\mathcal I]/m_RR[\mathcal I]$ and the analytic spread of the filtration $\mathcal  I$ is defined to be
\begin{equation}\label{eqN10}
\ell(\mathcal I)=\dim F[\mathcal I]=\dim R[\mathcal I]/m_RR[\mathcal I].
\end{equation}
Thus if $I$ is an ideal in $R$ and $\mathcal I=\{I^n\}$ is the $I$-adic filtration, then the classical analytic spread $\ell(I)$ of $I$ of $R$ is  $\ell(I)=\ell(\mathcal I)$.
The essential new phenomena in the case of arbitrary filtrations $\mathcal I$ is that the Rees algebra $R[\mathcal I]$ may not be Noetherian.  

An essential property of the analytic spread $\ell(I)$ of an ideal is the inequality (Proposition 5.1.6 and Corollary 8.3.9 [SH]) 
\begin{equation}\label{eqI1}
{\rm ht}(I)\le \ell(I)\le \dim R.
\end{equation}

If $\mathcal I$ is a filtration, then 
${\rm ht}(I_n)={\rm ht}(I_1)$ for all $n$ (Equation (7) [CSa2]) so it is natural to define ${\rm ht}(\mathcal I)={\rm ht}(I_1)$.


We always have (Lemma 3.6 [CSa2]) that 
\begin{equation}\label{eqZ1}
\ell(\mathcal I)\le \dim R
\end{equation}
so the second inequality of (\ref{eqI1}) always holds for filtrations. 
 However, the first inequality of (\ref{eqI1}),
$\mbox{ht}(\mathcal I)\le \ell(\mathcal I)$,
fails spectacularly, even attaining the condition that $\ell(\mathcal I)=0$ (Examples 1.2,  6.1 and 6.6 of [CSa2]). The last two of these examples are ``naturally occuring'' filtrations of symbolic powers of space curves, which are  divisorial filtrations (defined below). These last two examples are difficult to construct. 

 %The condition that  a filtration has  analytic spread zero has a simple ideal theoretic interpretation (\cite[Lemma 3.8]{CPS}).
 %Suppose that  $\mathcal I=\{I_n\}$ is a filtration in a local ring $R$. Then the analytic spread
%$\ell(\mathcal I)=0$ if and only if
%$$
%\mbox{For all $n>0$ and $f\in I_n$, there exists $m>0$ such that $f^m\in m_RI_{mn}$.}
%$$

%Let $\mathcal I=\{I_n\}$ be a filtration of a Noetherian local ring $R$.
%The integral closure  of the Rees algebra $R[\mathcal I]=\sum_{n\ge 0}I_nt^n$ in $R[t]$ is  denoted by $\overline{R[\mathcal I]}$.


Let $R$ be a  Noetherian local domain of dimension $d$ with quotient field $K$.  Let $\nu$ be a discrete valuation of $K$ with valuation ring $\mathcal O_{\nu}$ and maximal ideal $m_{\nu}$.  Suppose that $R\subset \mathcal O_{\nu}$. Then for $n\in \NN$, define valuation ideals
$$
I(\nu)_n=\{f\in R\mid \nu(f)\ge n\}=m_{\nu}^n\cap R.
$$
%We can extend this definition to $\alpha\in \RR$, defining 
%$$
%I(\nu)_{\alpha}=\{f\in R\mid \nu(f)\ge \alpha\}=I(\nu)_{\lceil\alpha\rceil}.
%$$
 
  A divisorial valuation of $R$ (Definition 9.3.1 [SH]) is a valuation $\nu$ of $K$ such that if $\mathcal O_{\nu}$ is the valuation ring of $\nu$ with maximal ideal $m_{\nu}$, then $R\subset \mathcal O_{\nu}$ and if $p=m_{\nu}\cap R$ then $\mbox{trdeg}_{\kappa(p)}K\nu={\rm ht}(p)-1$, where $\kappa(p)$ is the residue field of $R_{p}$ and $K\nu$ is the residue field of $\mathcal O_{\nu}$.   If $\nu$ is divisorial valuation of $R$ such that $m_R= m_{\nu}\cap R$, then $\nu$ is called an $m_R$-valuation.
   
 
 
 Let $I$ be an ideal in a Noetherian local domain $R$.
There exist divisorial valuations $\nu_1,\ldots,\nu_r$ of $R$ called the Rees valuations of $I$ and $a_1\ldots,a_r\in \ZZ_{>0}$ such that 
$$
\overline{I^n}=I(\nu_1)_{na_1}\cap\cdots\cap I(\nu_r)_{na_r}
$$
for $n\in \NN$. This expression is irredundant, in the sense that for all $j$, there exists an $n$ such that  $\overline{I^n}\ne \cap_{n\ne j}I(\nu_i)_{na_i}$. This is proven (even for a generalization of this statement to arbitrary Noetherian rings) in [Re] and Theorem 10.2.2 [SH]. This filtration is an example of a divisorial filtration.
 
 An integral  divisorial filtration of $R$ (which we will also refer to as a divisorial filtration) is a filtration $\mathcal I=\{I_m\}$ such that  there exist divisorial valuations $\nu_1,\ldots,\nu_s$ and $a_1,\ldots,a_s\in \ZZ_{\ge 0}$ such that for all $m\in \NN$,
$$
I_m=I(\nu_1)_{ma_1}\cap\cdots\cap I(\nu_s)_{ma_s}.
$$

Given an ideal $I$ in $R$, the filtration $\{\overline{I^n}\}$ is an example of a divisorial filtration of $R$. The filtration 
$\{\overline{I^n}\}$ is Noetherian if $R$ is universally Nagata. The symbolic powers of an ideal in a regular local ring is another example of a divisorial filtration. This filtration is often not Noetherian. 


The filtration $\mathcal I$ is called an $\RR$-divisorial filtration if $a_1,\ldots,a_s\in \RR_{>0}$ and $\mathcal I$ is called a $\QQ$-divisorial filtration if $a_1,\ldots,a_s\in \QQ$. If $a_i\in \RR_{>0}$, then
$$
I(\nu_i)_{na_i}:=\{f\in R\mid \nu_i(f)\ge na_i\}=I(\nu_i)_{\lceil na_i\rceil},
$$
where $\lceil x\rceil$ is the round up of a real number $x$. 

If $R$ is an excellent domain and $\mathcal I=\{I_n\}$ is a divisorial filtration of $R$, then there exists a projective normal $\mbox{Spec}(R)$-scheme $X$ which birationally dominates $\mbox{Spec}(R)$ such that the center of each $\nu_i$ on $X$ is a prime divisor $E_i$ on $X$. Thus for all $n$,
$$
I_n=\Gamma(X,\mathcal O_X(-\lceil nD\rceil))\cap R
$$
 where $D=\sum_{i=1}^sa_iE_i$ is an $\RR$-Weil divisor on $X$. If resolution of singularities holds above $\mbox{Spec}(R)$, we can take $X$ to be nonsingular so that  the $E_i$ are Cartier divisors. 
We say that $D$ respresents the filtration $\mathcal I$ and write $I_n=I(nD)$ and $\mathcal I=\mathcal I(D)$.




\begin{Question}  Is it possible to extend Theorem \ref{Theorem2}
 to more general rings, for instance excellent local domains?
\end{Question}

The implication 1) implies 2) of Theorem \ref{Theorem2}  is proven for arbitrary Noetherian local domains in [CSa2]. To extend the proof of 3) implies 1)  it is necessary to extend results on sectional decomposition for Cartier divisors on nonsingular schemes in [Nak], [C10] to sectional decomposition of Weil divisors on a normal projective scheme birationally dominating $\mbox{Spec}(R)$. There is a theory of sectional decomposition of Weil divisors, but the results are much less strong than for Cartier divisors on nonsingular schemes.

%{Theorem}\label{ThmI1}(Theorem \ref{Cor2}) Let $R$ be a two dimensional normal excellent local ring. The following are equivalent for a $\QQ$-divisorial filtration $\mathcal I(D)$ on $R$.
%\begin{enumerate}
%\item[1)]
%The analytic spread $\ell(\mathcal I(D))=\dim R[D]/m_RR[D]=2$.
%\item[2)] $m_R\in \mbox{Ass}(R/I(nD))$ for some $n$.
%\item[3)] There exists $n_0\in \ZZ_{>0}$ such that $m_R\in \mbox{Ass}(R/I(nD))$ for all $n\ge n_0$.
%\end{enumerate}
%\end{Theorem}

\begin{Question}\label{Q20} Suppose that $\mathcal I(D)$ is a $\QQ$-divisorial filtration. Does  the fiber cone of $\mathcal I(D)$,
$$
F[\mathcal I(D)]=R[\mathcal I(D)]/m_RR[\mathcal I(D)]
$$
have only finitely  many associated primes?
\end{Question}

  If $R$ is a two dimensional excellent normal local ring, then the question has a positive answer (Proposition 6.7 [C6]).
Lemma 4.5 [CSa2] and the proof of Theorem 4.6 [CSa2] show that if $R$ is an excellent local domain, then $F[\mathcal I(D)]$ can have only finitely many associated primes $P$ such that $\dim F[\mathcal I(D)]/P=\dim R$. We propose to use the techniques of sectional decomposition used in [C10] to prove Theorem \ref{Theorem2} in all dimensions to answer Question \ref{Q20}.



%The more general result in dimension 2 in  [C6] relies on Zariski decomposition which does not extend to higher dimensions as remarked above, although weaker tools such as sectional decomposition [Nak], [C10] are available. This technique is used to prove Theorem \ref{Theorem2} in [C10] in all dimensions 
%More generally, is there an example of a filtration for which there exist infinitely many minimal primes .....



\subsection{Hilbert functions of fiber cones of filtrations}\thispagestyle{empty}
Let $R$ be an analytically unramified local ring and $\mathcal I=\{I_n\}$ be a filtration of $R$.  Let $R[\mathcal I]=\oplus_{n\ge 0}I_n$ be the Rees algebra of $\mathcal I$. The fiber cone of $\mathcal I$ is $F=F[\mathcal I]=R[\mathcal I]/m_RR[\mathcal I]$. The $k=R/m_R$-algebras $R[\mathcal I]$ and  $F[\mathcal I]$ are generally not finitely generated. We have an associated Hilbert function
$$
H_{\mathcal I}(n)= \dim_kF_n=\dim_kI_n/m_RI_n=\lambda_R(I_n/m_RI_n).
$$
\thispagestyle{empty}

In the classical case of the $I$-adic filtration $\mathcal I=\{I^n\}$ of powers of an ideal $I$, we have that $H_{\mathcal I}(n)$ is a polynomial of degree $n^{d-1}$ where $d$ is the dimension of $R$, so that 
\begin{equation}\label{eq10}
\lim_{n\rightarrow\infty}\frac{H_{\mathcal I}(n)}{n^{d-1}}\mbox{ exists}
\end{equation}
and 
\begin{equation}\label{eq11}
 \lim_{n\rightarrow\infty}\frac{H_{\mathcal I}(n)}{n^{d-1}}>0
 \mbox{ if and only if }\ell(\mathcal I)=d.
 \end{equation}
 \thispagestyle{empty}
 
We can easily construct examples of filtrations with highly non polynomial like Hilbert functions. Let $R=k[[x,y]]$, the power series ring in two variables over a field $k$, with maximal ideal $m_R=(x,y)$. Let $\sigma:\NN\rightarrow \NN$ be a function such that 
\begin{equation}\label{eq2}
\sigma(0)=0\mbox{ and }\sigma(n+1)\ge \sigma(n)\mbox{ for all $n$}
\end{equation}
and
\begin{equation}\label{eq3}
\sigma(m+n)\le\sigma(m)+\sigma(n)\mbox{ for all $m$, $n$.}
\end{equation}
Then $\mathcal I=\{I_n\}$ where $I_n=m_R^{\sigma(n)}$ is a filtration of $R$. The Hilbert function $H_{\mathcal I}$ satisfies
$$
H_{\mathcal I}(n)=\dim_km_R^{\sigma(n)}/m_R^{\sigma(n)+1}=\sigma(n)+1
$$
for all $n$.
 Let $\lambda\in \RR$ be such that $0\le\lambda\le1$. Then the round up $\sigma(n)=\lceil n^{\lambda}\rceil$ satisfies (\ref{eq2}) and (\ref{eq3}). (Equation \ref{eq3}) follows from Theorem 19 on page 28 of [HLP]). Then
 \begin{equation}\label{eq4}
 H_{\mathcal I}(n)=\lceil n^{\lambda}\rceil+1.
 \end{equation}
 
 We have that (\ref{eq10})  holds for these examples, but (\ref{eq11}) holds if and only if $\lambda$ is rational.  
 
 
 An important class of filtrations where we can hope to find better behaved Hilbert functions are the $\QQ$-divisorial filtrations $\mathcal I(D)=\{I(nD)\}$. There are many naturally occuring filtrations which are $\QQ$-divisorial. For instance, 
 the filtrations of (generalized) symbolic powers of an ideal in a regular  local ring  are  $\QQ$-divisorial filtrations.
 
  
 A case where the Hilbert functions are very well behaved is  that of the divisorial filtrations in a two dimensional normal local ring.
 
 \thispagestyle{empty}
 
 
\begin{Theorem}(Theorem 1.6 [C6]) \label{HilbThm} Suppose that $R$ is a two dimensional normal excellent local ring and $\mathcal I(D)$ is a $\QQ$-divisorial filtration on $R$.
Then there exist a nonnegative rational number $\alpha$ and a bounded function $\sigma:\NN\rightarrow\QQ$ such that 
\begin{equation}\label{eq6}
H_{\mathcal I(D)}(n)=\lambda_R(I(nD)/m_RI(nD))=\lambda_R((R[D]/m_RR[D])_n)=n\alpha+\sigma(n)
\end{equation}
for $n\in \NN$.
The constant $\alpha$ is positive if and only if the analytic spread 
$$
\ell(\mathcal I)=\dim(R[D]/m_RR[D])=2.
$$
\end{Theorem}
In particular, formulas (\ref{eq10}) and (\ref{eq11}) hold for the filtrations of Theorem \ref{HilbThm}.
 
 
 It is unlikely that the function $\sigma(n)$ of Theorem \ref{HilbThm} will always be eventually periodic. 
 It is shown in Theorem 9 [CS] that if $D$ has exceptional support then
 the Hilbert function of ${\rm gr}_{\mathcal I}(R)=\sum_{n\ge 0}I(nD)/I((n+1)D)t^n$ has an expression
 $$
 \lambda_R(I(nD)/I((n+1)D))=n\beta+\tau(n)
 $$
 where $\beta\in \QQ$ and $\tau(n)$ is a bounded function. If $R$ has equicharacteristic zero then it is shown in Theorem 9 [CS] that $\tau(n)$ is eventually periodic, but Example 5 [CS] gives an example where  $R$ has equicharacteristic $p>0$ and $\tau(n)$ is not  eventually periodic. Thus it is likely that the following  question has a positive answer.
 \thispagestyle{empty}
 \begin{Question} Does there exist  a $\QQ$-divisorial filtration $\mathcal I(D)=\{I(nD)\}$ on a normal  two dimensional local ring $R$ such that  the function $\sigma(n)$ of (\ref{eq6})   is not  eventually periodic? 
  \end{Question}
 
 Suppose that $R$ is an excellent normal local ring of dimension 3. Let $Z\rightarrow \mbox{Spec}(R)$ be  a resolution of singularities and $D$ be an effective divisor on $Z$, all of whose components contract to the maximal ideal $m_R$.  Then the Hilbert function
$$
h(n)=\lambda_R(I(nD)/I((n+1)D))
$$
  may be far from being polynomial like. The examples Example 6 [CS] and Theorem 1.4 [C2] have the property that
$$
\lim_{n\rightarrow \infty}\frac{h(n)}{n^2}
$$
is an irrational number. These examples are in three dimensional equicharacteristic rings $R$ of any characteristic. The reason for this irrational behavior in dimension three is because of the lack of existence of Zariski decompositions in dimension three. 
The existence of Zariski decompositions of divisors on nonsingular projective $\mbox{Spec}(R)$-schemes birationally dominating $\mbox{Spec}(R)$ when $\dim R=2$ plays an essential role in the proof of Theorem \ref{HilbThm}.

We expect the answers to all parts of the following question to be positive. 
  \thispagestyle{empty} 
   \begin{Question}\label{Q10} Suppose that  $\mathcal I(D)=\{I(nD)\}$ is an $\RR$-divisorial filtration   on an excellent  normal  local ring $R$  of dimension $d\ge 3$. The associated Hilbert function is 
   $$
   H_{\mathcal I(D)}(n)= \lambda_R(I(nD)/m_RI(nD)).
   $$ 
 \begin{enumerate}
 \item[1)] Is 
 \begin{equation}\label{eq7}
 \limsup_{n\rightarrow\infty}\frac{H_{\mathcal I(D)}(n)}{n^{d-1}}<\infty? 
 \end{equation}
 \item[2)]
Can (\ref{eq7}) be an irrational number? 
\item[3)]
Does (\ref{eq7}) exists as a limit; that is, is
$$
\limsup_{n\rightarrow\infty}\frac{H_{\mathcal I(D)}(n)}{n^{d-1}}
=\lim_{n\rightarrow\infty}\frac{H_{\mathcal I(D)}(n)}{n^{d-1}}?
$$
\end{enumerate}
\end{Question} 
A positive solution to both 1) and 3) is just the statement that (\ref{eq10}) holds for the filtration. 




We now outline our approach to answering Question \ref{Q10}. Even if we are only interested in answering the question for $\QQ$-divisorial filtrations, in dimensions $\ge 3$ we must solve the problem for $\RR$-divisors (as explained below). The first case to be solved then is for $\RR$-divisorial filtrations on a two dimensional normal excellent local ring. This case involves determining if the proof for $\QQ$-divisors in  Theorem \ref{HilbThm} can be extended to $\RR$-divisors. This case will certainly be possible to solve.
Assuming that this has a positive answer,  
 we restrict  to the case of an  
 excellent normal local ring of equicharacteristic zero and dimension $\ge 3$, so that resolution of singularities holds (resolution of singularities  already follows in this situation from the original paper of Hironaka [H], where he actually proves resolution of singularities for  reduced schemes which are of finite type over an excellent equicharacteristic zero local ring). 
 We propose to answer Question \ref{Q10} by three reductions. 
 We first consider the case where there exists a birational projective morphism $f:X\rightarrow \mbox{Spec}(R)$ such that $X$ is nonsingular, $\mathcal I(D)$ is represented by a divisor on $X$, and $D$ has a rational Zariski decomposition on $X$. This allows us to reduce to the case where $D$ is itself a nef $\QQ$-divisor on $X$, and this case should follow from an extension of the methods used in proving Theorem \ref{HilbThm}, using good properties of nef $\QQ$-divisors such as Fujita's vanishing theorem which can be found in [Laz].  The next case to be considered is when there exists $X$ as above such that $D$ has a real (but not necessarily rational) Zariski decomposition.  
 In this case one can compare with nef rational approximations of the positive part of the Zariski decomposition. The general case can be approached using the theory of antipositive intersection products in [C15] to approximate 
 $\frac{H_{\mathcal I(D)}(n)}{n^{d-1}}$.
for large $n$ on a suitable infinite sequence of blowups above $\mbox{Spec}(R)$, where we have nef divisors which are increasingly close to the pullback of $D$. Antipositive intersection products are a local form of the positive intersection products on projective varieties of [BFJ].

 Another question we can ask is the following. 
\begin{Question}\label{Q11} Suppose that $\mathcal I(D)=\{I(nD)\}$ is a $\QQ$-divisorial filtration   on an excellent  normal  local ring $R$  of dimension $d\ge 3$. Suppose that 
there is a positive solution to both 1) and 3) of Question \ref{Q10}; that is,  (\ref{eq10}) holds for $\mathcal I(nD)$. 
Then does (\ref{eq11}) hold? That is, do we have
\begin{equation}
 \lim_{n\rightarrow\infty}\frac{H_{\mathcal I(D)}(n)}{n^{d-1}}>0
 \mbox{ if and only if }\ell(\mathcal I(D))=d?
 \end{equation}
 \end{Question}

This question is unlikely to have a positive answer since it is false for $\RR$-divisorial filtrations even on a one dimensional regular local ring. In fact, $\mathcal I(D)=\{I_n(D)\}$ where $I_n(D)=(x^{\lceil n\pi\rceil})$ in the ring $R=k[[x]]$ gives such an example. Since $d=\dim R=1$, we have that 
 $H_{\mathcal I(D)}(n)=1$ for all $n$ but $\ell(\mathcal I(D))=0$.  Examples of $\QQ$-divisorial filtrations in dimension 3 which have real (but not rational) Zariski decomposition do exist ([C12], [C15], [C2]).  It is probable that Question \ref{Q11} has a negative answer for such examples.

  
 






In [BS], [HKTT] and [DM],   an alternate definition of the analytic spread of a filtration is given. 
% In the case of a Noetherian filtration, their definitino an two definitions are the same.
They define
$\ell^*(\mathcal I)=\min\{t\in \RR\mid \mu(I_n)=O(n^{t-1})\}$.
If $\mathcal I$ is a noetherian filtration (the Rees algebra $R[\mathcal I]$ is a finitely generated $R$-algebra) then 
  this definition agrees with our definition of analytic spread. That is, 
  $\ell^*(\mathcal I)=\ell(\mathcal I)$.
However, in the case of a non Noetherian filtration $\mathcal I$, $\ell(\mathcal I)$ and $\ell^*(\mathcal I)$ can be different.
In the example after Question \ref{Q11}, we have that $\ell(\mathcal I(D))=0$ but $\ell^*(\mathcal I)=1$.
Theorem \ref{HilbThm} shows that if $D$ is a $\QQ$-divisorial filtration on a two dimensional normal local ring, then
$\ell(\mathcal I(D))=\ell^*(\mathcal I(D))$. However, there are likely to be examples of $\QQ$-divisorial filtrations $\mathcal I(D)$ in dimension $\ge 3$ such that $\ell(\mathcal I(D))\ne \ell^*(\mathcal I(D))$
(by the discussion after Question \ref{Q11}).


 Dao and Montano conjecture in Section 5 of [DM], that in the case of a symbolic filtration $\{I^{(n)}\}$, 
we always have that $\ell^*(\{I^{(n)}\})\le \dim R$.
% and they ask if we even have that $\ell*(\{I^{(n)}\})\le \ell(I)+1$. 
 Dao and Montano prove that these bounds hold in some special cases. They show that if $R/I_n$ is a Cohen-Macaulay $R$-module for all $n\gg 0$, then $\ell^*(\mathcal I)\le \dim R$.

% have that $\ell^*(\mathcal I)=\ell(\mathcal I)$ in the case of this theorem. It also follows from Corollary 3.12 [DM] that these two numbers are the same in this case.




%\begin{Question} Suppose that $\mathcal I(D)$ is a divisorial filtration on a normal excellent domain. Is 
%$\ell^*(\mathcal I(D))=\ell(\mathcal I(D))$?
%\end{Question}

%Recall that symbolic filtrations are divisorial on regular local rings.   Since we always have that $\ell(\mathcal I)\le \dim R$ for a filtration $R$, this would give a proof of the conjecture of [DM] in the case of symbolic filtrations of regular local rings. 

\subsection{Generalized Multiplicities}\thispagestyle{empty}
We propose to explore the relationship between various generalized multiplicities, including ordinary multiplicities, Buchsbaum-Rim multiplicities, $j$-multiplicities and $\epsilon$-multiplicities. A motivating theorem is the following volume=multiplicity formula.  

\begin{Theorem}\label{Theorem15}(Theorem 6.5 [C13])  Suppose that $R$ is a $d$-dimensional analytically unramified local ring and  $\{I_i\}$ is a graded family of $m_R$-primary ideals in $R$. Then 
$$
\lim_{n\rightarrow \infty}\frac{\lambda_R(R/I_n)}{n^d/d!}=\lim_{p\rightarrow \infty}\frac{e_R(I_p)}{p^d}
$$
exists.
Here $e_R(I_p)$ is the ordinary multiplicity of an ideal
$$
e_R(I_p)=\lim_{k\rightarrow \infty} \frac{\lambda_R(R/I_p^k)}{k^d/d!}.
$$
\end{Theorem}


Theorem \ref{Theorem15} is proven for valuation ideals associated to an Abhyankar valuation in a regular local ring which is essentially of finite type over a field in  [ELS], for general families of $m_R$-primary ideals when $R$ is a regular local ring containing a field in [Mus] and when $R$ is a local domain which is essentially of finite type over an algebraically closed field $k$ with $R/m_R=k$ in Theorem 3.8 [LM]. 
\thispagestyle{empty}



Let $R$ be a Noetherian local ring, $A=\oplus_{n\ge 0} A_n\subset B=\oplus_{n\ge 0}B_n$ be graded inclusions of standard graded $R$-algebras. The relative $\epsilon$-multiplicity of $A$ and $B$ is defined by Ulrich and Validashti in [UV],\begin{equation}\label{eq1}
\epsilon(A|B)=\limsup_n(\dim B-1)!\frac{\lambda_R(P_n/A_n)}{n^{\dim B-1}}.
\end{equation}
where $P_n=A_n:{B_n}m_R^{\infty}$. This generalizes the epsilon multiplicity $\epsilon(I)$ of an ideal defined in Subsection \ref{SubsecEp}.


Suprajo Das showed in [Da] that If $R$ is excellent and reduced, $B$ is reduced and if $P$ is a minimal prime of ideal of $B$ then $P\cap R\ne m_R$, then the limsup in (\ref{eq1}) actually exists as a limit. 

Suppose that $B$ is a birational extension of $A$, $r=\dim A-1=\dim B-1$.
Let $P(n)=\sum_{k\ge 0}P_n^k\subset A(n)=\sum_{k\ge 0}A_n^k$, so that $P(n)$ and $A(n)$ are respective Veronese algebras.
 For large $n$, the Hilbert polynomial of $P(n)/A(n)$ is a polynomial
 $$
 e_R(P(n)_k,A(n))\frac{k^r}{r!}+\cdots
 $$
 where $e_R(P(n),A(n))$ is the Buchbaum-Rim multiplicity of $A(n)$ and $B(n)$. This multiplicity  is explored in the context of equisingularity in [KT].
 
The following question was posed by  Antoni Rangachev; he has shown that a positive answer to Question \ref{Q1}  would  prove local volume formulas. The theory of local volume in algebraic geometry is developed in [Ful]  and [Ra].
\thispagestyle{empty}
\begin{Question}\label{Q1} Suppose that $R$ is a reduced local equidimensional excellent local ring with $\dim R\ge 2$. 
Suppose that $A\subset B$ are  reduced equidimensional standard graded $R$-algebras such that $A\rightarrow B$ is birational and that $A_1$ contains a nonzero divisor in $A$. Assume that the minimal primes of $A$ contract to minimal primes of $R$. Then is
$$
\epsilon(A|B)=\lim_{n\rightarrow\infty}\frac{e(A(n),P(n))}{n^r}?
$$
\end{Question} 

 
Using methods from [C13] and [Da], we expect to be able to  settle this conjecture. We have proven the following related theorem using methods of [C13]. The methods in [C13] are a mixture of Okounkov body methods, commutative algebra and algebraic geometry. Theorem \ref{Thm5} gives a solution to Question \ref{Q1} in the case that $A=\sum_{n\ge 0}I^nt^n\subset B=R[t]$ is  the Rees algebra of an ideal $I$ of $R$.

\begin{Theorem}\label{Thm5}
Suppose that $R$ is a $d$-dimensional analytically unramifed local ring and $I$ is an ideal in $R$. Then
for all $m>0$, the limits
\begin{equation}\label{eq9}
e_R(I^m,(I^m)^{\rm sat})=d!\lim_{k\rightarrow\infty}\frac{\lambda_R([(I^m)^{\rm sat}]^k/I^{km})}{k^d}
\end{equation}
exist. Further,
$$
\epsilon(I)=\lim_{m\rightarrow\infty}\frac{e_R(I^m,(I^m)^{\rm sat})}{m^d}.
$$
\end{Theorem}
Here $(I^n)^{\rm sat}=I^n:m_R^{\infty}$ and $\epsilon(I)$ is the epsilon multiplicity of $I$ defined in Subsection \ref{SubsecEp}.

Another direction to be explored is  $j$-multiplicity. The $j$-multiplicity of an ideal was defined by 
Achilles and Manaresi [AM]. It extends the multiplicity of an $m_R$-primary ideal to arbitrary ideals.  Suppose that $\mathcal I=\{I_n\}$ is a filtration in a local ring $R$. The $j$-multiplicity of an ideal  can be extended to filtrations by the formula
$$
j(\mathcal I)=\limsup (d-1)!\frac{\lambda_R(H^0_{m_R}(I_n/I_{n+1})}{n^{d-1}}.
$$
Using methods from [C13], we can show that this limsup is a limit if $R$ is  analytically unramified. 
As for the case of ideals, we have  the inequality for filtrations, $\epsilon(\mathcal I)\le j(\mathcal I)$.
We will explore generalizations of theorems about $j$-multiplicity of ideals to filtrations. The theory of multiplicity of $m_R$-primary ideals extends remarkably well to filtrations (as shown in [C13] and our subsequent papers) so it is to be expected that there is a   good answer here.

\begin{Question} Is there a good generalization of the theory of $j$-multiplicity for ideals to filtrations?
\end{Question}



\subsection{Inseparable local uniformization}\thispagestyle{empty}
Resolution of singularities  for characteristic zero algebraic varieties was proven by Hironaka [H] in the mid 1960s. However, resolution of singularities in positive and mixed characteristic is still open in dimension larger than 3. The case of dimension 3 was proven in characteristic $p>5$ by Abhyankar in the mid 1960s [Ab], and in 2019 for excellent reduced three dimensional schemes by Cossart and Piltant [CoP]. De Jong [dJ] proved resolution of singularities, after blowing up and taking a finite extension (an alteration) in 1996. This proof uses the theorem on semistable reduction of families of curves. A stronger theorem of this type has been proven by Gabber.  A local form of resolution of singularities was proposed by Zariski and solved by him for algebraic varieties over a field of characteristic zero in the late 1930s and early 1940s [Z1], [Z2], [Z3]. Let $R$ be an excellent local domain with quotient field $K$ and  $v$ be a valuation of $K$ which dominates $R$. Then a local uniformization of $v$ along $R$ is an extension $R\rightarrow S$ such that $S$ is a regular local ring with quotient field $K$ which is dominated by $v$ and $S$ is essentially of finite type over $R$.  If $S$ is weakened to have a quotient field which is a finite extension $L$ of $K$, and $S$ is dominated by an extension of $v$ to $L$, then $S$ is called a weak local uniformization.   The existence of weak local uniformization follows from the work of de Jong [dJ] and Gabber. 
Knaf and Kuhlmann  [KK2] proved in 2009  that weak local uniformization holds after taking a separable extension.
Temkin [T2] has shown in 2013 that weak local uniformization holds after taking a purely inseparable extension. 
% Deeply ramified fields play an important role in Temkin's work, which is one of the focuses of this proposal. 
The proofs of Knaf and Kuhlmann and of  Temkin start with local uniformization of Abhyankar valuations in an algebraic function field such that the residue field of the valuation is separable over the ground field. This result  was proven by Knaf and Kuhlmann in [KK1]  and by  Temkin in [Te].  Local uniformization of Abhyankar valuations in an algebraic function field with no restrictions on the residue field is proven in [C3].

It is conjectured by Abramovich and Oort and by Temkin that resolution of singularities of varieties is true after taking an inseparable alteration.


There are three main techniques for understanding extensions of valuations in simple extensions. The first method is that of constructing generating sequences for a valuation by successive construction of  key polynomials approximating the valuation. This method was  introduced by MacLane [M1], [M2] in the 1930s for discrete valuations, and extended to arbitrary valuations by Vaqui\'e in [V1], [V2].
An improvement on this method has recently been introduced by Novacoski and Spivakovsky [NS1]. 
The approximants describing the extension of valuations are finite (in the case that the extension of valuations is unique)
if there is no defect. Defect is discussed in Subsection \ref{SecRam}. Defect can only occur in an extension where the residue characteristic is positive, and is the essential obstruction to local uniformization. 
The second method is that of pseudo convergent sequences, introduced by Ostrowski [Os] and Kaplansky [Kap]. The third method is that of classification of norms on the Berkovich affine line and curves by Berkovich [Be]. 
Although the construction of key polynomials and pseudo convergent sequences are closely related, the approach of Berkevich is quite different, and there is a need of giving a translation between these  methods.  

One of our goals  is to provide a bridge between the known methods of resolution, specifically between those arising from algebraic geometry, valuation theory, non archimedean analysis and almost mathematics. A specific direction is to find a  proof of inseparable local uniformization (Temkin's theorem) which does not rely on nonarchimedean analysis, in the hope that this will give some new understanding and progress on the local uniformization problem. 
A motivating problem is to   prove inseparable local uniformization for algebraic varieties.

\begin{Question} Is it possible to find a translation between the theory of key polynomials of simple extensions of valued fields and the theory of the  Berkovich affine line?
\end{Question}

\begin{Question} Is it possible to give an alternate proof of inseparable local uniformization which does not require the Berkovitch space techniques of  Temkin's proof?
\end{Question}

\subsection{Ramification, K\"ahler differentials and almost \'etale extensions}\label{SecRam} \thispagestyle{empty}

By $(L|K,v)$ we denote a 
field extension $L|K$ where $v$ is a
valuation on $L$ and $K$ is endowed with the restriction of $v$. The valuation
ring of $v$ on $L$ will be denoted by $\cO_L\,$, and that on $K$ by
$\cO_K\,$. Similarly, $m_L$ and $m_K$ denote the maximal ideals of $\mathcal O_L$ and $\mathcal O_K$.
The value group of the valued field $(L,v)$ will be denoted by $vL$, and
its residue field by $Lv$. The value of an element $a$ will be denoted by $va$, and 
its residue by~$av$.
A valued field extension $(L|K,v)$ is unibranched if the
extension of $v$ from $K$ to $L$ is unique. 
%Note that a unibranched extension is automatically algebraic, since every 
%transcendental extension always admits several extensions of the valuation.
%
If in addition the extension is finite, then by the Lemma of Ostrowski
(Corollary to Theorem 25, Section G, page 78 of [ZS2]),
\begin{equation}                    \label{feuniq}
[L:K]\>=\> p^{\nu }[vL:vK][Lv:Kv]\>,
\end{equation}
where $\nu$ is a non-negative integer and $p=\chara Kv$ if it is 
positive and $p=1$ otherwise. The factor $d(L|K,v):=p^{\nu }$ is 
the defect of the extension $(L|K,v)$.
In [Ku] and [KR], Kuhlmann and Rzepka classify defect Artin-Schreier and Kummer extensions are classified as dependent or independent. 



Almost mathematics was introduced by Faltings in [Fa] and given an extensive categorical development by Gabber and Ramero in [GR]. A nice survey of almost mathematics, perfectoid spaces and  applications to commutative algebra is given in the lecture notes [Bh] of  Bhargav Bhatt. Some of the basic theory is explained in the language of ordinary commutative algebra in [NaSh].  The basic set up of almost mathematics is a pair $(R,I)$ where $R$ is a ring and $I$ is an ideal in $R$ such that $I^2=I$. We also make the commonly made assumption that $I$ is a flat $R$-module. We will be concerned with the case where $R$ is a valuation ring $V$, so that all all ideals and extension rings of $V$ are flat $V$-modules.

Fix a basic setup $(R,I)$ and let $\phi:A\rightarrow B$ be an $R$-algebra homomorphism. The concept of $\phi$ being $I$-almost  finite \'etale is defined in [Fa], [GR], [NaSh] generalizing the classical notion of this concept in [EGAIV]. Their definition is through an almost version of the notion of a separable algebra. This formulation in ordinary mathematics $(I=R)$ is explained in [F].

We ask the following question.
\begin{Question}\label{Q6} Let $(L|K,v)$ be a finite Galois extension of valued fields and $(\mathcal O_K,I)$ be a basic set up such that $I$ is a prime ideal in $\mathcal O_K$ with $I^2=I$ or $I=\mathcal O_K$. Is there a good characterization of  when $\mathcal O_K\rightarrow \mathcal O_L$ is almost finite \'etale, or of some related concept such as almost \'etale or almost weakly \'etale?
\end{Question}


For $\mathcal O_K\rightarrow \mathcal O_L$ to be almost finite \'etale, $\mathcal O_K\rightarrow \mathcal O_L$ must be unibranched so we assume this condition from now on in the discussion of this problem.  

If $I=\mathcal O_K$ (ordinary mathematics) there is a simple characterization answering Question \ref{Q6}. In this case, $\mathcal O_K\rightarrow \mathcal O_L$ is almost finite \'etale if and only if it is finite \'etale if and only if $\mathcal O_K\rightarrow \mathcal O_L$ is an inertial extension (the Galois group is the inertia group).

If $v$ has rank 1 and is not discrete and $I=m_K$, then a finite separable unibranched extension $(L|K,v)$ is almost finite \'etale if and only if $\Omega_{\cO_L|\cO_K}$ is almost zero. This follows from the theory developed in [GR], but can also proven directly, by  relatively simple arguments.  In our papers [CKR] and [CK] we compute the K\"ahler differentials of finite Galois extensions, by reducing this computation to the study of Artin-Schreier and Kummer extensions. We compute the K\"ahler differentials of these extensions, and we deduce when they are zero. If the K\"ahler differentials of an Artin-Schreier  or Kummer extension are almost zero, in almost all cases the K\"ahler differentials are actually zero, but it is possible for them to be almost zero but not zero. 
In the case of defect Artin-Schreier and Kummer extensions such that $v$ has rank 1 and $I=m_K$, we have the characterization that $\mathcal O_K\rightarrow \mathcal O_L$ is almost finite \'etale  if and only if 
$\Omega_{\mathcal O_L|\mathcal O_K}$ is almost zero  if and only if $\Omega_{\mathcal O_L|\mathcal O_K}=0$ if and only if $L/K$ is an independent defect extension. The defect extensions are the most troublesome ones; these extensions are the obstruction to proving local uniformization in positive and mixed characteristic.

The assumption that the valuation $v$ has rank 1 is used to prove that in a  finite unibranched separable extension, $\mathcal O_K\rightarrow \mathcal O_L$ is almost finitely generated, a requirement for the extension to be almost \'etale. Probably almost finite generation can fail in finite unibranched separable extensions if the rank of  $v$ is larger than 1. Determining if this is indeed the case is an interesting question to be addressed. If almost finite generation can fail in higher rank, then the question becomes to find the correct generalization of almost finite \'etale to such extensions.

\subsection{Generating sequences of simple extensions of valuation rings and K\"ahler differentials}\thispagestyle{empty}
Extensions of valuation rings in simple extensions are studied by construction of sequences of key polynomials by MacLane [M1], [M2] for discrete (rank 1) valuations and this construction is extended to arbitrary valuations by Vaqui\'e [V1], [V2].  A refinement of this theory has  been developed by Novacoski and Spivakovsky [NS1]. 

Suppose that $K\rightarrow L$ is a finite simple extension of valued fields. Recently Novacoski and Spivakovsky have given a construction  in [NS2] of the K\"ahler differentials $\Omega_{\mathcal O_L|\mathcal O_K}$ using properties of a sequence of key polynomials defining the extension. Their results mostly assume that the valuations have rank 1. We seek to interpret this result in terms of the construction of $\Omega_{\mathcal O_L|\mathcal O_K}$ in [CKR] and [CK]. 

\begin{Question} Let $K\rightarrow L$ be a unibranched Kummer or Artin-Schreier extension  of (rank 1) valued fields of prime degree. Interpret the results of [NS2]   in terms of the construction of $\Omega_{\mathcal O_L|\mathcal O_K}$ given in [CKR] and [CK].
\end{Question}


\vskip .1truein
\noindent{\bf Broader Impacts.}\thispagestyle{empty}

Because of the diversity of the people who are attracted by the subject matter of my research I have had the good fortunate to have had students and postdoctoral mentees  across all  races and religions.
This diversity  has enriched both my research and collaborations. 

Thirteen  graduate students  have earned their PhDs  with me and two students are working with me currently towards a PhD. 

Russell Goward (PhD in 2001) was the first African American to be awarded a PhD in Mathematics from the University of Missouri. 
Olga Kashcheyeva (PhD in 2003), Samar El-Hitti (PhD in 2008), Razieh Ahmadian (PhD in 2015 from IPM, Tehran), Smita Praharaj (PhD in 2023) are Women. Roberto N\'u\~nez (PhD in 2021) is Hispanic. Krishna Hanumanthu (PhD in 2008), Soumya Sanyal (PhD in 2014), Vinh Pham An (PhD in 2014), Arpan Dutta (PhD in 2018), Suprajo Das (PhD in 2020),  Navaneeth Chenicheri (PhD in 2022) and Smita Praharaj (PhD in 2023) are Asian. 
 
 My recent contributions to Broader Impacts are documented in the Broader Impacts section of the Results from Prior NSF Support. 


 %My book ``Introduction to Algebraic Geometry'', based on my lectures at the University of Missouri,  was published by the  American Mathematical Society in 2018. My book ``Resolution of Singularities'' was published by the American Mathematical Society in 2004, and is based on lectures that I gave at Chennai Math. Institute. 


\thispagestyle{empty}

\vskip .1truein
\noindent{\bf Upcoming conferences that I am coorganizing:}
\begin{enumerate}
\item Conference on Resolution of Singularities, Valuation Theory and Related Topics, to be held August 5 - 9, 2024, Centro de Ciencies Matem\'aticas, UNAM, Morelia, Mexico. Member of Scientific Committee.
\item Program on Commutative Algebra, to be held Spring semester, 2024 at MSRI (now SLMath), Berkeley. Coorganized with Aldo Conca, Claudia Poloni, Claudiu Raicu, Steven Sam and Kevin Tucker. 
\item KUMUNU Conference on Commutative Algebra, to be held at the University of Missouri, Sept. 23 - 24, 2023. Coorganized with Ian Aberbach, Rankeya Datta, Hema Srinivasan and  Dana Weston.
\end{enumerate}
\vskip .1truein

I have been a coorganizer, from October 2022,  of The Fellowship of the Ring, a virtual seminar in commutative algebra, with Christine Berkesch, David Eisenbud, Mel Hochster, Srikanth Iyengar, Claudia Miller, Irena Peeva, Karl Schwede and Bernd Ulrich.

\vskip .1truein
 I currently have applications pending or in preparation with colleagues for upcoming international workshops and programs.  
  \vskip .1truein
  
 I have been an editor of the Journal of Algebra since 2005. 
 My books ``Introduction to Algebraic Geometry''  and ``Resolution of Singularities'' 
 were published by the  American Mathematical Society.  Both of these books are  accessible to a large community of mathematicians, especially to graduate students. 
 

  
%\vskip 2truein

%$j$-multiplicity of filtrations: an extension of multiplicity to arbitrary (not necessarily $m$-primary) ideals. 

%Existence for filtrations (true in many cases). What extends to filtrations?  What is the relation to epsilon multiplicity?  

%State Conjecture 1.1 of Rangachev. 

%Solve using techiniques of my earlier papers and the thesis of Suprajo. Explain the connection and what needs to be done. 

%Rangachev has shown that a proof of this conjecture would give a simpler proof of the ``local volume formula'' (of Fulger, Rangachev). 

%Describe relationship with Buschbaum-Rim and j-multiplicity.
%Does this give an interpretation of epsilon multiplicity in terms of j-multiplicity?  Theory in terms of divisorial valuations and their approximations. 

%Is there an approximation formula for epsilon multiplicity (of divisorial valuations) in terms of their epsilon multiplicities?

%Is there a limit formula for epsilon/j-multiplicities in terms of finitely generated approximations like the Lazarsfeld/Mustata formula for ordinary multiplicities?







\end{document}




